% !TEX root = ../Dissertation.tex

% ist ist anzunehmen / lässt sich daraus schließen 
% mit einer gewissen vorsicht
% gemessen an...
% nur moderierende Faktoren, die mit mehr oder weniger XY in Verbindung gebracht werden können
% (keine) kausalen, also ursächlichen Faktoren / kausaler Zusammenhang
% XY lässt Vorhersagen zu, ob...
% XY kann mit YZ in Verbindung gebracht werden
% Es handelt sich dabei zunächst um eine Korrelation, keine Kausalität
% Nichtsdestotrotz liegt die Hypothese nahe, dass...
% hart und belastbar
% interessiert an relativen Interaktionen bzw. Unterschieden: Satz A wird signifikant höher gerankt als Satz B
% auf Floor-/Ceiling-Effekte eingehen
% skalenbasierte Beschreibung
%Im Verfahren des Chi-Quadrats wurde die Nullhypothese höchstsignikant (p < 0,0005) widerlegt, ersichtlich
% Bezüglich des empirischen Vorhabens der vorliegenden Arbeit wurden zwei Ziele formuliert. Erstens sollten dXY statistisch geprüft werden, zweitens sollten statistische Tests auch den Einfluss XY auf YZ feststellen


%parametrische Größen (Größe, Länge, Breite) (vgl. Stopyra 2020: 272)




\chapter{{Experimentelle Studien -- Teil I}} \label{Experiment0}
Dieses Kapitel führt in den zweiten empirischen Teil ein, der sich mit der experimentellen Untersuchung der in der Korpusstudie berücksichtigten expressiven Intensivierer beschäftigt. Die zentrale Idee ist es, Intensivausdrücke von Versuchspersonen auf eine Skala abbilden zu lassen, um den in der Literatur angenommenen wahrnehmbaren Unterschied des durch sie ausgedrückten Grades an Skalarität und Sprechereinstellung herauszuarbeiten (vgl. {\citealt[315]{Lang.1983}}; {\citealt[22]{Paradis.1997}}; \citealt[57]{Ortner.2014}; {\citealt[133]{Gutzmann.2019a}}). In den ersten Abschnitten wird zunächst auf den theoretischen Rahmen, den Untersuchungsaufbau und die Untersuchungshypothesen eingegangen. Daraufhin folgt eine Präsentation des Untersuchungsgegenstands, der Personenstichprobe und der Auswertungskriterien. Der darauffolgende Ergebnisteil setzt sich aus deskriptiv- und inferenzstatistischen Auswertungen zusammen. Nachdem die aus dem Experiment gewonnenen Resultate diskutiert wurden, bildet ein Fazit den Schlussteil des Kapitels.




\section{Theoretischer Rahmen}\label{abs:}
Kern des in diesem Kapitel vorgestellten Experiments sind die Fragestellungen und das lexikalische Material, das im Rahmen der Korpusstudie einer diachronen Frequenzanalyse unterzogen wurde. Dieser lag die Idee zugrunde, dass ein inflationärer Gebrauch eines Intensivierers zulasten seiner Expressivität geht und er in der Folge entweder seltener oder -- ähnlich wie \textit{sehr} -- nur noch mit skalarer Bedeutung verwendet wird {\citep[vgl.][2]{Hauschild.1899}}. {\citet[274]{Baumgarten.1908}} sieht darin die Ursache für die Wandelbarkeit von Intensivierern, führt der ständige Verschleiß alter Ausdrücke doch zu einem kontinuierlichen Bedarf an neuen, die aus verschiedenen semantischen Herkunftsbereichen Einzug in die Partikelebene halten (vgl. {\citealt[126f.]{Biedermann.1969}}; \citealt[18]{Bolinger.1972}; \citealt[]{Claudi.2006}).


Um die longitudinale Entwicklung sichtbar zu machen, wurden die Gebrauchsfrequenzen einer umfangreichen Stichprobe an gegenwärtig verwendeten Intensivierern in einem Zeitschriftenkorpus erhoben. Weil es sich bei den durch die Untersuchung gewonnenen Ergebnissen um frequenzbasierte Daten handelt, die bloß den diachronen Verlauf der Ausdrücke im Korpus abbilden, war es nicht möglich, direkte Rückschlüsse auf den Expressivitätsgrad zu ziehen. Stattdessen wurde davon ausgegangen, dass die Frequenz als Spiegelbild des potenziellen Expressivitätsverlusts dienen kann. Vor diesem Hintergrund wurde den Expressivitätsmaximen nach {\citet[]{Kirschbaum.2002b}} (vgl. {Abschnitt \ref{Auslöser}}) folgend angenommen, dass eine quantitativ rückläufige Entwicklung ein hinreichendes Indiz für den Schwund der Expressivität eines Intensivierers darstellt, der u.\,a. für \citet[2]{Hauschild.1899} eine unmittelbare Folge von inflationärem Gebrauch ist (vgl. H1). Daneben wurde die Möglichkeit in Betracht gezogen, dass auch die beobachtete stetige Frequenzzunahme der hochfrequenten Ausdrücke in den Zeitschriften auf einen fortgeschrittenen Expressivitätsverlust hindeutet. In diesem Kontext wurde diskutiert, dass ein Rückgang der Expressivität im schriftlichen Sprachgebrauch eventuell zu einer besseren Eignung, einer größeren Flexibilität im Gebrauch und einer entsprechend höheren Verwendungshäufigkeit führt. Das Auftreten neuer Intensivierer weist überdies auf einen tatsächlichen Bedarf an neuen expressiven Ausdrücken zur Kompensation hin, wie schon \citet[274]{Baumgarten.1908} annimmt {(vgl. H2)}. Erste Evidenz für die Gültigkeit der Hypothese H1 zeigte sich neben hochfrequenten \glqq{}schrecklichen\grqq{} Intensivierern wie \textit{verdammt} und \textit{wahnsinnig} insbesondere durch die diachron abnehmende Entwicklung von \textit{derb}, während die neuen Intensivierer \textit{hammer}, \textit{krass} und \textit{mega} ferner die Korrektheit der Hypothese H2 nahelegten.


Da bislang unklar ist, welchen Expressivitätsgrad die im Korpus analysierten Intensivierer aufweisen, soll die folgende Untersuchung diesen bestimmen und die Schlussfolgerungen aus Kapitel \ref{Korpus} überprüfen. Die zuvor formulierten Annahmen bilden den Ausgangspunkt für dieses Experiment, das zugleich die Brücke zwischen korpuslinguistischer und experimenteller Vorgehensweise schlägt.




%Neben den Häufigkeiten lieferten die Daten weiterhin auch Informationen über das Erstvorkommen und die Existenzdauer der begutachteten Ausdrücke. 


%Den theoretischen Hintergrund von Experiment A bildet zum einen die in Kapitel \ref{Theorie} angeführte Idee der mentalen Skala, die Intensivierern zugrunde gelegt wird und dazu dient, den durch die Partikel hervorgerufene erhöhte Ausprägungsgrad einer Eigenschaft oder eines Zustands in Beziehung zum Ausprägungsgrad eines unintensivierten Ausdruck zu setzen.


%Einen weiteren Teil der Grundlage bildet zum anderen...


%Gottfried R. Marschall in Schmale (2011: 106): "Was man nachprüfbar messen und physikalisch bestimmen kann, ist der Inbegriff der Objektivität. --> Idee: etwas Subjektives wie Sprecherbewertung in Form einer ordinalen Skala objektiv quantifizierbar zu machen


\section{Hypothesen und Untersuchungsaufbau}\label{abs:Hypothese und Forschungsaufbau2}
Die Idee hinter dieser experimentellen Untersuchung besteht darin, den gegenwärtigen Expressivitätsgrad der Intensiverer zu ermitteln. Wie oben erwähnt, steht dabei in erster Linie die im Rahmen der Korpusstudie dargelegte Frequenzhypothese H1 nach {\citet[2]{Hauschild.1899}} im Fokus, die als ursächlich für die hohe Wandelbarkeit von Intensivierern gilt und sich unter Einbezug des Parameters Zeit \citep[vgl.][59]{Tobler.1868} wie folgt {interpretieren lässt}:

\begin{itemize} 
\setlength{\itemindent}{1cm}
\item[\textbf{\textsc{H1-Freq}}] Eine hochfrequente Verwendung von Intensivierern führt zu einer pragmatischen Abnutzung und demzufolge sukzessive zum Verlust ihrer Expressivität. 

\item[\textbf{\textsc{H1-Temp}}] Eine persistente Verwendung von Intensivierern führt zu einer pragmatischen Abnutzung und demzufolge sukzessive zum Verlust ihrer Expressivität. 
\end{itemize}

\noindent Für die empirische Grundlage dieses Experiments bedeutet die Spezifizierung nach H1-Freq und H1-Temp, dass sowohl im Korpus hochfrequente (z.\,B. \textit{super} oder \textit{verdammt}) als auch lange existierende Ausdrücke (z.\,B. \textit{arg} oder \textit{furchtbar}) möglicherweise weniger expressiv sind als diejenigen Ausdrücke, die weniger frequent (z.\,B. \textit{sau} und \textit{voll}) oder neuer bzw. erst im Laufe des Untersuchungszeitraums erstmals im Korpus in intensivierender Funktion in Erscheinung getreten sind (z.\,B. \textit{krass} oder \textit{mega}). Im Sinne von \citet[87]{Biedermann.1969} und \citet[408]{Breindl.2007} ist daher davon auszugehen, dass die korpuslinguistisch untersuchten Ausdrücke graduelle Abstufungen ihres Intensitätsgrades aufweisen, die sich jedoch nicht mit quantitativen Daten erfassen lassen. Die potenziellen Abstufungen zu ermitteln, ist das primäre Ziel dieser Untersuchung. Von besonderem Interesse ist das expressive \textit{derb}, das aufgrund seiner rückläufigen Entwicklung im Korpus als einziger Intensivierer hinreichende Evidenz für ein grundsätzliches Zutreffen der Hypothese H1 lieferte. Bei diesem Ausdruck ist vorstellbar, dass er von den Versuchspersonen als nur wenig expressiv eingeschätzt wird und auf einer Expressivitätsskala dementsprechend niedrig verortet wird.


%Wie an mehreren Stellen ausgeführt, wird Expressivität hier grundsätzlich als ein erhöhter Grad an emotional-evaluativer Sprechereinstellung und als Motor für pragmatischen Wandel verstanden, mit dem Ziel, \glqq{}to express 'old things in new ways'\grqq{} (\citealt[171]{Androutsopoulos.1999}).


%Zusammengefasst dient die Untersuchung also dazu, die in Kapitel \ref{Korpus} vorgenommene quantitative Einordnung der Intensivierer um eine Einschätzung ihres Expressivitätsgrades zu erweitern. Durch experimentelle Verfahren wie das in diesem Experiment durchgeführte Skalenrating kann die Expressivität der Ausdrücke ermittelt und das Zutreffen der zu testenden Hypothesen H1-Freq und {H1-Temp} überprüft werden. Die Grundidee liegt darin, die Expressivität von Sätzen mit Intensivausdrücken, bestehend aus den berücksichtigten Intensivierern und Adjektiven, von Versuchspersonen auf eine numerische Skala abbilden zu lassen. Die zugewiesenen Skalenwerte können hinterher gemappt und miteinander verglichen werden.


Neben den Hypothesen H1-Freq und H1-Temp wird im Experiment des Weiteren der Gültigkeit der in der Literatur vorherrschenden Annahmen nachgegangen, dass Expressivität gleichermaßen mit einem höheren Grad an Skalarität (vgl. {\citealt[22]{Paradis.1997}}; {\citealt[133]{Gutzmann.2019a}}) und Sprechereinstellung (vgl. {\citealt[315]{Lang.1983}}; {\citealt[57]{Ortner.2014}}) einhergeht. Die in diesem Kontext zu überprüfenden Hypothesen lauten folgendermaßen:

\begin{itemize}
\item[\textbf{\textsc{H3}}] Expressive Intensivierer werden mit einem höheren Grad an Skalarität
assoziiert als deskriptive Intensivierer. 

\item[\textbf{\textsc{H4}}] Expressive Intensivierer geben eine Einstellungsbekundung des Sprechers
zu dem im Satz beschriebenen Sachverhalt wieder und werden in der Folge mit einem höheren Grad an Sprechereinstellung assoziiert als deskriptive Intensivierer. 
\end{itemize}

\noindent Um die beiden Merkmale getrennt voneinander abzufragen, wurden der Untersuchung zwei Skalen zugrunde gelegt, die den Versuchspersonen als Instrument für die Einschätzung des jeweils zutreffenden Grades dienten. Die Idee bestand darin, beide Skalen getrennt und damit unabhängig voneinander zu präsentieren, um eine gerichtete Abhängigkeit der Merkmale nachweisen zu können. Weiterführend zur Hypothese H4, nach der Expressivität der Einstellungsspiegelung des Sprechers dient, wird die Überprüfung einer weiteren in diesem Zusammenhang relevanten Hypothese angestrebt:

\begin{itemize} 
\item[\textbf{\textsc{H5}}] Die Einstellungsbekundung, die expressive Intensivierer zu dem im Satz beschriebenen Sachverhalt wiedergeben, ist evaluativer Natur und wird daher als positiv oder negativ wahrgenommen.
\end{itemize}

\noindent Diese Annahme ergibt sich im Wesentlichen aus einer Schlussfolgerung von {\citet[135]{Gutzmann.2019a}}, nach der die Einstellungsspiegelung, die expressive Intensivierer kommunizieren, evaluativen Charakters ist und sich in Form einer positiven oder negativen Bewertung des Sachverhalts niederschlägt. Deskriptive Intensivierer drücken dagegen keine Einstellung des Sprechers aus und werden demnach dahingehend mutmaßlich vorwiegend als neutral beurteilt. Diese Überlegung wurde mithilfe einer dritten Skala untersucht.


% in einer dritten Aufgabe untersucht, in der die Versuchspersonen dazu aufgefordert waren, die Polarität der Sprechereinstellung mithilfe einer komparativen Konstant-Summen-Skala einzuschätzen und 100 virtuelle Taler auf die drei zur Option stehenden Auswahlfelder 'positiv', 'neutral' und 'negativ' zu verteilen.



%{\citet[681]{McCready.2008}}


%Zusammengefasst dient die Untersuchung also dazu, die in Kapitel \ref{Korpus} vorgenommene quantitative Einordnung der Intensivierer um eine Einschätzung ihres Expressivitätsgrades zu erweitern. Durch experimentelle Verfahren wie das in diesem Experiment durchgeführte Skalenrating kann die Expressivität der Ausdrücke ermittelt und das Zutreffen der zu testenden Hypothesen H1-Freq und {H1-Temp} überprüft werden. Die Grundidee liegt darin, die Expressivität von Sätzen mit Intensivausdrücken, bestehend aus den berücksichtigten Intensivierern und Adjektiven, von Versuchspersonen auf eine numerische Skala abbilden zu lassen. Die zugewiesenen Skalenwerte können hinterher gemappt und miteinander verglichen werden.




%Dabei wird die von \citet[46]{Ruf.1996} vorgeschlagene Definition von 'Verstärkung' zugrunde gelegt, die den \glqq{}Ausdruck des hohen Grades oder der hohen Intensität einer bezeichneten Eigenschaft oder eines Vorganges \textelp{}, sowie die Bezeichnung der besonderen Größe oder Ausdehnung des Genannten\grqq{} meint. Dies bedeutet, dass Intensivierer die durch das nachfolgende Element ausgedrückte Eigenschaft insofern verstärken, als diese durch sie einen höheren Wert auf der Intensitätsskala erlangen. Nach \citet[951]{Weinrich.1993} komme es dabei zu einem impliziten Vergleich, ohne dass der Vergleichsstandard (im Unterschied zum Komparativ) explizit verbalisiert wird. Davon ausgehend ist ein Satz wie \textit{Das Kind ist mega laut} immer in Bezug zu einem normal lauten Kind zu verstehen, das in der intensivierten Aussage zwar nicht mitangegeben, aber dennoch unterschwellig thematisiert ist. Die Verstärkung des Adjektivs führt in dem Fall dazu, dass der Sprecher den Grad seiner Assoziationen, die er mit der bezeichneten Eigenschaft (hier: das Lautsein) verbindet, auf einer entsprechenden Skala nach oben verschiebt. Demzufolge sollte es in einer experimentellen Untersuchung möglich sein, den durch ein Adjektiv ausgedrückten Grad mithilfe von Intensivierern zu manipulieren und die Versuchspersonen so in ihrer Antwort zu lenken, dass sie der Aussage einen höheren Intensitätsgrad zuweisen, als sie dies einem Satz ohne Intensivierer und somit ohne impliziten Vergleich zuschreiben würden: \textit{Das Kind ist mega laut} vs. \textit{Das Kind ist laut}.



%derivierte Adjektive: -ig und -lich Suffixe zur Bildung von Adjektiven





\section{Untersuchungsgegenstand}\label{abs:Untersuchungsgegenstand99}
Um Rückschlüsse auf die Expressivität der Intensivierer ziehen zu können, wurde eine Skalierungsaufgabe erstellt, bei der die Versuchspersonen die Ausdrücke, in einen Kontextsatz eingebettet, zu bewerten hatten. Dies beruhte auf der Idee, dass Intensivierer die durch das nachfolgende Element ausgesagte Eigenschaft verstärken und sich der Grad dieser Verstärkung auf eine kontinuierliche, längs angeordnete Skala abbilden lässt.


Infolgedessen wurden 18 Testsätze,\footnote{Die für die Auswertung relevanten Testsätze werden anschließend auch als Items bezeichnet.} bestehend aus verschiedenen Nomen-Adjektiv-Kombinationen (wie z.\,B. \textit{Der Berg ist hoch} oder \textit{Der Flur ist lang}), konstruiert, die mit den in der Korpusstudie berücksichtigten expressiven Intensivierern (n = 16) gekreuzt wurden: \textit{arg}, \textit{derb}, \textit{furchtbar}, \textit{fürchterlich}, \textit{hammer}, \textit{krass}, \textit{mega}, \textit{ober}, \textit{sau}, \textit{scheiß}, \textit{schrecklich}, \textit{super}, \textit{unheimlich}, \textit{verdammt}, \textit{voll} und \textit{wahnsinnig}. Neben dem deskriptiven Intensivierer \textit{sehr} (n = 1) wurde ebenso die unmarkierte Variante des Satzes (n = 1) als Kontrollstufe in das Experiment miteingebunden. Dadurch konnte überprüft werden, wie sich Sätze mit expressivem Intensivierer relativ zu Sätzen mit dem deskriptiven Intensivierer \textit{sehr} oder unmodifizierten Sätzen verhalten. In diesem Kontext wurde \citet[133]{Gutzmann.2019a} folgend vermutet, dass diese beiden Stufen auf der Skala insgesamt niedriger verortet werden als die expressive Stufe.


Ebensolche Evaluationsexperimente, bei denen Versuchspersonen Gegenstände im Hinblick auf bestimmte Merkmale miteinander vergleichen, sind aus methodischer Sicht allerdings nicht ganz unproblematisch. So ist davon auszugehen, dass die Semantik der verwendeten Ausdrücke zwangsläufig -- wenn auch für die Versuchsperson unbewusst -- die Bewertung beeinflusst. Die Beeinflussbarkeit trifft auf die Klasse der graduierbaren Adjektive in besonderem Maße zu, gelten diese aufgrund ihrer Vagheit doch gemeinhin als stark kontextabhängig \citep[vgl.][11f.]{Bierwisch.1987a}. Um dieses Problem zu umgehen und zu gewährleisten, dass sich die erforderlichen Nomen und Adjektive nicht unerwünscht auf die Antwort der Versuchspersonen auswirken, wurde in diesem Experiment auf bedeutungsleere Pseudowörter zurückgegriffen. Diese befolgten zwar die phonotaktischen Regeln des Deutschen,\footnote{Nach dem Kriterium der sprachspezifischen Phonotaktik lassen sich phonotaktisch legale und illegale Pseudowörter unterscheiden: \textit{Zunsel} vs. \textit{Kxzefurl}. Bei ersteren spricht {\citet[5]{Juergenson.2009}} von einer \glqq{}zufälligen Lücke\grqq{} im lexikalischen Inventar einer Sprache.} waren aber in dem Sinne semantisch leer, dass sie keine existenten deutschen Wörter sind.


Pseudowörter kommen bspw. bei sogenannten Jabberwocky-Sätzen zum Einsatz. Diese Sätze dienen vor allem in der Neurolinguistik dazu, der kognitiven Verarbeitung von Syntax ohne Einfluss von Semantik auf den Grund zu gehen. Dabei werden neben elektrophysiologischen Verfahren wie Elektroenzephalografie (EEG) und behavioralen Verfahren wie Self-Paced Reading ebenso bildgebende neuroanatomische Methoden wie funktionelle Magnetresonanztomografie (fMRT), Magnetenzephalografie (MEG) oder Positronen-Emissions-Tomografie (PET) angewandt (vgl. u.\,a. \citealt[]{Mazoyer.1993}; \citealt[]{Hahne.2001}; \citealt[]{Humphries.2006}; \citealt[]{Kharkwal.2014}). Grundsätzlich lassen sich bei Pseudowörtern zwei Arten unterscheiden: einerseits diejenigen, deren Wortursprung transparent ist und die Ähnlichkeit zu existenten Wörtern haben (z.\,B. \textit{Elefons} und \textit{Tamute}, ähnlich wie  \textit{Elefant} und \textit{Tomate}), und andererseits solche, die lexikalisch opak sind und somit keinerlei Ähnlichkeit mit existenten Wörtern aufweisen (z.\,B. \textit{Hützel} und \textit{Zünums}) {\citep[vgl.][1420]{Raettig.2008}}.


Um im Sinne der Zielsetzung Priming durch Ähnlichkeit mit realen deutschen Wörtern zu vermeiden und die Gradeinschätzung ausschließlich durch den Intensivierer zu evozieren, wurden in diesem Experiment legale opake Pseudowörter in Form von Nomen und Adjektiven eingesetzt. Denn nur durch opake Pseudowörter kann semantischer wie lexikalischer Zugriff vollständig verhindert werden, wie \citet[1426]{Raettig.2008} betonen: \glqq{}[O]nly opaque pseudowords reliably prevent successful lexical access and consequently inhibit subsequent semantic processing since they cannot be related to any existing entry in the mental lexicon.\grqq{} Bei transparenten Pseudowörtern ist die semantische und lexikalische Zugänglichkeit dagegen abhängig von der zugrunde liegenden Struktur der Wörter. Durch welche Vorgehensweise die im Experiment verwendeten Pseudowörter generiert und selektiert wurden, wird im Anschluss ausführlich geschildert.\footnote{Generierung und Auswahl der Pseudowörter erfolgte in Kooperation mit meinem Kollegen Fabian Ehrmantraut, der in seinem Forschungsprojekt dieselben {Pseudonomen verwendet}.}


\subsection{Generierung der Pseudowörter}\label{abs:Generierung}
Die in diesem Experiment verwendeten Pseudowörter basierten auf realen deutschen Wörtern, die unter Zuhilfenahme der seit 2007 online frei verfügbaren lexikalischen Datenbank dlexDB (\citealt[]{Heister.2011}) extrahiert wurden. Diese wurde deswegen ausgewählt, weil sie mit mehr als 100 Millionen Textwörtern und 290\,000 verschiedenen Wortformen erheblich umfangreicher ist als die in der Linguistik als Standard geltende Datenbank CELEX ({\citealt[]{Baayen.1996}}), die auf dem Mannheimer Korpus des Instituts für deutsche Sprache beruht und sich aus lediglich etwa 6 Millionen Textwörtern zusammensetzt {\citep[vgl.][10]{Heister.2011}}.
%Vor dem Hintergrund des erforderlichen Umfangs eines Korpus verweisen \citet[12]{Heister.2011} auf \citet[980]{Brysbaert.2009}, die bemerken, dass 16 bis 30 Millionen Wörter zu zuverlässigen Häufigkeitsschätzungen führen.
Nach Auffassung der Entwickler von dlexDB liegt das größte Defizit von {CELEX} in der einseitigen und veralteten Textauswahl, wohingegen die dlexDB zugrunde liegenden Textsorten, deren Korpusgrundlage das Digitale Wörterbuch der deutschen Sprache [DWDS] (\citealt[]{Geyken.2007}) ist, wesentlich aktueller und ausbalancierter seien {\citep[vgl.][11]{Heister.2011}}.\footnote{Zu diesem Schluss kommt auch \citet[50]{Gmoser.2013}, die in ihrer Arbeit die Worthäufigkeitsangaben der Korpusdatenbanken \textsc{DeReKo} und CELEX miteinander vergleicht.}


Als Vorbereitung des Datenexports wurde mittels der Filterabfrage des Webinterfaces von dlexDB\footnote{\citet[]{dlex.2011}: \textit{http://www.dlexdb.de/query/kern/typposlem/} [letzter Aufruf: 25.05.2021].} systematisch nach allen zweisilbigen Nomen und Adjektiven gesucht, deren Treffer in frequenziell abnehmender Reihenfolge angezeigt wurden. Diese Listen wurden exportiert und die häufigsten Ausdrücke anschließend in den Pseudowort-Generator \textit{Wuggy} (\citealt[]{Keuleers.2010}) eingespeist. Dies geschah vor dem Hintergrund der Annahme, dass die graphische und phonetische Struktur frequenter Ausdrücke vermutlich als am natürlichsten empfunden wird, sie also eine in hohem Maße natürliche Abfolge von Buchstaben und Lauten aufweisen. Bei {\textit{Wuggy}} handelt es sich um einen kostenfreien, computerbasierten Pseudowort-Generator, der für verschiedene Sprachen (z.\,B. Englisch oder Französisch) erhältlich ist. Mit dessen Hilfe lassen sich durch einen computerbasierten Algorithmus Pseudowörter nach individuellen, selbstdefinierten Kriterien entwickeln. In diesem Experiment wurden sowohl für die Nomen als auch für die Adjektive folgende Kriterien festgelegt: Die generierten Wörter sollten zweisilbig sein und möglichst wenige orthografische Nachbarn aufweisen, um eine Nähe zu existierenden Wörtern zu vermeiden. So handelt es sich bei \textit{Wonne} oder \textit{Tonne} bspw. um direkte Wortnachbarn des Nomens \textit{Sonne}, die sich in bloß einem Graphem voneinander unterscheiden. Für die Pseudoadjektive wurde als weiteres Kriterium definiert, dass deren Stämme über typische Adjektivendungen verfügen, was dazu diente, den Adjektivstatus der Wörter für die Versuchspersonen salienter zu machen. Hierfür wurden die adjektivbildenden Suffixe \textit{-ig} und \textit{-lich} ausgewählt, die im Deutschen als frequent und produktiv gelten (vgl. {\citealt[530]{Wellmann.1998}}; {\citealt[298]{Fleischer.2012}}).


Auf dieser Grundlage wurden von \textit{Wuggy} pro Nomen und Adjektiv fünf Wortvarianten erzeugt, die wiederum manuell selektiert wurden. Im Zuge der Selektion wurden zunächst doppelte, homographe und homophone Pseudowörter aussortiert, die lexikalischen Bezug zu realen Wörtern hatten (z.\,B. \textit{Löste} oder \textit{Wonder}). Pseudowörter, deren Bestandteile im Deutschen als eigenständige Lexeme existieren, wurden ebenfalls nicht verwendet (z.\,B. \textit{Sargtuch}). Im zweiten Schritt wurden die verbliebenen Pseudowörter in Hinsicht auf potenzielle orthografische Nachbarn gefiltert: Je weiter das Pseudowort von existenten Wörtern entfernt war, je weniger orthografische Nachbarn es durch den Austausch von Graphemen also hatte, desto besser eignete es sich für das Experiment. Demzufolge war \textit{Masen} mit einer Nachbarschaftsdichte von 17 Wörtern weniger gut geeignet (vgl. \textit{Rasen}, \textit{Vasen}, {\textit{Magen} etc.}) als \textit{Zotzer} mit null orthografischen Nachbarn.


Durch dieses sorgfältige Vorgehen wurden die Pseudoadjektive auf einen Kandidaten pro Wort und die Pseudonomen auf zwei Kandidaten pro Wort reduziert. Der Grund für den Unterschied lag in der Befürchtung, dass die Pseudonomen möglicherweise mehr Assoziationen zu realen Wörtern wecken und der Ausschuss demnach vermutlich größer sein würde als der der Pseudoadjektive. Somit blieben letztlich 80 Pseudoadjektive (40\,x \textit{-ig} und {40\,x \textit{-lich}}) und 100 Pseudonomen übrig, deren Buchstabenabfolge möglichst wenig Ähnlichkeit zu existierenden Wörtern zeigte. Um auszuschließen, dass die verbliebenen Pseudowörter mit realen Wörtern assoziiert werden, wurden vor dem Hauptexperiment für beide Pseudowort-Gruppen gemäß einem in der Literatur häufig beschriebenen Vorgehen (bspw. \citealt[668]{Kemp.2009}; \citealt[458]{Riedl.2014}) schriftliche Vortests durchgeführt, die im nächsten Abschnitt näher erläutert werden.


\subsection{Auswahl der Pseudowörter}\label{abs:Auswahl}
\subsubsection{Pseudoadjektive}\label{abs:Auswahl-Adj}
Im Rahmen des Vortests der Pseudoadjektive wurde eine Tabelle, bestehend aus Pseudoadjektiv und Freitext, per E-Mail an zehn freiwillige Versuchspersonen geschickt, mit der Aufforderung, die mit den Wörtern verbundenen Assoziationen anzugeben. Bei der Personenstichprobe handelte es sich um muttersprachliche Studierende des Faches Germanistik der Universität des Saarlandes, die für die Teilnahme keine Aufwandsentschädigung erhielten. Die Dauer des Vortests betrug etwa fünf Minuten. Die Ergebnisse wurden ohne personenbezogene Daten gesammelt und manuell auf Übereinstimmungen hin überprüft: Wenn zwei oder mehr Personen die gleiche Assoziation angaben, wurde das Pseudowort vom Experiment ausgeschlossen. Als gleiche Assoziation wurden diejenigen Wörter gewertet, die entweder identisch waren oder aus dem gleichen Themengebiet stammten (bspw. \textit{Schule}\,--\,\textit{Schüler}) bzw. den gleichen Wortstamm aufwiesen (bspw. \textit{Geifer}\,--\,\textit{geifern}). Nach Auswertung des Vortests blieben folgende 18 Pseudoadjektive übrig: \textit{fongig}, \textit{göcktig}, \textit{laukrig}, \textit{ontig}, \textit{ormig}, \textit{röngig}, \textit{rümmig}, \textit{wübig}, \textit{zänsig} und \textit{zarwig} sowie \textit{bühllich}, \textit{galklich}, \textit{kieflich}, \textit{ralklich}, \textit{schühnlich}, \textit{speiglich}, \textit{waumlich} {und \textit{wosplich}}.



\subsubsection{Pseudonomen}\label{abs:Auswahl-Nom}
Die Pseudonomen machten vor dem Assoziationstest einen zusätzlichen Schritt erforderlich, nämlich die Bestimmung des Artikels. Dieser war vonnöten, um im Hauptexperiment einen grammatischen Satz bilden zu können. Aus diesem Grund wurden für die Pseudonomen zwei unabhängige Versuchspersonengruppen rekrutiert, die für die Teilnahme nicht entlohnt wurden. Die Aufgabe der Versuchspersonen aus {Gruppe 1} bestand darin, jedem Pseudonomen über das Verfahren der Mehrfachauswahl nach eigener Einschätzung der Natürlichkeit einen der drei Definitartikel des Deutschen zuzuweisen: \textit{der Zünums}, \textit{die Zünums} oder \textit{das Zünums}. Durch dieses Vorgehen wurde sichergestellt, dass die Kombination aus Artikel und Pseudonomen im Hauptexperiment als möglichst natürlich wahrgenommen wird.
%- eventuell konkretisieren + und keinerlei negative Auswirkungen im eigentlichen Experiment hat.
Dieser Test wurde mittels der frei zugänglichen Umfrage-Applikation \textit{LimeSurvey} ({Version 3.22.9+200317})\footnote{\citet[]{LimeSurvey.2012}: \textit{https://www.limesurvey.org/de/} [letzter Aufruf: 03.05.2021].} erstellt. Die Items waren randomisiert und wurden den Versuchspersonen daher in der gleichen zufälligen Reihenfolge präsentiert. Für die Teilnahme erhielten 13 freiwillige deutsche Muttersprachler:innen -- Mitarbeiter:innen des Fachbereichs Germanistik der Universität des Saarlandes -- per E-Mail einen Link zugesandt, über den sie zum Vortest kamen. Diesen konnten sie daraufhin individuell online durchführen, ohne dass personenbezogene Daten erhoben wurden. Die Dauer des Vortests betrug im Schnitt zehn Minuten. Ausgewählt wurde schließlich derjenige Artikel, der im Vortest die meisten Übereinstimmungen hatte. Im Falle eines Gleichstands (n = 2) wurde der Artikel per Zufallsprinzip zugewiesen.


Im zweiten Vortest wurden die Versuchspersonen aus Gruppe 2 dazu aufgefordert, in einem dafür vorgesehenen Textfeld die Assoziationen anzugeben, die sie mit dem jeweiligen Nominalausdruck, d.\,h. der zuvor ermittelten Kombination aus Artikel und Pseudonomen, verbinden. Auch diese Umfrage wurde mit \textit{LimeSurvey} erstellt, woraufhin weitere freiwillige deutsche Muttersprachler:innen -- Studierende des Faches Germanistik der Universität des Saarlandes -- den Link per E-Mail zugesandt bekamen. Vollständig teilgenommen haben 28 Personen,\footnote{Die divergierenden Teilnehmerzahlen bei den Vortests kamen dadurch zustande, dass jeweils eine unterschiedliche Anzahl von Personen ihr Interesse an der Teilnahme bekundete. Dieses Ungleichgewicht wurde bei der anschließenden Selektion der Wörter berücksichtigt.} wobei auch hier keine personenbezogenen Daten erhoben wurden. Die Dauer des Vortests betrug ca. zehn Minuten.\footnote{Neben den Personen, die den Vortest bis zum Ende durchgeführt haben, gab es ebenso Drop-outs (n = 6), die den Test zwar begonnen, ihn jedoch mittendrin abgebrochen haben, was zu einem leichten Ungleichgewicht bei der Gesamtzahl der Antworten pro Nomen führte. Diese Angaben wurden beim maximalen Prozentsatz je Pseudonomen einkalkuliert.} Erneut wurden diejenigen Pseudowörter aussortiert, die bei zwei oder mehr Personen die gleichen Assoziationen weckten. Die aus der anschließenden Selektion resultierenden Nominalausdrücke (n = 18) sahen wie folgt aus: \textit{der Cobe}, \textit{der Conen}, \textit{der Darster}, \textit{der Gedalst}, \textit{der Jordruck}, \textit{der Nümig}, \textit{der Pulpart}, \textit{der Wonger}, \textit{der Zegeff}, \textit{der Zotzer}, \textit{die Itne}, \textit{die Itse}, \textit{das Genavs}, \textit{das Hensuln}, \textit{das Tusern}, \textit{das Zafche}, \textit{das Zeseet} und \textit{das Zülugs}.\footnote{Auch wenn in Bezug auf die Verteilung des definiten Artikels ein ausbalanciertes Verhältnis angestrebt wurde, war dies aufgrund der Ergebnisse des Assoziationstests nicht möglich. Daher mussten hier Abstriche zugunsten der Assoziationslosigkeit der Ausdrücke in Kauf genommen werden, die für das Experiment als bedeutsamer erachtet wurde.}


%Dabei handelt es sich um 18 Ausdrücke, deren Genus sich aus 10\,x Maskulinum, 2\,x Femininum und 6\,x Neutrum zusammensetzte.


%Auch wenn hier ein ausbalancierteres Verhältnis angestrebt wurde, war dies aufgrund der Ergebnisse des Assoziationstests unglücklicherweise nicht möglich. Infolgedessen mussten an dieser Stelle Abstriche zugunsten der Assoziationslosigkeit der Ausdrücke in Kauf genommen werden, die für das Experiment als bedeutsamer erachtet wurde. 
Als Ergebnis der schrittweisen Auswahl blieben letztlich nur solche Pseudowörter übrig, die keinerlei einheitliche Assoziationen bei den Versuchspersonen der Vortests hervorriefen. Auf diese Weise konnte gewährleistet werden, dass die Versuchspersonen des Hauptexperiments in keiner Weise durch semantische oder lexikalische Eigenschaften der Pseudowörter beeinflusst werden und der jeweils zugeschriebene Skalenwert ausschließlich durch die Gradbedeutung des Intensivierers zustande kam. Die Versuchspersonen der Vortests nahmen nicht am Hauptexperiment teil.


\subsection{Untersuchungsmaterial}\label{abs:Untersuchungsmaterial}
Das Untersuchungsmaterial des Hauptexperiments basierte auf den durch die Vortests gewonnenen Ausdrücken. Beispiele für die im Experiment abgefragten Kopulasätze mit verschiedenen Kombinationen aus Pseudonomen, Intensivierern und Pseudoadjektiven sind in \ref{items0} gelistet.

\ea   \label{items0}
\ea Der Nümig ist arg fongig.
\ex Der Wonger ist mega zänsig.
\ex Die Itse ist schrecklich rümmig.
\ex Das Zülugs ist sehr wosplich.
\ex Das Zafche ist göcktig.
\z
\z

\noindent Untersucht wurde eine unabhängige Variable mit 18 Stufen: 16 expressive Intensivierer, 1 deskriptiver Intensivierer sowie 1 unmodifizierte Variante. Bei den Pseudonomen und Pseudoadjektiven handelte es sich um feste Kombinationen, die über die Listen hinweg nicht variiert wurden. Exemplarisch bedeutet dies, dass \textit{Der Nümig} immer zusammen mit \textit{fongig} auftrat, sich je Liste jedoch der Intensivierer unterschied. Demnach rotierte jede feststehende Kombination aus Pseudonomen und Pseudoadjektiv (n = 18) im Experiment über die Intensivierer, wodurch für jeden Intensivierer gleich viele Datenpunkte gesammelt werden konnten. %Somit tauchte jede Pseudonomen-Pseudoadjektiv-Kombination im Experiment genau ein Mal mit einem anderen Intensivierer bzw. in einem unmodifizierten Satz auf.
Die abzufragenden Items wurden nach dem lateinischen Quadrat auf 18 Versuchspersonengruppen verteilt. Somit hatte jede Person eine dieser 18 Listen mit jeweils 18 Testsätzen zu bewerten. Aufgrund des gewählten Within Subjects-Designs sah jede Versuchsperson jedes mögliche Item, d.\,h. jeden Intensivierer sowie die Nullvariante, exakt ein Mal. Um Positionseffekte zu vermeiden, wurde die Abfolge der Items pro Liste randomisiert, sodass die Testsätze im Experiment in der gleichen zufälligen Reihenfolge präsentiert wurden.


Ergänzend zu den Items wurden vier Übungssätze mit gleicher syntaktischer Struktur konzipiert, an denen sich die Versuchspersonen anfangs mit der Aufgabe vertraut machen und ein Gespür für den Umgang mit den Pseudowörtern und Skalen entwickeln konnten. Dass es sich um Trainingssätze handelt, wurde nicht kommuniziert. Wie in den Items kamen auch in den Übungssätzen Intensivierer vor. Dabei wurde auf Ausdrücke zurückgegriffen, die nicht Gegenstand der Korpusstudie waren (z.\,B. \textit{außerordentlich} oder \textit{extrem}). Die Beschreibung der Objekte erfolgte zudem mithilfe derjenigen Pseudowörter, die infolge der Assoziationstests ausgesiebt wurden. Vom Gebrauch von Distraktoren wurde abgesehen: Da die Versuchspersonen die Intensivierer relativ zueinander in Beziehung setzen und den durch sie ausgedrückten Grad an Skalarität und Sprechereinstellung bewerten sollten, wurden Ablenkersätze als überflüssig angesehen.
%Die Übungssätze sahen wie folgt aus: 
%
%\ex.   \label{uebung1}
%\a. Der Zälper ist ganz fietslich.
%\b. Die Hünnel ist recht zücklich.
%\b. Das Länen ist extrem röngig.
%\b. Der Pites ist außerordentlich schnausig.





\section{Untersuchungsdurchführung}\label{abs:Untersuchungsdurchführung155}
Jede der Versuchspersonen hatte pro Item drei Aufgaben zu erledigen. Dazu wurden ihnen am Computer kurze und syntaktisch einfach konstruierte Sätze präsentiert, die verschiedene Objekte beschrieben. Wie bereits ausgeführt, erfolgte die Beschreibung der Objekte durch die mittels \textit{Wuggy} generierten Pseudowörter, wie z.\,B. \textit{Der Nümig ist sehr fongig}. Im ersten Schritt wurden die Personen darum gebeten, unter Zuhilfenahme eines Schiebereglers auf einer kontinuierlichen Skala den Ausprägungs- bzw. Skalaritätsgrad der im Satz beschriebenen Eigenschaft einzuschätzen. Übertragen auf das Beispiel also, wie fongig der Nümig ist. Dabei war der Regler durch Klicken oder Ziehen mit der Maus auf diejenige Skalenposition zu setzen, die ihrer Einschätzung nach dem Grad der Eigenschaft am besten entspricht. Zur Illustration ist die Skala von {Aufgabe I} in Abbildung \ref{Aufg1} dargestellt. Darüber hinaus findet sich im {Anhang \ref{Screen}} in {Abbildung \ref{Screen1}} ein Bildschirmfoto mit einer beispielhaften Komplettansicht des angeführten Items.



%
\begin{figure}
  \includegraphics[width=\textwidth]{Aufg1_22}
  \caption[Darstellung der in Experiment 1 verwendeten Bewertungsskala zu Aufgabe I] {Darstellung der in Experiment 1 verwendeten Bewertungsskala zu Aufgabe I -- Skalarität.} 
  \label{Aufg1}
\end{figure}

%\noindent Die Skalenrange reichte von 0 bis 100 Punkte ohne dass die Werte eingeblendet wurden. Vom Anzeigen sowohl des Werts der Skalenendpunkte, von festgelegten Skalenmarken wie auch der gewählten Skalenposition wurde abgesehen, um unerwünschte Priming-Effekte zu umgehen. So wurde bspw. die Möglichkeit in Betracht gezogen, dass Binnendifferenzierungen wie Skalenmarken oder segmentale Abstufungen die Versuchspersonen dazu verleiten würden, alle Ausdrücke gleichermaßen um ebendiese Marke(n) herum anzuordnen, statt sich bei der Beurteilung gänzlich auf ihre persönliche Intuition zu verlassen.
\noindent Die zu Beginn farblose Skala war am linken Rand mit `normal' und am rechten Rand mit `maximal' beschriftet. Diese neutralen Bezeichnungen wurden aus der Intention heraus gewählt, numerische Skalenbeschriftungen (bspw. $0$ und $100$) oder dimensional gefärbte Ausdrücke wie \textit{Intensität} oder \textit{Stärke} zu vermeiden und unerwünschte Priming-Effekte zu verhindern. Dadurch sollten die Versuchspersonen so wenig wie möglich bei ihrer Wahrnehmung bzw. der Wahl der Reglerposition beeinflusst werden. Wurde der Schieberegler betätigt, so färbte sich der Bereich vom Nullpunkt bis zum ausgewählten Punkt hin farbig ein, mit dem Ziel den Versuchspersonen die zugeschriebene Intensität visuell zu verdeutlichen. Links neben der Bewertungsskala stand die Aufgabenstellung, die im Wortlaut dem jeweiligen Itemsatz angepasst war. Exemplarisch also: \textit{Geben Sie an, wie fongig der Nümig ist}.


Die zweite Aufgabe, die eine optisch beinahe identische Skala inkludierte, erfragte den im Satz ausgedrückten Grad an Sprechereinstellung. Diese Aufgabe befand sich auf der gleichen Webseite unmittelbar unter der ersten. Die Skala von {Aufgabe II} befindet sich in {Abbildung \ref{Aufg2}}.



%
\begin{figure}
  \includegraphics[width=\textwidth]{Aufg2_22}
  \caption[Darstellung der in Experiment 1 verwendeten Bewertungsskala zu Aufgabe II] {Darstellung der in Experiment 1 verwendeten Bewertungsskala zu Aufgabe II -- Sprechereinstellung.} 
  \label{Aufg2}
\end{figure}

\noindent Die Skala zu Aufgabe II war zu Beginn ebenfalls farblos und entsprach der ersten in Gestalt und Betätigung. Auch hier war der Schieberegler intuitiv auf die zutreffende Skalenposition zu bewegen. Lediglich die Beschriftungen und die Aufgabenstellung unterschieden sich: Diese Skala war links mit `gar nicht' und rechts mit `maximal' gelabelt und erfragte den ausgedrückten Grad an Sprechereinstellung: \textit{Geben Sie an, wie stark der Sprecher seine persönliche Haltung zum beschriebenen Sachverhalt ausdrückt}. In der Aufgabenstellung wurde vom Gebrauch jeglicher Intensivierer (bspw. \textit{...\,, wie sehr der Sprecher\,...}) abgesehen und stattdessen auf das Adjektiv \textit{stark} zurückgegriffen.



%Die Skalenrange reichte jeweils von 0\,\% bis 100\,\%, ohne dass die Prozentwerte eingeblendet wurden. Vom Anzeigen sowohl des Werts der Skalenendpunkte wie auch der gewählten Skalenposition wurde abgesehen, um Priming-Effekte jeglicher Art zu umgehen. Die zu Beginn farblose Skala von Aufgabe I, die den Grad an Skalarität erfragte, war am linken Rand mit 'normal' und am rechten Rand mit 'maximal' beschriftet. Diese neutralen Bezeichnungen wurden aus der Intention heraus gewählt, dimensional gefärbte Ausdrücke wie \textit{Intensität} oder \textit{Stärke} zu vermeiden und die Versuchspersonen so wenig wie möglich bei ihrer Antwort bzw. der Wahl der Reglerposition zu beeinflussen. Wurde der Schieberegler betätigt, so färbte sich der Bereich vom Nullpunkt bis zum ausgewählten Punkt hin farbig ein, mit dem Ziel, den Versuchspersonen die zugeschriebene Intensität visuell zu verdeutlichen. Links neben der Skala befand sich die Aufgabenstellung, die im Wortlaut dem Itemsatz angepasst war: \textit{Geben Sie an, wie fongig der Nümig ist}. Zur Illustration ist die Skala von Aufgabe I in Abbildung 
%\ref{Aufg1} dargestellt.


%
%%
%\begin{figure}
%
%  \includegraphics[width=1.005\textwidth]{Aufg1_22}
%  \caption[Darstellung der im Brückenexperiment verwendeten Bewertungsskala zu Aufgabe I] {Darstellung der im Brückenexperiment verwendeten Bewertungsskala zu Aufgabe I -- Skalarität.} 
%  \label{Aufg1}
%\end{figure}
%
%\noindent Die Skala von Aufgabe II, die auf den Grad an Sprechereinstellung abzielte, war zu Beginn ebenfalls farblos und entsprach der ersten in Gestalt und Betätigung. Lediglich die Beschriftungen und die Aufgabenstellung unterschieden sich: Diese Skala war links mit 'gar nicht' und rechts mit 'maximal' gelabelt und erfragte, wie stark der Sprecher seine persönliche Haltung zum beschriebenen Sachverhalt ausdrückt. In der Aufgabenstellung wurde vom Gebrauch jeglicher Intensivierer abgesehen und stattdessen auf das Adjektiv \textit{stark} zurückgegriffen.
%%Um vom Gebrauch jeglicher Intensivierer in der Aufgabenstellung abzusehen, wurde an der Stelle das Adjektiv \textit{stark} verwendet. 
%Die Skala von {Aufgabe II} befindet sich in {Abbildung \ref{Aufg2}}.



%%
%\begin{figure}
%
%  \includegraphics[width=1.01\textwidth]{Aufg2_22}
%  \caption[Darstellung der in Experiment 1 verwendeten Bewertungsskala zu Aufgabe II] {Darstellung der in Experiment 1 verwendeten Bewertungsskala zu Aufgabe II -- Sprechereinstellung.} 
%  \label{Aufg2}
%\end{figure}

Die dritte Aufgabe war mithilfe einer komparativen Konstant-Summen-Skala zu erfüllen und erfragte die evaluative Einstellung des Sprechers zum bezeichneten Sachverhalt in Form von Abweichung von einer neutralen Mitte. Dabei wurden die Versuchspersonen dazu aufgefordert, mit einem festgelegten, virtuellen Geldbetrag zu wetten, ob die Satzaussage positiv, neutral oder negativ wertend zu verstehen ist bzw. ob der Sprecher dem Sachverhalt gegenüber positiv, neutral oder negativ eingestellt ist. Um die Aufgabe zu erledigen, sollten sie 100 Taler gemäß persönlicher Intuition restlos auf die zur Auswahl stehenden Felder `Positiv', `Neutral' und `Negativ' setzen. Diese Aufgabe wurde ebenso auf der gleichen Webseite direkt unter der zweiten angezeigt. Die im Experiment verwendete komparative Konstant-Summen-Skala zur Ermittlung der Einstellungspolarität ist mitsamt Aufgabenstellung in {Abbildung \ref{Aufg3} illustriert}.\footnote{Zu beachten ist, dass die Aufgabenstellung in der Abbildung aufgrund größtmöglicher Seitenskalierung verzerrt ist. Im Experiment umfasste diese zwei vollständig ausgefüllte Zeilen.}



%
\begin{figure}
  \includegraphics[width=\textwidth]{Aufg3_23}
  \caption[Darstellung der in Experiment 1 verwendeten Konstant-Summen-Skala zu Aufgabe III] {Darstellung der in Experiment 1 verwendeten Konstant-Summen-Skala zu Aufgabe III -- Einstellungspolarität.} 
  \label{Aufg3}
\end{figure}


\noindent Die noch einzusetzende Restmenge (`Verbleibend') sowie die Summe des bereits verteilten Geldes (`Gesamt') waren unter dem letzten Auswahlfeld eingeblendet. Die Taler mussten letztlich vollständig und in ganzen Zahlen auf die Felder aufgeteilt werden. Der Mehrwert einer komparativen Konstant-Summen-Skala liegt darin, dass die Versuchspersonen qua Abfragedesign dazu gezwungen sind, die Talersumme restlos aufzubrauchen und sie erst danach zur Beantwortung des nächsten Items übergehen können. Wie sie die Taler auf die Auswahlfelder verteilen, steht ihnen vollkommen frei. Durch dieses Vorgehen erhält man feinkörnige Abstufungen zwischen den Optionen.


Bevor das Experiment begann, wurden die Versuchspersonen in einer ausführlichen Instruktion samt beispielhafter Skalenabbildungen über den Anlass, die Dauer und den Ablauf des Experiments sowie datenschutzrechtliche Aspekte informiert. Außerdem wurden sie darum gebeten, das Experiment an einem Endgerät mit physischer Tastatur vorzunehmen, um zu verhindern, dass die Skalen aufgrund unterschiedlicher Bildschirmgröße (Smartphone/Tablet vs. Notebook/Desktop-PC) nicht korrekt angezeigt oder nicht wie vorgesehen bedient werden können.\footnote{Zur Kontrolle wurde im Experiment die Bildschirmgröße der {Versuchspersonen erfasst}.} Im Anschluss an die Einleitung folgten zunächst die Übungssätze, die gemäß dem Vorgehen bei den Testsätzen einzeln abgefragt wurden und nach der Beantwortung nahtlos und ohne weitere Erklärung in die für die Auswertung relevanten {Items übergingen}.


%Bei den abhängigen Variablen handelt es zum einen sich um die Skalenwerte, die die Versuchspersonen den Testsätzen zuschreiben sollten, sowie um die Taler, die sie auf die persönliche Einstellung des Sprechers wetten sollten. 



\section{Versuchspersonen}\label{abs:Versuchspersonen1}
Am Experiment haben insgesamt 71 deutsche Muttersprachler:innen im Alter von 18 bis 40 Jahren teilgenommen, die über die Crowdsourcing-Plattform \textit{Clickworker}\footnote{\citet[]{Rozsenich.2020}: \textit{https://www.clickworker.de} [letzter Aufruf: 23.06.2021].} rekrutiert und (bemessen an dem zu der Zeit gültigen deutschen Mindestlohn) mit je 1,70\,Euro entlohnt wurden. Über das Ziel der Untersuchung wurden die Versuchspersonen nicht informiert.


Wie die Vortests wurde die Untersuchung mittels \textit{LimeSurvey} erstellt und von den Versuchspersonen individuell online durchgeführt. Diesen wurde eine der 18 Listen zugeteilt. Die Dauer des Experiments betrug im Schnitt etwa zehn Minuten. Die Teilnahme erfolgte völlig anonym ohne Erhebung personenbezogener Angaben wie Alter, Geschlecht oder Herkunft, sodass keinerlei Rückschlüsse auf die Identität der Personen gezogen werden {konnten}.


Angestrebt wurde die Teilnahme von 72 Versuchspersonen, um zu gewährleisten, dass von jedem Intensivierer sowie der Nullvariante pro Liste vier Datenpunkte und somit eine gleichmäßige Verteilung vorliegt. Durch ein Problem von \textit{Clickworker} konnten nach Abschluss des Experiments jedoch nur Daten von 71 Personen exportiert werden. Aus diesem Grund standen für die Auswertung letztlich in Summe 17 Listen mit Daten von je vier Versuchspersonen sowie eine Liste mit Daten von drei Versuchspersonen zur Verfügung, was ein leichtes Ungleichgewicht nach sich zog. Da einzelne Personen von der Analyse ausgeschlossen werden mussten und die verwendeten linearen Regressionsmodelle als tolerant gelten, wurde angenommen, dass das bestehende Ungleichgewicht nicht relevant und vernachlässigbar ist. Zudem lässt sich bei Auswertungen häufig nicht vermeiden, dass einzelne Versuchspersonen unberücksichtigt bleiben, was ebenfalls zu Datenungleichgewichten führen kann.




\section{Ausschlusskriterien}
Im Vorfeld der Untersuchung wurden verschiedene Kriterien definiert, deren Zutreffen den Ausschluss einzelner Versuchspersonen zur Folge hatte. Ein Ausschluss von Personen kann dann erforderlich sein, wenn sich Anhaltspunkte für eine absichtliche Manipulation der Ergebnisse finden und der Einbezug ihrer Antworten gegebenenfalls zu verzerrten Resultaten führen würde. Zu den Ausschlusskriterien gehörten vorrangig eine überschnelle, schier unrealistische Reaktionszeit und ein uniformes Antwortmuster, die zusammen auf ein oberflächliches Durchführen des Experiments {schließen ließen}.


Von ursprünglich 71 Versuchspersonen wurden die Daten von vier Personen ausgeschlossen, weil sich bei ihnen der Verdacht erhärtete, dass sie das Experiment in Teilen oder auch als Ganzes nicht sorgfältig bearbeitet hatten. Dies resultierte in einem vollkommen einheitlichen Antwortmuster, das keinerlei Variation hinsichtlich der Skalenbewertung aufzeigte. So wurden z.\,B. alle Intensivierer inklusive der unmarkierten Variante gleichermaßen am linken oder rechten Skalenrand positioniert. Gestützt wurde diese Vermutung durch eine Begutachtung der durchschnittlichen Reaktionszeit je Item: Betrug diese bei allen Personen über die Dauer des Experiments hinweg im Schnitt 22,07 Sekunden (Standardabweichung (Englisch: \textit{standard deviation}, hier und im Folgenden abgekürzt mit \textit{SD} = 15,8), lag diese bei den ausgeklammerten Personen deutlich darunter (bspw. bei 5,94 Sekunden ({\textit{SD} = 1,8)} oder 6,59 Sekunden ({\textit{SD} = 1,0))}. Diese Schnelligkeit scheint angesichts der Menge der zu bewältigenden Aufgaben pro Item (zwei Mal Skaleneinschätzung und Aufteilung der Taler mittels Konstant-Summen-Skala) unrealistisch. Stattdessen ist davon auszugehen, dass diese knappe Zeitspanne in der Praxis kaum dazu ausreicht, um den zu bewertenden Satz sowie die Aufgabenstellungen zu lesen und zu verarbeiten; vor allem dann nicht, wenn unbekannte Pseudowörter darin enthalten sind und der durch den Satz ausgedrückte Grad auf eine Skala abgebildet werden soll. Um im Sinne der verfolgten Fragestellung repräsentative Ergebnisse zu erzielen, wurde der Ausschluss dieser Versuchspersonen als gerechtfertigt angesehen.





\section{Ergebnisse}\label{abs:Ergebnisse238}
%In diesem Abschnitt werden die Ergebnisse des Experiments vorgestellt. Dieses wurde zum einen mit dem Ziel durchgeführt, die Gültigkeit der noch ungeklärten Hypothesen H1-Temp und H1-Freq aus Kapitel \ref{Korpus} zu überprüfen, nach denen die Dauer oder Häufigkeit des Gebrauchs eines Intensivierers abträglich für dessen Expressivität ist (vgl. \citealt[59]{Tobler.1868}; {\citealt[2]{Hauschild.1899}}). Zum anderen wurde die Korrektheit der theoretischen Annahme getestet, dass Expressivität gleichermaßen mit einem höheren Grad an {Skalarität (vgl. H3)} und Sprechereinstellung (vgl. H4) einhergeht (vgl. {\citealt[315]{Lang.1983}}; {\citealt[22]{Paradis.1997}}; \citealt[57]{Ortner.2014}; {\citealt[133]{Gutzmann.2019a}}). Im Kontext letzterer wurde ferner die Polarität der Sprechereinstellung in die Untersuchung miteinbezogen, die bei expressiven Intensivierern aufgrund ihrer Einstellungsbezogenheit nach Auffassung von {\citet[135]{Gutzmann.2019a}} vielmehr positiv oder negativ als neutral interpretiert wird ({vgl. H5}).


Dieses Teilkapitel setzt sich aus deskriptiv- und inferenzstatistischer Datenanalyse zusammen und beginnt mit der Begutachtung der Ergebnisse von {Aufgabe I}, die die Skalarität ermittelte. Daraufhin werden die Daten von {Aufgabe II}, die auf die Sprechereinstellung abzielte, inspiziert. Vor diesem Hintergrund wird überprüft, ob expressive Intensivierer wirklich mit einem höheren Grad an Skalarität (vgl. H3) und Sprechereinstellung {(vgl. H4)} einhergehen, wie nach \citet[22]{Paradis.1997} und \citet[133]{Gutzmann.2019a} sowie {\citet[315]{Lang.1983}} und \citet[57]{Ortner.2014} anzunehmen ist. Da im Sinne von {\citet[133]{Gutzmann.2019a}} die Möglichkeit besteht, dass die Variablen \textsc{Skalarität} und \textsc{Sprechereinstellung} insofern korrelieren, als die Expressivität eines Intensivierers auf dessen Skalaritätsgrad einwirkt, wird im weiteren Verlauf ebendiesem möglichen Zusammenhang nachgegangen. Den Abschluss macht eine Betrachtung der Ergebnisse, die im Hinblick auf die in {Aufgabe III} abgefragte Einstellungspolarität erzielt wurden. In diesem Kontext wird untersucht, ob sich die Expressivität eines Intensivierers signifikant auf die Einschätzung der Einstellungspolarität auswirkt. So wird die Sprechereinstellung, die expressive Intensivierer ausdrücken, {\citet[135]{Gutzmann.2019a}} zufolge vielmehr positiv oder negativ als neutral interpretiert (vgl. H5).


Die statistischen Auswertungen wurden mithilfe von R ({Version 1.1.456}, \mbox{\citealt[]{R.2018}}) und den Paketen \textit{lme4} (\citealt[]{Bates.2015}) und \textit{lmerTest} ({\citealt[]{Kuznetsova.2017}}) vorgenommen. Die Visualisierung der Daten erfolgte außerdem mit \textit{ggplot2} ({\citealt[]{Wickham.2016}}). In der Untersuchung werden die drei klassischen Signifikanzniveaus `signifikant' ({\textit{p} < 0,05 $\star$}), `sehr signifikant' (\textit{p} < 0,01 $\star$$\star$) und `hochsignifikant' (\textit{p} < 0,001 $\star$$\star$$\star$) angenommen.\footnote{Der Asterisk (*) dient dazu, den jeweiligen \textit{p}-Wert gemäß dem zutreffenden Signifikanzniveau zu kennzeichnen.}
Bei der inferenzstatistischen Auswertung werden jeweils die im Zuge der Datenmodellierung berechneten linearen gemischten Regressionsmodelle vorgestellt. Um den einleitenden Fragen hinsichtlich dem Einfluss von Frequenz und Persistenz Rechnung zu tragen, wurden für die im Rahmen der Korpusuntersuchung eingeteilten Frequenz- und Persistenzgruppen ebenfalls Modelle berechnet, die jeweils im Anschluss an Aufgabe I und {Aufgabe II} präsentiert werden. Dadurch wird getestet, ob eine hohe Frequenz (vgl. H1-Freq) oder ein persistenter Gebrauch (vgl. H1-Temp) tatsächlich abträglich für die Expressivität eines Intensivierers ist, wie \citet[2]{Hauschild.1899} und {\citet[59]{Tobler.1868}} postulieren. Weil letztere Berechnungen lediglich die Gruppe der expressiven Intensivierer betreffen, wurden die deskriptive und unmarkierte Stufe bei diesen Regressionsanalysen außen vor gelassen. Die für die beiden Faktoren Frequenz und Persistenz berechneten Modelle schließen sich jeweils an die Modelle zu Aufgabe I und Aufgabe II an. An dieser Stelle sei erwähnt, dass die Faktoren nicht unabhängig voneinander, sondern inhärent zusammenhängend sind. Dabei wird lediglich der Faktor Frequenz variiert, während der Faktor Persistenz in dem Sinne konstant bleibt, dass die drei neuen Intensivierer grundsätzlich mit einer niedrigen Frequenz auftreten.








\subsection{Aufgabe I -- Skalarität} \label{abs:VP-Vergleich1251}


\subsubsection{Vergleich der Stufen} \label{abs:I-Vergleich254}


In Tabelle \ref{SkaWe1} sind die gemittelten Skalenwerte, d.\,h. das arithmetische Mittel, sowie die Standardabweichungen, basierend auf 67 Datenpunkten pro Intensivierer bzw. der unintensivierten Variante, zu Aufgabe I aufgeführt. Die 16 expressiven Intensivierer werden zunächst zu einer Stufe zusammengefasst und mit den beiden Kontrollstufen kontrastiert. Die Tabelle bildet daher die Werte von drei Stufen (d.\,h. deskriptiv, expressiv und unmarkiert)  in alphabetischer Reihenfolge ab.

%
\begin{table}
%\setlength{\tabcolsep}{16pt}
\begin{tabularx}{.8\textwidth}{lYY}
\lsptoprule
%            & \textbf{Aufgabe I}  & \\
%            & \textbf{Skalarität} & \\
Stufe      & Skalenwert     &\textit{SD}\\
\midrule
deskriptiv & 63,90          & 17,9\\
expressiv  & 67,70          & 22,4\\
unmarkiert & 35,27          & 27,4\\
\lspbottomrule
\end{tabularx}
\caption[Gemittelte Skalenwerte der Stufen bei {Aufgabe I} von Experiment 1]{Gemittelte Skalenwerte und Standardabweichungen der Stufen bei Aufgabe I -- Skalarität (in alphabetischer Reihenfolge).}
\label{SkaWe1}
\end{table}


\noindent Die Werte aus Tabelle \ref{SkaWe1} weisen darauf hin, dass die expressiven Intensivierer nach Einschätzung der Versuchspersonen insgesamt eine höhere Skalenposition einnehmen als der deskriptive Intensivierer oder die unmarkierte Variante. So wurde die expressive Stufe diesbezüglich mit 67,70 Punkten auf der Skala am höchsten verortet, wohingegen der deskriptiven Stufe von den Versuchspersonen ein um etwa 3,80 Punkte niedrigerer Skalenwert zugewiesen wurde. Weit abgeschlagen dahinter befindet sich die unmarkierte Stufe, deren Wert sich von der expressiven Stufe um ca. 32,43 Punkte, von der deskriptiven Stufe um {ca. 28,63 Punkte} unterscheidet. Dies deutet erwartungsgemäß auf eine in der Skalarität abnehmende Rangfolge der {Stufen hin}.


Dieses Ergebnis ist in Abbildung \ref{SkalarBox} mittels Boxplots, auch Kastengrafik genannt, dargestellt. Die von der Box nach oben und unten reichenden Whiskers verdeutlichen die Varianz der Skalenbewertung und geben an, wie stark das Skalenspektrum von den Versuchspersonen genutzt wurde. Die Querstriche oben und unten symbolisieren das Maximum und Minimum der Whiskers und somit die Extremwerte, die noch keine Ausreißer sind, während die Einkerbung innerhalb der Box für den Median steht. Die beiden Quartile, aus denen die Box besteht, geben die Verteilung der Skalenwerte um den Median herum an. Darüber hinaus stellen die Kreise Ausreißer dar, deren Abstand mehr als anderthalb Interquartile von der Box entfernt liegen (vgl. \citealt[77f.]{Mittag.2016};; \citealt[579]{Wollschlaeger.2017}).


%
\begin{figure}
  \includegraphics[width=1\textwidth]{BoxSkalar_1}
	\caption[Gemittelte Skalenwerte der Stufen bei {Aufgabe I} von Experiment 1]{Gemittelte Skalenwerte der Stufen bei Aufgabe I -- Skalarität ({x-Achse =} Stufe, y-Achse = Skalenwert).}
  \label{SkalarBox}
\end{figure}


%
\subsubsection{Regressionsanalysen zum Vergleich der Stufen}


Mit dem Ziel die soeben beschriebenen deskriptivstatistischen Ergebnisse auf ihre Signifikanz hin zu überprüfen, wurden sie einer zusätzlichen inferenzstatistischen Auswertung unterzogen. Wie eingangs beschrieben, wurden die Daten mittels R und den Paketen \textit{lme4} (\citealt[]{Bates.2015}) und \textit{lmerTest} ({\citealt[]{Kuznetsova.2017}}) ausgewertet. Hierzu wurden lineare gemischte Modelle berechnet, die aus der unabhängigen Variable die abhängige Variable voraussagen. Im Rahmen der Regressionsanalyse für Aufgabe I wurde zunächst ein volles Modell mit den Effekten aller unabhängigen Variablen und deren Zwei-Wege-Interaktionen berechnet, die eventuell einen Einfluss auf die abhängige Variable hatten. Das volle Modell beinhaltete \textit{Fixed Effects} für den \textsc{Skalenwert}, die \textsc{Stufe}, die \textsc{Position}\footnote{\textit{Position} meint hier und im weiteren Verlauf die Position eines Items im Experiment.} sowie deren Zwei-Wege-Interaktionen. Der \textsc{Skalenwert} wurde dabei als metrische abhängige Variable modelliert, die \textsc{Stufe} und die \textsc{Position} als kategoriale Prädiktoren. Die Variable \textsc{Skalenwert} steht bei dieser Aufgabe für den Skalaritätsgrad, den die Versuchspersonen mit der jeweiligen Stufe oder dem jeweiligen Intensivierer verbunden haben. Um einen möglichen graduellen Effekt und Unterschiede zwischen den drei Stufen zu finden, wurde das Referenzlevel `Expressiv' jeweils mit `Deskriptiv' und `Unmarkiert' verglichen. Dazu wurden die Ausprägungsformen der mehrkategorialen Variable \textsc{Stufe} (d.\,h. deskriptiv, expressiv und unmarkiert) per Dummy-Codierung in zwei Schritten (d.\,h. `Expressiv' vs. `Deskriptiv' und `Expressiv' vs. `Unmarkiert') als $0$ und $1$ kodiert und somit nach binären Kontrasten gesplittet {\citep[vgl.][407, 442]{Sauer.2019}}. Dadurch wurden aus einer unabhängigen Variable mit drei Stufen zwei unabhängige Variablen mit je zwei Stufen. Zudem wurden im Modell versuchspersonen- und itemweise \textit{Random Intercepts} und \textit{Random Slopes} für die Variablen \textsc{Skalenwert}, \textsc{Stufe} und \textsc{Position} miteinbezogen.  Das bedeutet, dass im Modell ebenso potenzielle Unterschiede zwischen den Versuchspersonen und Pseudowörtern berücksichtigt wurden. Ausgehend vom vollen Regressionsmodell wurden durch das Prinzip der \textit{Backward Model Selection}, auch Rückwärtsselektion genannt, schrittweise alle Effekte getilgt, die den Modell-Fit nicht signifikant verbesserten. Begonnen wurde jeweils mit dem Effekt, der nach der Berechnung mit {\textit{lmerTest} (\citealt[]{Kuznetsova.2017})} den höchsten \textit{p}-Wert aufwies. Dieser wurde im nächsten Schritt beim neuen Modell entfernt und das Modell so sukzessive vereinfacht. Daraufhin wurden die beiden Modelle mithilfe von Likelihood-Quotienten-Tests bzw. Chi$^{2}$-Tests mit der Funktion \textit{anova} in R kontrastiert: Wenn das komplexere Modell die Daten nicht signifikant besser erklärte, wurde das einfachere Modell bevorzugt. Auf diese Weise konnte das optimale Regressionsmodell gefunden werden.


Nachdem durch dieses Vorgehen mehrere Modelle mittels Rückwärtsselektion und der \textit{anova}-Funktion in Bezug auf die statistische Signifikanz des Unterschiedes miteinander verglichen wurden, konnte ein finales Modell mit ausschließlich signifikanten Effekten abgeleitet werden. Dieses gibt den Einfluss der unabhängigen Variablen auf die abhängige Variable an. In der Tabelle, die das jeweilige finale Regressionsmodell abbildet, ist neben dem \textit{$\chi$$^{2}$}-Wert und \textit{p}-Wert der Standardfehler (englisch: \textit{standard error}, hier und im Folgenden abgekürzt mit \textit{SE}) angegeben, der die Streuung der Daten wiedergibt {\citep[vgl.][273]{Sauer.2019}}. Zudem ist der Regressionskoeffizient {(= \textit{Estimate}}) gelistet, der den \glqq{}Einfluss des Prädiktors\grqq{} ({\citealt[351]{Sauer.2019}}) {anzeige}. Die vollen und finalen Regressionsmodelle zu diesem Experiment sind im {Anhang \ref{Regr1}} separiert nach Aufgabe I (Skalarität), Aufgabe II (Sprechereinstellung) und {Aufgabe III} (Einstellungspolarität) gesammelt.


Das finale Modell enthält signifikante Haupteffekte für die Prädiktoren \textsc{Stufe} und \textsc{Position}. Diese Koeffizienten sind in {Tabelle \ref{LinReg0001}} zu finden.


%
\begin{table}
%\setlength{\tabcolsep}{9.375pt}
\begin{tabularx}{\textwidth}{Xccccr}
\lsptoprule
Prädiktor & \textit{Est.} & \textit{SE} & \textit{$\chi$$^{2}$} & \textit{p} & \\
\midrule
(Intercept) & \textcolor{white}{-}67,88 & 1,98 & 124,15 & < 0,001 & $\star$$\star$$\star$\\
Expressiv vs. Unmarkiert & -32,56 & 3,39 & \textcolor{white}{-}54,42 & < 0,001 & $\star$$\star$$\star$\\
Expressiv vs. Deskriptiv & \textcolor{white}{-}\textcolor{white}{-}-5,11 & 2,56 & \textcolor{white}{-}\textcolor{white}{-}3,74 & < 0,05 & $\star$ \\
Position & \textcolor{white}{-}\textcolor{white}{-}-1,70 & 0,63 & \textcolor{white}{-}\textcolor{white}{-}6,90 & < 0,01 & \multicolumn{1}{r}{$\star$$\star$} \\
\lspbottomrule
\end{tabularx}
\caption[Finales Regressionsmodell zu Aufgabe I von Experiment 1 (mit `Expressiv' als Referenzlevel)]{Finales Regressionsmodell zu Aufgabe I -- Skalarität (mit `Expressiv' als Referenzlevel).}
\label{LinReg0001}
\end{table}

\noindent Dem in Tabelle \ref{LinReg0001} dargelegten Modell nach haben zwei unabhängige Variablen einen hohen Einfluss auf die abhängige Variable, d.\,h. den Skalenwert: \textsc{Stufe} und \textsc{Position}. Zum einen ist der festgestellte Unterschied zwischen der expressiven und der unmarkierten Stufe signifikant {($\chi$$^{2}$\,(1) = 54,42}, {\textit{p} < 0,001}), sprich Sätzen mit expressivem Intensivierer wird im Allgemeinen eine signifikant höhere Skalenposition zugeschrieben als unintensivierten Sätzen. Zum anderen ist ein signifikanter Unterschied zwischen der expressiven und der deskriptiven Stufe auszumachen {($\chi$$^{2}$\,(1) = 3,74}, \textit{p} < 0,05). Somit werden neben unintensivierten Sätzen auch Sätze mit deskriptivem Intensivierer auf der Skala signifikant niedriger gerankt als Sätze mit expressivem Intensivierer. Daneben zeigt sich ein signifikanter Positionseffekt ({$\chi$$^{2}$\,(1) = 6,90}, {\textit{p} < 0,01}), der wie folgt zu interpretieren ist: Je später ein Item im Experiment auftritt, desto niedriger ist seine Bewertung auf der Skala. Dies deutet auf einen strategischen Gewöhnungseffekt hin, in dessen Folge der zu Beginn als stärker wahrgenommene Unterschied zwischen den Stufen im Laufe des Experiments nach und nach kleiner geworden ist.


In einer weiteren Analyse wurde das Referenzlevel der mehrkategorialen Variable \textsc{Stufe} auf `Unmarkiert' gesetzt, was bedeutet, dass `Unmarkiert' jeweils mit `Expressiv' und `Deskriptiv' kontrastiert wurde. Vor diesem Hintergrund drängt sich die Frage auf, ob die Kontrollstufe, d.\,h. die unmarkierten Sätze, prinzipiell mit einem signifikant niedrigeren Skalaritätsgrad assoziiert werden und sich dahingehend sowohl von der expressiven als auch der deskriptiven Stufe unterscheiden. Das methodische Vorgehen entsprach dem obigen: Wie zuvor wurde zunächst ein volles Regressionsmodell berechnet, in dem nach Rückwärtsselektion und Modellvergleich sukzessive all diejenigen Effekte weggelassen wurden, die sich nicht signifikant auf die abhängige Variable ausgewirkt haben. Die \textit{Random-Effects}-Struktur wurde mit Ausnahme des Referenzlevels der beiden binären Vergleiche nicht verändert. Das finale Modell beinhaltet signifikante Haupteffekte für die Prädiktoren \textsc{Stufe} und \textsc{Position}. Die berechneten Koeffizienten sind in {Tabelle \ref{LinReg0010}} {aufgeführt}.

%
\begin{table}
\begin{tabularx}{\textwidth}{Xccccr}
\lsptoprule
Prädiktor & \textit{Est.} & \textit{SE} & \textit{$\chi$$^{2}$} & \textit{p} & \\
\midrule
(Intercept) & 35,33 & 3,18 & 70,60 & < 0,001 & $\star$$\star$$\star$\\
Unmarkiert vs. Expressiv & 32,55 & 3,39 & 54,35 & < 0,001 & $\star$$\star$$\star$\\
Unmarkiert vs. Deskriptiv & 27,43 & 3,33 & 44,02 & < 0,001 & $\star$$\star$$\star$\\
Position & \textcolor{white}{-}-1,71 & 0,63 & \textcolor{white}{-}6,96 & < 0,01 & \multicolumn{1}{r}{$\star$$\star$} \\
\lspbottomrule
\end{tabularx}
\caption[Finales Regressionsmodell zu Aufgabe I von Experiment 1 (mit `Unmarkiert' als Referenzlevel)]{Finales Regressionsmodell zu Aufgabe I -- Skalarität (mit `Unmarkiert' als Referenzlevel).}
\label{LinReg0010}
\end{table}


\noindent Die in Tabelle \ref{LinReg0010} zusammengetragenen Werte erhärten die zu testende Annahme: Die unmarkierte Stufe unterscheidet sich nicht nur signifikant von der expressiven Stufe ($\chi$$^{2}$\,(1) = 54,35, \textit{p} < 0,001), sondern auch von der deskriptiven Stufe ({$\chi$$^{2}$\,(1) = 44,02}, \textit{p} < 0,001). Dies zeigt, dass die Methode mit unintensivierten Sätzen als Kontrollstufe in dem Sinne funktioniert hat, dass sie im Vergleich zu mittels deskriptivem oder expressivem Intensivierer modifizierten Sätze von den Versuchspersonen auf der Skala grundsätzlich niedriger positioniert wurden. Darüber hinaus findet sich ein signifikanter Positionseffekt {($\chi$$^{2}$\,(1) = 6,96}, {\textit{p} < 0,01)}: Je fortgeschrittener die Position eines Items im Experiment ist, desto niedriger wird der Grad an Skalarität evaluiert.


Kurz gefasst belegen die Ergebnisse von Aufgabe I den in der Literatur angenommenen Unterschied zwischen der deskriptiven und expressiven Intensiviererklasse (vgl. \citealt[22]{Paradis.1997}; {\citealt[133]{Gutzmann.2019a}}). So wurde der Skalaritätsgrad von Sätzen mit expressivem Intensivierer von den Versuchspersonen in der Tat als signifikant höher empfunden als der von Sätzen, die im Experiment mit dem deskriptiven oder ohne Intensivierer präsentiert wurden. Dass sich der Unterschied zwischen den drei Stufen derart deutlich zeigt, lässt außerdem den Schluss zu, dass das Skalenrating als Untersuchungsmethode funktioniert hat.


\subsubsection{Vergleich der Intensivierer} \label{abs:I-Vergleich339}


In Tabelle \ref{SkaWe2} werden der gemittelte Skalenwert und die Standardabweichung für die einzelnen Intensivierer bei Aufgabe I dargelegt. Dies bietet einen detaillierten Überblick über die Einschätzung der Versuchspersonen und ermöglicht, die im Experiment abgefragten Ausdrücke in Bezug auf den ihnen zugeschriebenen Skalaritätsgrad miteinander zu vergleichen. Die expressiven Intensivierer sind in ansteigendem Skalenwert angeordnet.


%Diese Werte sind inklusive Standardabweichungen in Tabelle \ref{SkaWe2} aufgeführt.





\begin{table}
%\setlength{\tabcolsep}{32.75pt}
\begin{tabularx}{.8\textwidth}{llYY}
\lsptoprule
%       &  \multicolumn{1}{c}{\textbf{Aufgabe I}} &  \\
%       &  \multicolumn{1}{c}{\textbf{Skalarität}} &  \\  [0.5em]
Stufe & Intensivierer & Skalenwert &\textit{SD}\\
\midrule
\rowcolor[gray]{.90}
deskriptiv & \textit{sehr} &  63,90 & 17,9\\
expressiv  & \textit{arg}  &  59,87 & 20,8\\
\rowcolor[gray]{.90}
           & \textit{voll} &  62,40 & 22,4\\
           & \textit{ober} &  62,84 & 22,9\\
\rowcolor[gray]{.90}
           & \textit{derb} &  63,36 & 20,3\\
           & \textit{sau}  &  65,14 & 22,7\\
\rowcolor[gray]{.90}
           & \textit{unheimlich} &  65,78 & 24,1\\
           & \textit{verdammt}   &  67,05 & 20,5\\
\rowcolor[gray]{.90}
           & \textit{fürchterlich} &  68,05 & 23,0\\
           & \textit{furchtbar}    &  68,99 & 22,4\\
\rowcolor[gray]{.90}
           & \textit{krass}       &  69,02 & 22,4\\
           & \textit{schrecklich} &  70,09 & 20,0\\
\rowcolor[gray]{.90}
           & \textit{hammer} &  70,70 & 22,9\\
           & \textit{scheiß} &  71,36 & 25,0\\
\rowcolor[gray]{.90}
           & \textit{wahnsinnig} &  71,37 & 23,2\\
           & \textit{super}      &  72,51 & 19,0\\
\rowcolor[gray]{.90}
           & \textit{mega}        &  74,66 & 22,4\\
\multicolumn{1}{l}{unmarkiert} & \textit{$\emptyset$} &  35,27 & 27,4\\
\lspbottomrule
\end{tabularx}
\caption[Gemittelte Skalenwerte der einzelnen Intensivierer bei Aufgabe I von Experiment 1]{Gemittelte Skalenwerte und Standardabweichungen der einzelnen Intensivierer bei Aufgabe I -- Skalarität (angeordnet nach ansteigendem Skalenwert der expressiven Intensivierer).}
\label{SkaWe2}
\end{table}


%
\noindent Der Übersicht in Tabelle \ref{SkaWe2} ist zu entnehmen, dass sich der deskriptive Intensivierer \textit{sehr} bei Aufgabe I kaum von den expressiven Intensivierern unterscheidet. So liegt der gemittelte Skalenwert von \textit{sehr} mit 63,90 Punkten partiell sogar über dem Wert mancher expressiver Intensivierer: {\textit{arg} (59,87 Punkte)}, {\textit{voll} (62,40 Punkte)}, {\textit{ober} (62,84 Punkte)} und \textit{derb} (63,36 Punkte). Wie im Hinblick auf {Tabelle \ref{SkaWe1}} beschrieben, wurde die unmarkierte Variante mit durchschnittlich 35,27 Punkten niedrig bewertet. Schaut man sich das Skalenspektrum der expressiven Intensivierer an, ist festzustellen, dass sich deren Werte alle in einem mittleren Skalenbereich von {59,87 Punkten (\textit{arg})} und 74,66 Punkten (\textit{mega}) bewegen. Somit beträgt die Differenz zwischen dem am niedrigsten und dem am höchsten eingestuften expressiven Intensivierer etwa 14,79 Punkte. Diese Tendenzen sind im {Anhang \ref{Rat1}} in {Abbildung \ref{Skalar7}} mittels Säulendiagramm grafisch dargestellt. Zudem findet sich im {Anhang \ref{Pseu1}} in {Tabelle \ref{Pseudos1}} eine Übersicht der gemittelten Skalenwerte jeder Kombination aus Pseudonomen und Pseudoadjektiv. Darin ist zu erkennen, dass die Pseudowörter über ähnliche Werte verfügen. So weist die Kombination aus \textit{Das Gedalst} und \textit{wübig} mit durchschnittlich 62,55 Punkten den niedrigsten Skalenwert auf, wohingegen die Kombination aus \textit{Das Zafche} und \textit{göcktig} mit 69,69 Punkten auf der Skala am höchsten beurteilt wurde ($\Delta$ = ca. 7,14 Punkte).\footnote{Das Symbol Delta ($\Delta$) steht hier und im Folgenden für die Differenz zwischen zwei Werten.} Zur besseren Anschaulichkeit sind die gemittelten Skalenwerte je Kombination aus Pseudonomen und Pseudoadjektiv im Anhang \ref{Happy0} in {Abbildung \ref{Happy3}} mithilfe eines Säulendiagramms visualisiert. Die nur geringe Differenz zwischen den einzelnen Pseudowort-Kombinationen deutet darauf hin, dass die Verwendung von Pseudowörtern einen Einfluss der semantischen Eigenschaften verhindert und die Untersuchungsmethode funktioniert hat. Wäre der Unterschied zwischen ihnen an dieser Stelle ausgeprägter, so ließe sich spekulieren, dass die zugeschriebenen Skalenwerte durch die Pseudowörter beeinflusst gewesen sein könnten. Dies scheint jedoch nicht der Fall zu sein.


\subsubsection{Regressionsanalyse zu Frequenz und Persistenz}


Im Folgenden wird dem möglichen Einfluss von Frequenz (vgl. H1-Freq) und Persistenz (vgl. H1-Temp) auf den Skalaritätsgrad nachgegangen. Diese Überprüfung richtet sich nach der Gruppeneinteilung (d.\,h. `Hochfrequent' vs. `Niedrigfrequent' für H1-Freq sowie `Alt' vs. `Neu' für {H1-Temp}), die sich auf Grundlage der Untersuchungsergebnisse der in Kapitel \ref{Korpus} präsentierten Korpusstudie vornehmen ließ. Da die Korpusdaten klar für die Annahme zweier unterschiedlicher Frequenz- und Persistenzgruppen sprachen, wurde bei diesen Regressionsanalysen mit kategorialen statt kontinuierlichen Variablen gerechnet.


Zunächst wird die Hypothese H1-Freq beschrieben. Nach dieser sind oft vorkommende Intensivierer weniger expressiv als seltene Intensivierer \citep[vgl.][2]{Hauschild.1899}. Die Häufigkeiten konnten im Rahmen der in Kapitel \ref{Korpus} vorgestellten Korpusstudie abgeleitet werden. Auf Basis der Korpusfrequenz des letzten Zeitabschnitts (d.\,h. dem Zeitraum von 1990 bis 2017) wurden die expressiven Intensivierer in eine niedrigfrequente (n = 9) und eine hochfrequente {(n = 7)} Gruppe eingeordnet (vgl. Tabelle \ref{Fre3.1}). Diese Gruppen wurden im vollen Regressionsmodell als Ausgangspunkt gewählt, um den potenziellen Einfluss von Frequenz und somit die Gültigkeit von H1-Freq zu überprüfen.


Im Hinblick auf H1-Temp wurden in Abhängigkeit des Erstvorkommens ebenfalls zwei Gruppen gebildet. Dies geschah vor dem Hintergrund der Einschätzung, dass lange gebrauchte Intensivierer weniger expressiv sind als neuere Intensivierer \citep[vgl.][59]{Tobler.1868}. Anhand der in der Korpusstudie abgeleiteten diachronen Entwicklungsverläufe konnten drei Intensivierer detektiert werden, die erst im Laufe des Untersuchungszeitraums im Korpus in Erscheinung getreten sind, nämlich \textit{hammer}, \textit{krass} und \textit{mega}. Diese wurden bei der Datenmodellierung von den restlichen Ausdrücken (n = 13) separiert und im vollen Modell als Gruppe der neuen Intensivierer berücksichtigt.


Wiederholt kam das oben beschriebene statistische Vorgehen zum Einsatz. Zunächst wurde ein volles Regressionsmodell berechnet, in dem mithilfe von Rückwärtsselektion und \textit{anova}-Modellvergleich sukzessive alle Effekte ausgeschlossen wurden, die den Modell-Fit nicht signifikant verbesserten. Das volle Modell enthielt \textit{Fixed Effects} für die kategorialen Prädiktoren \textsc{Frequenzgruppe}, \textsc{Persistenzgruppe}, \textsc{Position} sowie deren Zwei-Wege-Interaktionen. Der \textsc{Skalenwert} wurde erneut als metrische abhängige Variable modelliert; er steht in diesem Fall für den Skalaritätsgrad, den die Versuchspersonen der jeweiligen Stufe oder dem jeweiligen Intensivierer zugesprochen haben. Um mögliche Unterschiede zwischen hoch- und niedrigfrequenten Ausdrücken zu finden, wurde die binäre nominale Variable \textsc{Frequenzgruppe} dummy-gecoded. Dabei wurde der Stufe `Hoch' der {Wert $0$} und der Stufe `Niedrig' der {Wert $1$} zugewiesen. Daneben wurde auch die binäre nominale Variable \textsc{Persistenzgruppe} dummy-gecoded, wobei die Stufe `Alt' als $0$ und die Stufe `Neu' als $1$ codiert wurde. Diese Stufen wurden anschließend jeweils miteinander verglichen. Außerdem wurden im vollen Modell versuchspersonen- und itemweise \textit{Random Intercepts} und \textit{Random Slopes} für die Variablen \textsc{Skalenwert}, \textsc{Frequenzgruppe}, \textsc{Persistenzgruppe} und \textsc{Position} miteingerechnet, sodass auch potenzielle Unterschiede zwischen den Versuchspersonen und Pseudowörtern berücksichtigt wurden. Das finale Modell beinhaltet signifikante Haupteffekte für die Prädiktoren \textsc{Frequenzgruppe}, \textsc{Persistenzgruppe} und \textsc{Position}. Die berechneten Koeffizienten sind in {Tabelle \ref{LinRea0100}}.

 %
\begin{table}
%\setlength{\tabcolsep}{13.55pt}
\begin{tabularx}{.8\textwidth}{Xrrrrr}
\lsptoprule
Prädiktor & \textit{Est.} & \textit{SE} & \textit{$\chi$$^{2}$} & \textit{p} & \\
\midrule
(Intercept) & 70,96 & 2,29 & 30,94 & < 0,001 & $\star$$\star$$\star$\\
Alt vs. Neu & \textcolor{white}{-}5,91 & 1,36 & 16,58 & < 0,001 & $\star$$\star$$\star$\\
Hoch vs. Niedrig & -2,49 & 1,06 & \textcolor{white}{-}5,31 & < 0,05 & $\star$\\
Position & -0,31 & 0,13 & \textcolor{white}{-}5,95 & < 0,05 & $\star$\\
\lspbottomrule
\end{tabularx}
\caption[Finales Regressionsmodell zu Aufgabe I von Experiment 1 (Frequenz mit `Hoch' als Referenzlevel, Persistenz mit `Alt' als Referenzlevel)]{Finales Regressionsmodell zu Aufgabe I -- Skalarität (Frequenz mit `Hoch' als Referenzlevel, Persistenz mit `Alt' als Referenzlevel).}
\label{LinRea0100}
\end{table}

\noindent Das in Tabelle \ref{LinRea0100} dargelegte Regressionsmodell zeigt in Bezug auf Skalarität zum einen einen statistisch signifikanten Unterschied zwischen den beiden Persistenzgruppen an {($\chi$$^{2}$\,(1) = 16,58}, \textit{p} < 0,001). Demnach wird der Grad an Skalarität bei neuen Intensivierern auf der Skala insgesamt höher bewertet als bei älteren Intensivierern. Mit anderen Worten: Der eingeschätzte Skalaritätsgrad der neuen Intensivierer \textit{hammer}, \textit{krass} und \textit{mega} unterscheidet sich signifikant von dem der restlichen Intensivierer. Zum anderen unterscheiden sich auch die beiden Frequenzgruppen signifikant voneinander {($\chi$$^{2}$\,(1) = 5,31}, {\textit{p} < 0,05)}. Das bedeutet, dass die Versuchspersonen Sätze mit niedrigfrequenten Intensivierern mit einem niedrigeren Grad an Skalarität assoziiert haben als Sätze, die im Experiment mit hochfrequenten Intensivierern präsentiert wurden. Weiterhin ist ein signifikanter Positionseffekt auszumachen ({$\chi$$^{2}$\,(1) = 5,95}, \textit{p} < 0,05), der darauf hinweist, dass die eingeschätzte Skalarität mit fortschreitendem Experiment abnimmt. Dieser Befund legt einen Gewöhnungseffekt nahe, der sich bei den Versuchspersonen über die Dauer des Experiments hinweg eingestellt zu haben scheint.


Von den Ergebnissen der Regressionsanalyse ausgehend lässt sich zusammenfassen, dass die beiden Faktoren Frequenz und Persistenz anscheinend durchaus einen signifikanten Effekt auf die wahrgenommene Skalarität haben. So verbanden die Versuchspersonen die im Korpus rezent etablierten Intensivierer \textit{hammer}, \textit{krass} und \textit{mega} grundsätzlich mit einem höheren Skalaritätsgrad als die älteren Intensivierer. Darüber hinaus wird auch der Grad an Skalarität bei den hochfrequenten Intensivierern in Summe als deutlich stärker empfunden als bei den niedrigfrequenten Intensivierern. Dies lässt den Schluss zu, dass der Skalaritätsgrad mit zunehmender Frequenz ansteigt.




\subsection{Aufgabe II -- Sprechereinstellung} \label{abs:VP-Vergleich2423}


\subsubsection{Vergleich der Stufen} \label{abs:I-Vergleich426}


Tabelle \ref{SkaWex} bietet einen Überblick über die gemittelten Skalenwerte und Standardabweichungen für die  Stufen zu Aufgabe II. Wie zuvor sind die expressiven Intensivierer zunächst zu einer Stufe zusammengefasst.


\begin{table}
%\setlength{\tabcolsep}{16pt}
\begin{tabularx}{.8\textwidth}{lYY}
\lsptoprule
%   & \textbf{Aufgabe II} & \\
%  & \textbf{Sprechereinstellung} &\\
Stufe & Skalenwert &\textit{SD}\\
\midrule
deskriptiv & 42,34 & 23,3\\
expressiv  & 65,24 & 24,9\\
unmarkiert & 30,11 & 26,2\\
\lspbottomrule
\end{tabularx}
\caption[Gemittelte Skalenwerte der Stufen bei {Aufgabe II} von Experiment 1]{Gemittelte Skalenwerte und Standardabweichungen der Stufen bei Aufgabe II -- Sprechereinstellung (in alphabetischer Reihenfolge).}
\label{SkaWex}
\end{table}


\noindent Den Werten aus Tabelle \ref{SkaWex} nach wurden Sätze mit expressivem Intensivierer von den Versuchspersonen auf der Skala erheblich höher gerankt als Sätze mit deskriptivem Intensivierer oder die unmarkierte Variante. So wurden der expressiven Stufe im Schnitt {65,24 Punkte} zugeschrieben, wohingegen die deskriptive Stufe auf der Skala durchschnittlich bei 42,34 Punkten und somit um etwa 22,90 Punkte niedriger positioniert wurde. Schlusslicht hinsichtlich Sprechereinstellung bildet erwartungsgemäß die unmarkierte Variante, die eine Differenz von ca. 35,13 Punkten zur expressiven Stufe und {ca. 12,23} Punkten zur deskriptiven Stufe anzeigt. Die Stufen unterscheiden sich bezüglich der in Aufgabe II erfragten Sprechereinstellung also wesentlich stärker voneinander als bei der in Aufgabe I untersuchten Skalarität.


%Nichtsdestotrotz ist bei der unmarkierten Stufe auffällig, dass ihr sowohl im Kontext von Skalarität als auch von Sprechereinstellung von den Versuchspersonen eine vergleichsweise hohe Skalenposition zugeschrieben wurde. So würde man aufgrund des fehlenden Intensivierers vielmehr annehmen, dass unintensivierte Sätze weder als übermäßig skalar noch als stark expressiv wahrgenommen und diese in der Folge intuitiv im untersten Skalenbereich verortet werden. Die möglichen Gründe für diese unerwartet hohe Skaleneinschätzung werden in Abschnitt \ref{abs:Diskussion22} vertiefend betrachtet. 
Diese Tendenzen sind in Abbildung \ref{ExpressivBox} mittels Boxplots dargestellt. Auch darin ist  die soeben beschriebene Rangfolge zu sehen, nach der die expressive Stufe in Bezug auf Sprechereinstellung von den Versuchspersonen auf der Skala am höchsten verortet wurde. Darüber hinaus ist zu erkennen, dass der Abstand der expressiven Stufe zur deskriptiven und unmarkierten Stufe bei dieser Aufgabe merklich größer ist, als es in den in Abbildung \ref{SkalarBox} präsentierten Boxplots zu Aufgabe I der Fall ist.

%Nichtsdestotrotz wurde letztere, wie soeben beschrieben, auf der Skala deutlich höher verortet als infolge der fehlenden Skalarität und Expressivität zu erwarten war, waren in diesen Sätzen doch weder deskriptive noch expressive Intensivierer enthalten.
 
%Dabei ist auch hier das oben beschriebene Ranking zu erkennen, nach dem Sätzen mit expressiven Intensivierern von den Versuchspersonen der höchste Grad an Sprechereinstellung zugeschrieben wurde, wohingegen die unmarkierte Variante die niedrigste Skalenposition aufweist. 
%Diese Tendenzen sind in Abbildung \ref{ExpressivBox} in Form von Boxplots grafisch dargestellt. 


\begin{figure}
  \includegraphics[width=1\textwidth]{BoxExpr_2}
	\caption[Gemittelte Skalenwerte der Stufen bei {Aufgabe II} von Experiment 1]{Gemittelte Skalenwerte der Stufen bei Aufgabe II -- Sprechereinstellung (x-Achse = Stufe, y-Achse = Skalenwert).}
  \label{ExpressivBox}
\end{figure}


\subsubsection{Regressionsanalysen zum Vergleich der Stufen}


Nachfolgend werden die Regressionsmodelle für Aufgabe II vorgestellt. Auch diese beruhen auf dem jeweiligen vollen Modell, das durch Rückwärtsselektion und \textit{anova}-Funktion nach und nach solange vereinfacht wurde, bis bloß noch signifikante Effekte darin enthalten waren. Das volle Regressionsmodell beinhaltete \textit{Fixed Effects} für den \textsc{Skalenwert} als metrische abhängige Variable sowie für die kategorialen Prädiktoren \textsc{Stufe} und \textsc{Position} sowie deren Zwei-Wege-Interaktionen. In dieser Aufgabe steht die Variable \textsc{Skalenwert} für den Grad an Sprechereinstellung, den die Versuchspersonen der jeweiligen Stufe bzw. dem jeweiligen Intensivierer zugesprochen haben. Um einen möglichen graduellen Effekt und Unterschiede zwischen den drei Stufen zu finden, wurde das Referenzlevel `Expressiv' jeweils mit {`Deskriptiv'} und `Unmarkiert' kontrastiert. Wie schon bei Skalarität geschah dies in Form von Dummy-Codierung, bei der die Ausprägungsformen der mehrkategorialen Variable \textsc{Stufe} in zwei Schritten als $0$ und $1$ kodiert und dadurch nach binären Kontrasten separiert wurden {\citep[vgl.][407,\,442]{Sauer.2019}}. Auf diese Weise wurden auch hier aus einer unabhängigen Variable mit drei Stufen zwei unabhängige Variablen mit je zwei Stufen. Außerdem fanden sich im vollen Modell versuchspersonen- und itemweise \textit{Random Intercepts} und \textit{Random Slopes} für die Variablen \textsc{Skalenwert}, \textsc{Stufe} und  \textsc{Position}. Das finale Regressionsmodell enthält signifikante Haupteffekte für den Prädiktor \textsc{Stufe}. Die berechneten Koeffizienten sind in {Tabelle \ref{LinReg0002}} zusammengetragen.

%
\begin{table}
%\setlength{\tabcolsep}{9.5pt}
\begin{tabularx}{\textwidth}{Xccccr}
\lsptoprule
Prädiktor & \textit{Est.} & \textit{SE} & \textit{$\chi$$^{2}$} & \textit{p} & \\
\midrule
(Intercept) & \textcolor{white}{-}65,38 & 2,26 & 126,18 & < 0,001 & $\star$$\star$$\star$\\
Expressiv vs. Deskriptiv & -22,75 & 3,26 & \textcolor{white}{-}36,86 & < 0,001 & $\star$$\star$$\star$ \\
Expressiv vs. Unmarkiert & -34,47 & 3,83 & \textcolor{white}{-}51,58 & < 0,001 & $\star$$\star$$\star$\\
\lspbottomrule
\end{tabularx}
\caption[Finales Regressionsmodell zu Aufgabe II von Experiment 1 (mit `Expressiv' als Referenzlevel)]{Finales Regressionsmodell zu Aufgabe II -- Sprechereinstellung (mit `Expressiv' als Referenzlevel).}
\label{LinReg0002}
\end{table}

\noindent Die in Tabelle \ref{LinReg0002} gezeigten Koeffizienten deuten darauf hin, dass sich die expressive Stufe bei Aufgabe II signifikant von den beiden anderen Stufen unterscheidet. So ist einerseits ein statistisch signifikanter Unterschied beim Vergleich der expressiven Stufe mit der deskriptiven Stufe zu verzeichnen ({$\chi$$^{2}$\,(1) = 36,86}, \textit{p} < 0,001). Andererseits unterscheiden sich die expressive und die unmarkierte Stufe ebenso signifikant voneinander ($\chi$$^{2}$\,(1) = 51,58, {\textit{p} < 0,001}). Dieser Befund stimmt im Wesentlichen mit der Annahme überein, dass Sätzen mit expressivem Intensivierer grundsätzlich ein höherer Grad an Sprechereinstellung zugeschrieben wird als Sätzen mit deskriptivem Intensivierer oder unintensivierten Sätzen.


Auch bei Aufgabe II wurde das Referenzlevel der mehrkategorialen Variable \textsc{Stufe} in einer weiteren Analyse auf `Unmarkiert' gesetzt, um zu überprüfen, inwiefern sich die unintensivierten Sätze in Hinsicht auf die ausgedrückte Sprechereinstellung von Sätzen mit deskriptivem oder expressivem Intensivierer unterscheiden. Erneut wurden ausgehend vom vollen Modell durch Rückwärtsselektion und Modellvergleich alle Effekte getilgt, die den Modell-Fit nicht signifikant verbesserten. Die \textit{Random-Effects}-Struktur entsprach mit Ausnahme des Referenzlevels der beiden binären Vergleiche der des obigen Modells. Das finale Modell beinhaltet signifikante Haupteffekte für den Prädiktor \textsc{Stufe}. Die berechneten Koeffizienten sind in {Tabelle \ref{LinReg0020} angegeben}.


\begin{table}
%\setlength{\tabcolsep}{10pt}
\begin{tabularx}{\textwidth}{Xccccr}
\lsptoprule
Prädiktor & \textit{Est.} & \textit{SE} & \textit{$\chi$$^{2}$} & \textit{p} & \\
\midrule
(Intercept) & 31,02 & 3,05 & 94,75 & < 0,001 & $\star$$\star$$\star$\\
Unmarkiert vs. Expressiv & 34,39 & 3,81 & 51,58 & < 0,001 & $\star$$\star$$\star$ \\
Unmarkiert vs. Deskriptiv & 11,63 & 3,01 & 13,08 & < 0,001 & $\star$$\star$$\star$\\
\lspbottomrule
\end{tabularx}
\caption[Finales Regressionsmodell zu Aufgabe II von Experiment 1 (mit `Unmarkiert' als Referenzlevel)]{Finales Regressionsmodell zu Aufgabe II -- Sprechereinstellung (mit `Unmarkiert' als Referenzlevel).}
\label{LinReg0020}
\end{table}


\noindent Die Werte in Tabelle \ref{LinReg0020} verdeutlichen, dass unintensivierte Sätze mit einem niedrigeren Grad an Sprechereinstellung assoziiert werden als mittels Intensivierer modifizierte Sätze. So zeigt sich beim Vergleich der unmarkierten Stufe mit der expressiven Stufe ein signifikanter Unterschied ($\chi$$^{2}$\,(1) = 51,58, \textit{p} < 0,001), ebenso wie der beobachtete Unterschied zwischen der unmarkierten Stufe und der deskriptiven Stufe signifikant ist ($\chi$$^{2}$\,(1) = 13,08, {\textit{p} < 0,001}).


Zusammengenommen bestätigen also auch die Ergebnisse der Datenauswertung von Aufgabe II den in der theoretischen Literatur postulierten Unterschied zwischen der deskriptiven und expressiven Intensiviererklasse (vgl. {\citealt[315]{Lang.1983}}; \citealt[57]{Ortner.2014}). So wurde der Grad an Sprechereinstellung bei Sätzen mit expressivem Intensivierer von den Versuchspersonen tatsächlich als signifikant höher eingeschätzt als der von Sätzen, die mit dem deskriptiven Intensivierer oder unintensiviert präsentiert wurden.





\subsubsection{Vergleich der Intensivierer} \label{abs:I-Vergleich513}


Um einen besseren Überblick über die Einschätzung der einzelnen Ausdrücke zu erhalten und Vergleichbarkeit zu gewährleisten, werden die gemittelten Skalenwerte und Standardabweichungen im Folgenden für Aufgabe II im Detail betrachtet. Diese Werte sind {Tabelle \ref{SkaWey}} zu entnehmen. Auch hier sind die expressiven Intensivierer in ansteigendem Skalenwert angeordnet.




\begin{table}
%\setlength{\tabcolsep}{26.25pt}
\begin{tabularx}{.8\textwidth}{Xlrr}
\lsptoprule
%       &  \multicolumn{1}{c}{\textbf{Aufgabe II}}          &  \\
%       &  \multicolumn{1}{c}{\textbf{Sprechereinstellung}} &  \\
Stufe & Intensivierer  & Skalenwert                       & \textit{SD}\\
\midrule
 \rowcolor[gray]{.90}
\multicolumn{1}{l}{deskriptiv} & \textit{sehr} &  42,34 &  23,3\\
 \multicolumn{1}{l}{expressiv} & \textit{arg}  &  56,94 &  24,6\\
  \rowcolor[gray]{.90}
 & \textit{voll}         & 59,61 &   23,7\\
 & \textit{wahnsinnig}   & 60,33 &   26,4\\
    \rowcolor[gray]{.90}
  & \textit{ober}        & 60,43 &   25,2\\
 & \textit{unheimlich}   & 60,94 &   24,9\\
  \rowcolor[gray]{.90}
 & \textit{derb}         & 62,63 &   22,5\\
   & \textit{hammer}     & 66,06 &   27,5\\
  \rowcolor[gray]{.90}
 & \textit{mega}         & 66,58 &   24,9\\
  & \textit{super}       & 66,60 &   23,5\\
  \rowcolor[gray]{.90}
  & \textit{furchtbar}   & 66,64 &   25,3\\
  & \textit{verdammt}    & 66,66 &   22,0\\
  \rowcolor[gray]{.90}
  & \textit{schrecklich} & 66,76 &   23,6\\
 & \textit{sau}          & 68,24 &   25,1\\
  \rowcolor[gray]{.90}
   & \textit{krass}      & 68,63 &   26,9\\
  & \textit{fürchterlich} & 68,75 &  22,3\\
  \rowcolor[gray]{.90}
 & \textit{scheiß}       & 78,12 &   25,0\\
\multicolumn{1}{l}{unmarkiert} & \textit{$\emptyset$} & 30,11 &  26,1\\
\lspbottomrule
\end{tabularx}
\caption[Gemittelte Skalenwerte der einzelnen Intensivierer bei Aufgabe II von Experiment 1]{Gemittelte Skalenwerte und Standardabweichungen der einzelnen Intensivierer bei Aufgabe II -- Sprechereinstellung (angeordnet nach ansteigendem Skalenwert der expressiven Intensivierer).}
\label{SkaWey}
\end{table}



Auch die in Tabelle \ref{SkaWey} präsentierten Werte deuten darauf hin, dass sich der Unterschied zwischen den Stufen stark vergrößert hat. Hier liegt der gemittelte Skalenwert des deskriptiven Intensivierers \textit{sehr} mit 42,34 Punkten weit unter den Werten der expressiven Intensivierer. Abgeschlagen dahinter befindet sich die unmarkierte Variante, der von den Versuchspersonen ein Skalenwert von durchschnittlich 30,11 Punkten zugewiesen wurde. Wieder ist das Skalenspektrum der expressiven Intensivierer recht eng gefasst: So liegt der niedrigste expressive Intensivierer \textit{arg} auf der Skala im Schnitt bei 56,94 Punkten, während \textit{scheiß} mit durchschnittlich 78,12 Punkten am höchsten positioniert wurde {($\Delta$ = ca. 21,20 Punkte)}. Auffällig ist, dass \textit{arg} neben Skalarität auch im Hinblick auf Sprechereinstellung als am niedrigsten bewertet wurde. Diese Tendenzen sind im Anhang \ref{Rat2} in Abbildung \ref{Expressiv7} in Form eines Säulendiagramms visualisiert. Im Anhang \ref{Pseu2} findet sich in {Tabelle \ref{Pseudos2}} ferner eine Übersicht der gemittelten Skalenwerte jeder Kombination aus Pseudonomen und Pseudoadjektiv. Darin ist zu verzeichnen, dass sich die Pseudowörter hinsichtlich ihrer Skaleneinschätzung ähnlich verhalten. Wurde der Kombination aus \textit{Der Cobe} und \textit{speiglich} mit durchschnittlich 58,75 Punkten von den Versuchspersonen der niedrigste Skalenwert zugeschrieben, weist die Kombination aus \textit{Die Itne} und \textit{schühnlich} mit durchschnittlich 65,34 Punkten den höchsten Skalenwert auf ({$\Delta$ = ca. 6,59 Punkte}). Auch hier legt die bloß geringe Differenz nahe, dass die Verwendung von Pseudowörtern einem Einfluss der semantischen Eigenschaften entgegengewirkt hat. Erneut sind die Skalenwerte der Pseudowörter zur besseren Anschaulichkeit im Anhang \ref{Happy4} in Abbildung \ref{Happy5} mittels Säulendiagramm {grafisch dargestellt}.





\subsubsection[Einfluss von Sprechereinstellung auf Skalarität]{Beeinflusst die Sprechereinstellung die Skalarität?}


Im Anschluss wird untersucht, ob zwischen der in Aufgabe I erfragten Skalarität und der in Aufgabe II erfragten Sprechereinstellung eine signifikante Korrelation besteht. In diesem Kontext ist gemäß der Definition von Expressivität nach \citet[133]{Gutzmann.2019a} vorstellbar, dass ein höherer Grad an Sprechereinstellung immer von einem höheren Grad an Skalarität begleitet wird. Wenn der Skalenwert der Variable \textsc{Sprechereinstellung} zunimmt, steigt in der Folge mutmaßlich auch der Skalenwert der Variable \textsc{Skalarität} an. Zur Quantifizierung des potenziellen linearen Merkmalszusammenhangs  wurde auf Basis der Mittelwerte der expressiven Intensivierer der Korrelationskoeffizient \textit{r} nach Pearson berechnet,\footnote{Die Produkt-Moment-Korrelation nach Pearson ist eine parametrische Methode, die als Zusammenhangsmaß den Korrelationskoeffizienten \textit{r} nimmt. Dabei handelt es sich um die Wurzel aus dem Bestimmtheitsmaß, die zwischen -$1$ ($\hateq$ negative Korrelation) und +$1$ {($\hateq$ positive Korrelation)} liegt. Je mehr sich der Wert $0$ annähert, desto geringer ist der Zusammenhang zwischen zwei Variablen. In diesem Fall spricht man von Unkorreliertheit (vgl. {\citealt[128ff.]{Mittag.2016}}; \citealt[89]{Winter.2020}).} dessen \textit{Estimate} wie folgt aussieht: {|\textit{r}| = 0,60}. Das Ergebnis des zugehörigen \textit{t}-Tests\footnote{Mithilfe des Hypothesentests \textit{t}-Test kann festgestellt werden, ob sich die Mittelwerte zweier abhängiger Stichproben signifikant voneinander unterscheiden \citep[vgl.][280f.]{Sauer.2019}.} ist signifikant ({\textit{t}\,(14) = 2,78}, \textit{p} < 0,05), was bedeutet, dass Skalarität und Sprechereinstellung im Experiment in der Tat in einer Korrelation stehen: Wenn der Grad an Sprechereinstellung von den Versuchspersonen als hoch empfunden wurde, wurde ebenso der Skalaritätsgrad auf der Skala höher verortet.


Ergänzend zur Korrelationsanalyse wurde ein lineares Regressionsmodell gerechnet, mit dessen Hilfe überprüft werden kann, ob der Grad an Sprechereinstellung im Experiment den Grad an Skalarität bedingte. Auch hier ist denkbar, dass die Skalarität mit steigender Sprechereinstellung zunimmt. Es wurde also getestet, ob die unabhängige Variable \textsc{Sprechereinstellung} einen signifikanten Effekt auf die abhängige Variable \textsc{Skalarität} hat. Das lineare Modell beinhaltet \textit{Fixed Effects} für die \textsc{Skalarität} und die \textsc{Sprechereinstellung}. Die berechneten Koeffizienten sind in {Tabelle \ref{LM01}} aufgeführt.


\begin{table}
%\setlength{\tabcolsep}{6.5pt}
\begin{tabularx}{\textwidth}{Xccccr}
\lsptoprule
Prädiktor & \textit{Est.} & \textit{SE} & \textit{F} & \textit{p} & \\
\midrule
(Intercept) & 24,52 & \textcolor{white}{-}7,03 & 12,15 & < 0,001 & $\star$$\star$$\star$\\
Skalarität vs. Sprechereinstellung & 66,38 & 11,18 & 35,22 & < 0,001 & $\star$$\star$$\star$\\
\lspbottomrule
\end{tabularx}
\caption[Lineares Regressionsmodell zum Einfluss von Sprechereinstellung auf Skalarität bei Experiment 1]{Lineares Regressionsmodell zum Einfluss von Sprechereinstellung auf Skalarität.}
\label{LM01}
\end{table}

\noindent Den Werten aus Tabelle \ref{LM01} nach hat die Sprechereinstellung tatsächlich einen signifikanten Einfluss auf die Skalarität ({$F$ = 35,22}, \textit{p} < 0,001). Demzufolge liegt es nahe, dass der Grad an Skalarität umso größer wird, desto höher der Grad an Sprechereinstellung ist. Dieser Befund stimmt mit der Annahme im Sinne von \citet[133]{Gutzmann.2019a} überein. Darüber, ob die Annahme, dass sich die Sprechereinstellung auf die Skalarität eines Ausdrucks auswirkt, grundsätzlich korrekt ist, macht das Modell jedoch keine Aussage.


Dieses Ergebnis ist in Abbildung \ref{Korr00} in Form einer linearen Regression visualisiert. Die Regressionslinie illustriert die vorhergesagte, geschätzte Sprechereinstellung der im Experiment berücksichtigten expressiven Intensivierer. Dabei zeigt sie den systematischen Zusammenhang zwischen Skalarität und Sprechereinstellung an, der sich hier durch eine Steigung ausdrückt. Die Tatsache, dass der Grad an Skalarität mit steigendem Grad an Sprechereinstellung zunimmt, belegt ebendiesen korrelativen Zusammenhang. 


\begin{figure}
\includegraphics[width=\textwidth]{SkalExpr1}
\caption[Lineare Regression zum Einfluss von Sprechereinstellung auf Skalarität bei Experiment 1]{Lineare Regression zum Einfluss von Sprechereinstellung auf Skalarität (x-Achse = Sprechereinstellung, y-Achse = Skalarität).}
\label{Korr00}
\end{figure}






\subsubsection{Regressionsanalyse zu Frequenz und Persistenz}


Als Nächstes wird die in Aufgabe II erfragte Sprechereinstellung im Hinblick auf die oben erläuterten Frequenz- und Persistenzgruppen untersucht. Um dem möglichen Einfluss von Frequenz und Persistenz und somit der Gültigkeit der Hypothesen H1-Freq und H1-Temp nach {\citet[2]{Hauschild.1899}} und \citet[59]{Tobler.1868} auf den Grund zu gehen, wurden die durch die Korpusanalyse eingeteilten Gruppen im vollen Modell als Ausgangspunkt gewählt. Dieses wurde durch Rückwärtsselektion und Modellvergleich schrittweise vereinfacht, bis ausschließlich signifikante Effekte darin enthalten waren. Das volle Regressionsmodell beinhaltete \textit{Fixed Effects} für den \textsc{Skalenwert} als metrische abhängige Variable sowie für die kategorialen Prädiktoren \textsc{Frequenzgruppe}, \textsc{Persistenzgruppe} und \textsc{Position} sowie deren Zwei-Wege-Interaktionen. Um potenzielle Unterschiede zwischen hoch- und niedrigfrequenten Ausdrücken zu finden, wurde die binäre nominale Variable \textsc{Frequenzgruppe} dummy-gecoded, wobei der Stufe `Hoch' der {Wert $0$} und der Stufe `Niedrig' der Wert $1$ zugewiesen wurde. Auch wurde die binäre nominale Variable \textsc{Persistenzgruppe} dummy-gecoded. Im Zuge dessen wurde die Stufe `Alt' als $0$ und die Stufe `Neu' als $1$ codiert. Diese Stufen wurden anschließend jeweils miteinander verglichen. Außerdem wurden im vollen Modell versuchspersonenweise \textit{Random Intercepts} und \textit{Random Slopes} für die Variablen  \textsc{Skalenwert}, \textsc{Frequenzgruppe}, \textsc{Persistenzgruppe} und \textsc{Position} miteingerechnet und dadurch eventuelle Unterschiede zwischen den Versuchspersonen berücksichtigt.\footnote{Da das volle Regressionsmodell unter Einbezug der Pseudowörter nicht konvergierte, konnten bei dieser Berechnung lediglich versuchspersonenweise \textit{Random Intercepts} und \textit{Random Slopes} miteingebunden werden.} Das finale Regressionsmodell enthält einen signifikanten Haupteffekt für den Prädiktor \textsc{Frequenzgruppe}. Die berechneten Koeffizienten der Regressionsanalyse befinden sich in {Tabelle \ref{LinRegzaz0}}.

%
\begin{table}
%\setlength{\tabcolsep}{14pt}
\begin{tabularx}{.8\textwidth}{Xccccr}
\lsptoprule
Prädiktor & \textit{Est.} & \textit{SE} & \textit{$\chi$$^{2}$} & \textit{p} & \\
\midrule
(Intercept) & 63,55 & 2,34 & 2,17 & < 0,001 & $\star$$\star$$\star$\\
Hoch vs. Niedrig & \textcolor{white}{-}3,01 & 1,06 & 7,97 & < 0,01 & \multicolumn{1}{r}{$\star$$\star$} \\
\lspbottomrule
\end{tabularx}
\caption[Finales Regressionsmodell zu Aufgabe II von Experiment 1 (Frequenz mit `Hoch' als Referenzlevel)]{Finales Regressionsmodell zu Aufgabe II -- Sprechereinstellung (Frequenz mit `Hoch' als Referenzlevel).}
\label{LinRegzaz0}
\end{table}

\noindent Die in Tabelle \ref{LinRegzaz0} gesammelten Koeffizienten liefern Bestätigung für die Hypothese, dass eine hohe Frequenz zu einem als niedriger eingeschätzten Grad an Sprechereinstellung führt, unterscheiden sich die beiden Frequenzgruppen dahingehend doch signifikant voneinander ({$\chi$$^{2}$\,(1) = 7,97}, \textit{p} < 0,01). Dementsprechend werden die hochfrequenten Intensivierer grundsätzlich mit einem niedrigeren Grad an Sprechereinstellung assoziiert als die niedrigfrequenten Intensivierer. Mit anderen Worten: Bei einem Wechsel von `Hochfrequent' nach `Niedrigfrequent' nimmt der wahrgenommene Grad an Sprechereinstellung deutlich zu. Nachweislich keinen Einfluss auf die Einschätzung des Grades an Sprechereinstellung hat dagegen die Persistenz. Dieser Befund wird weiter unten in {Abschnitt \ref{abs:Diskussion22}} ausführlich besprochen.


Insgesamt untermauern die Ergebnisse der Regressionsanalyse die {Hypothese H1-Freq} nach \citet[2]{Hauschild.1899}. So wurden die im Korpus hochfrequenten Intensivierer von den Versuchspersonen mit einem bedeutend niedrigeren Grad an Sprechereinstellung assoziiert als die niedrigfrequenten Intensivierer. Dagegen findet sich keine empirische Evidenz, die die Gültigkeit der Hypothese H1-Temp nach \citet[59]{Tobler.1868} belegt. Aus welchem Grund dieses Resultat nur eingeschränkt aussagekräftig ist, wird ebenso in {Abschnitt \ref{abs:Diskussion22}} detailliert dargelegt.






\subsection{Aufgabe III -- Einstellungspolarität} \label{abs:Ev-Vergleich}
%Im nächsten Abschnitt werden die Resultate der in Aufgabe III ermittelten Einstellungspolarität präsentiert. Bei dieser Aufgabe ging es für die Versuchspersonen darum, mithilfe 100 virtueller Taler auf die evaluative Haltung zu setzen, die der Sprecher zum bezeichneten Sachverhalt hat. Dies diente der Beantwortung der Frage, ob expressive Intensivierer eine als positiv oder negativ wahrgenommene Interpretation der Sprechereinstellung nach sich ziehen. Vor diesem Hintergrund ist bspw. vorstellbar, dass Intensivausdrücke, die Mitglieder der \glqq{}schrecklichen\grqq{} Intensivierergruppe beinhalten, von den Versuchspersonen als verstärkt negativ beurteilt werden, wohingegen das semantisch entleerte \textit{sehr} möglicherweise zu einer als neutral wahrgenommenen Sprechereinstellung führt. Ob diese Mutmaßung durch die Ergebnisse der Untersuchung gestützt wird, wird im Anschluss erörtert.






\subsubsection{Vergleich der Stufen} \label{abs:VP-Vergleich636}


Im ersten Schritt werden die Mittelwerte der gesetzten Taler und Standardabweichungen für die drei Stufen mit Blick auf die Einstellungspolarität begutachtet. Diese Werte befinden sich in Tabelle \ref{Bedin}. Wiederholt werden die expressiven Intensivierer zunächst zu einer Stufe zusammengefasst. In der Tabelle ist jeweils diejenige Spalte farblich hervorgehoben, deren Stufe die Versuchspersonen die höchste Anzahl von Talern zugewiesen haben.

%
\begin{table}
%\setlength{\tabcolsep}{24.5pt}
\begin{tabularx}{\textwidth}{X rr rr rr}
\lsptoprule
%  & \multicolumn{3}{c}{\textbf{Aufgabe III}}  \\
% \multirow{2}{*}{Stufe}  & \multicolumn{3}{c}{\textbf{Einstellungspolarität}}  \\ [0.5em]
  & \multicolumn{2}{c}{ {positiv}} & \multicolumn{2}{c}{ {neutral}} & \multicolumn{2}{c}{ {negativ}} \\
  \cmidrule(lr){2-3}\cmidrule(lr){4-5}\cmidrule(lr){6-7}
   & Taler  &  SD& Taler  &  SD& Taler  &  SD\\
\midrule
{deskriptiv} & 30,60 & 24,5   & \cellcolor[gray]{.90}46,37 & \cellcolor[gray]{.90}29,3    &   23,03 & 23,5 \\
{expressiv} & 32,03 & 29,5   & 27,34 & 26,2   &   \cellcolor[gray]{.90} 40,63 & \cellcolor[gray]{.90}32,3 \\
{unmarkiert} & 22,34  & 27,4  & \cellcolor[gray]{.90}62,05 & \cellcolor[gray]{.90}33,7   &   15,61 & 20,6 \\
\lspbottomrule
\end{tabularx}
\caption[Mittel der gesetzten Taler der Stufen bei {Aufgabe III} von Experiment 1]{Mittel der gesetzten Taler und Standardabweichungen der Stufen bei Aufgabe III -- Einstellungspolarität (Talersumme je Stufe = 100).}
\label{Bedin}
\end{table}


%
\noindent Die in Tabelle \ref{Bedin} gelisteten Werte weisen auf den ersten Blick darauf hin, dass sich die expressive Stufe grundlegend von den beiden anderen abhebt: Während Sätze der unmarkierten und deskriptiven Stufe vornehmlich als neutral eingeschätzt wurden, wurden Sätze mit expressivem Intensivierer dagegen verstärkt als negativ evaluiert. Das starke Übergewicht zugunsten der Abweichung von `Neutral' bei der expressiven Stufe lässt sich auch dadurch erkennen, indem man den Wert von `Neutral' mit dem der Abweichung, also der Summe aus `Positiv' und `Negativ', kontrastiert. Hier ergibt sich hinsichtlich `Neutral' vs. `Davon abweichend' bei der deskriptiven Stufe ein Verhältnis von 46,37 Talern zu 53,63 Talern, das sich als relativ ausbalanciert bezeichnen lässt. Die unmarkierte Stufe zeigt dahingehend ein Verhältnis von 62,05 Talern zu 37,95 Talern. Dies deutet darauf hin, dass die Sprechereinstellung bei unintensivierten Sätzen tendenziell vermehrt als neutral wahrgenommen wird. Bei der expressiven Stufe beträgt das Verhältnis hingegen 27,34 Taler zu 72,66 Taler, was bedeutet, dass die Sprechereinstellung bei expressiven Intensivierern stärker als positiv oder negativ bewertet wurde. Die beobachtete negative Tendenz hängt unter Umständen mit der Semantik der Ursprungslexeme zusammen, verfügt eine Vielzahl der untersuchten expressiven Ausdrücke doch über einen transparenten negativen Ursprung und eine ebensolche Konnotation (bspw. \textit{furchtbar} oder \textit{scheiß}). Im Unterschied dazu ist der Anteil an Ausdrücken mit klar erkennbarer positiver Herkunft oder Konnotation deutlich geringer (bspw. \textit{mega} oder \textit{super}).




\subsubsection{Regressionsanalyse zum Vergleich der Stufen} \label{abs:VP-Vergleich669}


Dieser Abschnitt widmet sich der Regressionsanalyse zu Aufgabe III. 
% Diese Aufgabe erfragte mithilfe einer Konstant-Summen-Skala die Polarität der Sprechereinstellung. Vor diesem Hintergrund wurde im Zusammenhang mit der Hypothese H3 angenommen, dass sich die Semantik des Ursprungslexems auf die Einschätzung der Polarität einwirkt, wenn der Intensivierer mit einem semantisch leeren Pseudowort kombiniert wird. Insofern stellt sich die Frage, ob die wahrgenommene Skalarität und Sprechereinstellung der Ausdrücke Einfluss auf ebendiese Einschätzung nehmen.
%Vor diesem Hintergrund wurde angenommen, dass diejenigen Intensivierer, die mit einem hohen Grad an Skalarität und Sprechereinstellung verbunden sind, eine stärkere Abweichung von 'neutral' hervorrufen als diejenigen, die jeweils mit einem niedrigeren Grad einhergehen. 
Das oben beschriebene statistische Vorgehen fand hier ebenfalls Anwendung. Wiederholt wurde zunächst ein volles Modell berechnet, das per Rückwärtsselektion und \textit{anova}-Funktion Schritt für Schritt vereinfacht wurde, bis alle Effekte getilgt waren, die den Modell-Fit nicht signifikant verbesserten. Das volle Modell enthielt \textit{Fixed Effects} für die \textsc{Einstellungspolarität}, die \textsc{Skalarität}, die \textsc{Sprechereinstellung}, die \textsc{Position} sowie deren Zwei-Wege-Interaktionen. Als Referenzlevel der abhängigen Variable \textsc{Einstellungspolarität}, die die drei Stufen `Positiv', `Neutral' und `Negativ' umfasst, wurde `Neutral' gewählt, d.\,h., es wurde überprüft, inwiefern die \textit{Fixed Effects} zu einer Abweichung -- ganz gleich, ob positiv oder negativ -- von ebendiesem neutralen Punkt führten. Die Abweichung von `Neutral' bzw. die Stufen `Positiv' und `Negativ' wurden mittels Dummy-Codierung zu einer Stufe zusammengefasst. Dabei wurde ihnen der Wert $1$ und der Stufe `Neutral' der Wert $0$ zugewiesen. Dieses Coding-Schema erlaubt es im Modell, den Grad an Abweichung von der neutralen Interpretation zu quantifizieren. Neben den \textit{Fixed Effects}  wurden im vollen Modell versuchspersonen- und itemweise \textit{Random Intercepts} und \textit{Random Slopes} für die Variablen \textsc{Einstellungspolarität}, \textsc{Skalarität}, \textsc{Sprechereinstellung} und \textsc{Position} miteingerechnet. Das finale Modell beinhaltet Haupteffekte für die Prädiktoren \textsc{Sprechereinstellung}, \textsc{Skalarität} und \textsc{Position}. Die berechneten Koeffizienten {Tabelle \ref{LinReg0004}} zu entnehmen.

%
\begin{table}
%\setlength{\tabcolsep}{11.5pt}
%\begin{tabular}{p{0.425\textwidth} p{0.1\textwidth} p{0.1\textwidth} p{0.1\textwidth}p{0.1\textwidth} p{0.05\textwidth}}
\begin{tabularx}{.8\textwidth}{Xccccr}
\lsptoprule
Prädiktor & \textit{Est.} & \textit{SE} & \textit{$\chi$$^{2}$} & \textit{p} & \\
\midrule
(Intercept) & \textcolor{white}{-}66,17 & 3,51 & 111,50 & < 0,001 & $\star$$\star$$\star$\\
Sprechereinstellung & -41,19 & 5,27 & \textcolor{white}{-}39,67 & < 0,001 & $\star$$\star$$\star$ \\
Skalarität & -14,03 & 4,77 & \textcolor{white}{-}\textcolor{white}{-}9,25 & < 0,01 & $\star$$\star$\\ 
Position & \textcolor{white}{-}\textcolor{white}{-}-1,35 & 0,65 & \textcolor{white}{-}\textcolor{white}{-}4,11 & < 0,05 & $\star$\\
\lspbottomrule
\end{tabularx}
\caption[Finales Regressionsmodell zu Aufgabe III von Experiment 1 (mit `Neutral' als Referenzlevel)]{Finales Regressionsmodell zu Aufgabe III -- Einstellungspolarität (mit `Neutral' als Referenzlevel).}
\label{LinReg0004}
\end{table}

\noindent Den in Tabelle \ref{LinReg0004} aufgeführten Werten nach werden Intensivierer, die mit einem hohen Grad an Sprechereinstellung assoziiert werden, von den Versuchspersonen als signifikant positiver oder negativer wahrgenommen ($\chi$$^{2}$\,(1) = 39,67, \textit{p} < 0,001). Ähnliches trifft auf Intensivierer zu, denen eine hohe Skalarität zugeschrieben wird ($\chi$$^{2}$\,(1) = 9,25, \textit{p} < 0,01). Demzufolge führt ein hoher Grad an Sprechereinstellung oder Skalarität zu einer von `Neutral' abweichenden Einschätzung der Sprechereinstellung -- völlig egal, ob positiv oder negativ. Ein weiterer signifikanter Effekt geht mit der \textsc{Position} einher ($\chi$$^{2}$\,(1) = 4,11, {\textit{p} < 0,05}). Somit wird der Grad der Abweichung von `Neutral' umso stärker eingeschätzt, je fortgeschrittener die Position eines Items im Experiment ist.


Es lässt sich zusammenfassen, dass die Ergebnisse der Datenauswertung von Aufgabe III den in der Literatur vertretenen Unterschied zwischen der deskriptiven und expressiven Intensiviererklasse stützen (vgl. {\citealt[135]{Gutzmann.2019a}}). Je höher der Grad an Skalarität und Sprechereinstellung eines Intensivierers empfunden wurde, desto stärker tendierten die Versuchspersonen zu einer als evaluativ eingeschätzten {Sprechereinstellung}.




\subsubsection{Vergleich der Einstellungspolarität} \label{abs:VP-Vergleich700}



\noindent Tabelle \ref{ValenzInten} bildet den Mittelwert der gesetzten Taler und die Standardabweichung je Intensivierer ab und ermöglicht eine genauere Inspektion der den Ausdrücken zugewiesenen Einstellungspolarität. Die expressiven Intensivierer sind zur besseren Übersicht nach ansteigendem negativen Skalenwert angeordnet.




%
\begin{table}
%\setlength{\tabcolsep}{7pt}
\fittable{
\begin{tabularx}{\textwidth}{Xlrrrrrr}
\lsptoprule
% \textbf{Aufgabe III} &  &  &  & & & & \\
% \textbf{Einstellungspolarität} &  &  &  & & & & \\
\multirow{2}{*}{Stufe} & \multirow{2}{*}{Intensivierer} &  &  & & & & \\
 & & \multicolumn{2}{c}{ {positiv}} & \multicolumn{2}{c}{ {neutral}} & \multicolumn{2}{c}{ {negativ}} \\
  \cmidrule(lr){3-4}\cmidrule(lr){5-6}\cmidrule(lr){7-8}
 &  & Taler            & \textit{SD} &  Taler            & \textit{SD} &  Taler           &\textit{SD}\\
\midrule
 \rowcolor[gray]{.90}
deskriptiv & \textit{sehr}   &  30,60 & 24,5 & 46,37 &  29,3 & 23,03 &  23,5\\
expressiv  & \textit{hammer} &  53,87 & 35,4 & 25,34 &  28,1 & 20,79 &  24,1\\
  \rowcolor[gray]{.90}
           & \textit{super}  &  50,42 & 31,1 & 25,67 &  26,0 & 23,91 &  25,9\\
           & \textit{mega}   &  49,58 & 31,6 & 23,30 &  26,4 & 27,12 &  26,9\\
  \rowcolor[gray]{.90}
           & \textit{ober}   &  36,46 & 28,6 & 33,23 &  28,6 & 30,31 &  28,5\\
       & \textit{wahnsinnig} &  38,00 & 29,0 & 31,39 &  29,2 & 30,61 &  27,9\\
  \rowcolor[gray]{.90}
       & \textit{unheimlich} &  31,55 & 26,5 & 34,84 &  27,4 & 33,61 &  28,0\\
             & \textit{voll} &  33,18 & 26,3 & 32,07 &  26,1 & 34,75 &  25,6\\
  \rowcolor[gray]{.90}
            & \textit{krass} &  34,43 & 25,8 & 28,99 &  26,0 & 36,58 &  28,2\\
            & \textit{sau}   &  33,57 & 28,5 & 27,52 &  27,1 & 38,91 &  29,0\\
  \rowcolor[gray]{.90}
            & \textit{derb}  &  25,42 & 23,2 & 31,07 &  26,6 & 43,51 &  31,1\\
        & \textit{verdammt}  &  28,60 & 28,1 & 27,60 &  23,4 & 43,80 &  29,7\\
  \rowcolor[gray]{.90}
              & \textit{arg} &  23,19 & 21,9 & 31,03 &  26,3 & 45,78 &  28,7\\
      & \textit{schrecklich} &  22,40 & 28,8 & 21,93 &  22,9 & 55,67 &  34,0\\
  \rowcolor[gray]{.90}
        & \textit{furchtbar} &  19,89 & 24,7 & 23,24 &  21,5 & 56,87 &  34,6\\
     & \textit{fürchterlich} &  15,12 & 19,0 & 24,57 &  25,1 & 60,31 &  33,6\\
  \rowcolor[gray]{.90}
           & \textit{scheiß} &  16,76 & 26,4 & 15,67 &  23,5 & 67,57 &  35,3\\
\multicolumn{1}{l}{unmarkiert} & \textit{$\emptyset$} &  22,34 & 27,4 & 62,05 &  33,7 & 15,61 &  20,6\\
\lspbottomrule
\end{tabularx}
}
\caption[Mittel der gesetzten Taler der einzelnen Intensivierer bei Aufgabe III von Experiment 1]{Mittel der gesetzten Taler und Standardabweichungen der einzelnen Intensivierer bei Aufgabe III -- Einstellungspolarität (Talersumme je Intensivierer = 100; angeordnet nach ansteigendem negativen Skalenwert).}
\label{ValenzInten}
\end{table}


\noindent Die Übersicht in Tabelle \ref{ValenzInten} erhärtet die in Bezug auf Tabelle \ref{Bedin} dargelegte Vermutung: Während der opake deskriptive Intensivierer \textit{sehr} und unintensivierte Sätze von den Versuchspersonen überwiegend als neutral eingestuft wurden, ist bei Intensivierern mit klar negativem Ursprung eine ausgeprägte Tendenz für die negative Polarität festzustellen (bspw. \textit{fürchterlich}, \textit{schrecklich} und \textit{scheiß}). Dies zeigt sich in acht von 18 Fällen ($\hateq$ 44,44\,\%): \textit{arg}, \textit{derb}, \textit{furchtbar}, \textit{fürchterlich}, \textit{sau}, \textit{scheiß}, \textit{schrecklich} und \textit{verdammt}. In fünf Fällen ist dagegen eine Tendenz in Richtung `Positiv' zu erkennen ($\hateq$ 27,78\,\%): \textit{hammer}, \textit{mega}, \textit{ober}, \textit{super} und \textit{wahnsinnig}. Eine Tendenz in Richtung `Neutral' kristallisiert sich in zwei Fällen heraus ($\hateq$ 11,11\,\%): bei \textit{sehr} und der unmarkierten Variante. Daneben weisen drei Fälle ein vorwiegend ausgewogenes Verhältnis {auf {($\hateq$ 16,67\,\%)}}: \textit{krass}, \textit{unheimlich} und \textit{voll}. Dieses Ergebnis deutet darauf hin, dass die Interpretation der Sprechereinstellung bei expressiven Intensivierern -- anders als beim deskriptiven Intensivierer oder der unmarkierten Variante des Satzes -- vermehrt als positiv oder negativ beurteilt wurde. Zur  Visualisierung sind die Mittelwerte der gesetzten Taler für jeden Intensivierer in Abbildung \ref{Aufg3.1} in Form von Säulendiagrammen dargestellt. Zusätzlich ist im Anhang \ref{Pseu3} in Tabelle \ref{Pseudos3} eine Übersicht über die gemittelten Taler je Kombination aus Pseudonomen und Pseudoadjektiv beigefügt

\begin{figure}
  \includegraphics[width=0.975\textwidth]{Bewertung3_new}
	\caption[Säulendiagramme zu den Mitteln der gesetzten Taler je Intensivierer bei Aufgabe III von Experiment 1]{Säulendiagramme zu den Mitteln der gesetzten Taler je Intensivierer bei Aufgabe III -- Einstellungspolarität (x-Achse = Einstellungspolarität, y-Achse = Skalenwert).}
  \label{Aufg3.1}
\end{figure}


%Zusätzlich ist im Anhang \ref{Pseu3} in Tabelle \ref{Pseudos3} eine Übersicht über die gemittelten Taler je Kombination aus Pseudonomen und Pseudoadjektiv beigefügt. 
%Darin ist auszumachen, dass zwölf der insgesamt 18 Pseudowörter vermehrt in Richtung 'negativ' tendierten ($\hat{=}$ 66,67\,\%). Somit wurde die Mehrheit der Pseudowort-Kombinationen aufgrund des Intensivierers \textit{de facto} als von 'neutral' abweichend beurteilt. In Anbetracht der Menge an Intensivierern mit negativem Ursprung ist dieser Befund nicht sonderlich überraschend.



\section{Diskussion}\label{abs:Diskussion22}
Das dargelegte Experiment wurde einerseits mit dem Ziel durchgeführt, Rückschlüsse auf die noch ungeklärten Hypothesen aus der Korpusstudie ziehen zu können. Gemäß der Hypothese H1-Freq im Sinne von {\citet[2]{Hauschild.1899}} wirkt sich ein frequenter Gebrauch von Intensivierern negativ auf ihre Expressivität aus, die in der Konsequenz schwindet. Nach {\citet[59]{Tobler.1868}} spielt dabei jedoch auch der Faktor Zeit eine Rolle, erstreckt sich die Abnutzung von Intensivierern doch womöglich über eine gewisse Zeitspanne hinweg. Diese Annahme wurde in der Hypothese {H1-Temp} aufgegriffen. Andererseits wurden auch die {Hypothesen H3 und H4} auf ihre Richtigkeit hin überprüft, nach denen Expressivität gleichermaßen mit einem höheren Grad an Skalarität und Sprechereinstellung einhergeht, wie {\citet[22]{Paradis.1997}} und {\citet[133]{Gutzmann.2019a}} sowie {\citet[315]{Lang.1983}} und {\citet[57]{Ortner.2014}} anführen. Im Kontext der {Hypothese H5} wurde außerdem getestet, ob die Sprechereinstellung, die expressive Intensivierer mutmaßlich kommunizieren, evaluativ ist und vom Adressaten daher als positiv oder negativ interpretiert wird, wie {\citet[135]{Gutzmann.2019a}} ferner ergänzt.


Im Folgenden werden die im vorangegangenen Abschnitt präsentierten Untersuchungsergebnisse zusammengefasst und mit Blick auf die zentralen Hypothesen evaluiert. Dies geschieht in der Reihenfolge, in der die Merkmale Skalarität, Sprechereinstellung und Einstellungspolarität im Experiment berücksichtigt wurden. Darüber hinaus werden auch die Ergebnisse zur Rolle von Frequenz und Persistenz aufgearbeitet und besprochen.




\subsection{Aufgabe I -- Skalarität}
In Aufgabe I ging es für die Versuchspersonen darum, auf einer Skala von `normal' bis `maximal' anzugeben, wo sich der Grad der Eigenschaft des beschriebenen Objekts befindet. Die Fragestellung war im Wortlaut an den Itemsatz angepasst, z.\,B. \textit{Wie fongig ist der Nümig?} Die Idee dahinter bestand nach H3 darin, dass Sätze mit expressivem Intensivierer vermutlich einen höheren Grad an Skalarität ausdrücken als Sätze mit deskriptivem Intensivierer oder unintensivierte Sätze (vgl. {\citealt[22]{Paradis.1997}}; {\citealt[133]{Gutzmann.2019a})}. Um die durch die Untersuchung gewonnenen Daten auf ihre statistische Signifikanz hin zu überprüfen, wurden mithilfe von R lineare gemischte Modelle berechnet, die aus der unabhängigen Variable die abhängige Variable voraussagen.


Die Ergebnisse zu Aufgabe I weisen darauf hin, dass die Hypothese H3 im Kern zutreffend ist. So werden Sätze mit expressivem Intensivierer mit einem signifikant höheren Skalaritätsgrad assoziiert als Sätze mit deskriptivem Intensivierer oder unintensivierte Sätze (vgl. {Tabelle \ref{LinReg0001}}). Daneben unterscheidet sich auch die unmarkierte Variante des Satzes dahingehend signifikant von Sätzen, die mittels deskriptivem Intensivierer modifiziert werden (vgl. {Tabelle \ref{LinReg0010}}). Diese Befunde entsprechen in Kombination der eingangs angeführten Rangordnung nach {\citet[133]{Gutzmann.2019a}} und stehen infolgedessen im Einklang mit den im Vorfeld getroffenen Erwartungen: unmarkierte Variante $\prec$ deskriptive Intensivierer $\prec$ expressive Intensivierer. Dass der Skalaritätsgrad bei ersteren beiden Stufen von den Versuchspersonen nicht höher eingeschätzt wurde, legt außerdem nahe, dass das Skalenrating als Untersuchungsmethode wie vorgesehen funktioniert hat. Nichtsdestotrotz sind bei den Ergebnissen zwei eklatante Beobachtungen zu machen: Einerseits unterscheidet sich der Skalaritätsgrad des deskriptiven Intensivierers \textit{sehr} bei der Einzelbetrachtung nicht übermäßig stark von dem der expressiven Intensivierer. Stattdessen wird er partiell sogar als höher empfunden. Dies legt die Vermutung nahe, dass sein Grad an Skalarität infolge pragmatischer Abnutzung und damit einhergehender Frequenzzunahme sukzessive angestiegen ist, weswegen der Ausdruck heute bloß eine neutrale skalare Bedeutung trägt.\footnote{Wie bis hierhin mehrfach erläutert wurde, trifft auf die expressiven Intensivierer diesbezüglich eventuell eine gegenläufige Entwicklung zu, derzufolge eine ansteigende Frequenz zu einem Verlust der Expressivität und damit zu einem niedrigeren Grad an Sprechereinstellung -- und gegebenenfalls auch Skalarität -- führt.} Andererseits stellt sich in Bezug auf die unmarkierte Stufe die Frage, aus welchen Gründen der durch sie ausgedrückte Skalaritätsgrad nicht niedriger bewertet wurde, ist {\citet[133]{Gutzmann.2019a}} zufolge doch davon auszugehen, dass diese allenfalls einen geringen Grad einer Merkmalsausprägung wiedergibt. Die genauen Gründe für die Skaleneinschätzung lassen sich auf Basis der vorliegenden Daten allerdings nicht ohne Spekulation benennen, weshalb die Frage in diesem Punkt unbeantwortet bleiben muss.


An dieser Stelle sei erwähnt, dass eine Einschätzung des Skalaritätsgrades expressiver Intensivierer aus methodischer Sicht nicht ganz unproblematisch ist, kann in einer solchen experimentellen Erhebung doch weder eine trennscharfe noch eine isolierte Betrachtung der semantischen Ebenen erfolgen. So wirkt die Expressivität im Sinne von \citet[133]{Gutzmann.2019a} unweigerlich auf die Skalarität ein. Deswegen war es im Experiment trotz des Untersuchungsschwerpunkts Skalarität nicht möglich, die Merkmalsausprägung eines im Satz beschriebenen Gegenstands ohne jeglichen Einfluss des expressiven Gehalts des jeweiligen Intensivierers zu testen und den Expressivitätsgrad von der skalaren Grundbedeutung zu separieren. Dies lässt eine direkte Überprüfung der Hypothese H3 naturgemäß kaum zu. Insofern bleibt festzuhalten, dass die Grenzen einer solchen empirischen Validierung in diesem Punkt erreicht sind und das Untersuchungsdesign eine gute – und die bislang einzige – Approximation zur Ermittlung der Grade beider Bedeutungsdimensionen darstellt, auch wenn das tatsächliche Ausmaß der Auswirkung von Expressivität auf Grundlage der Daten unklar ist.


Die Ergebnisse zur in Aufgabe I untersuchten Skalarität bestätigen in Summe also die Hypothese H3 nach \citet[22]{Paradis.1997} und \citet[133]{Gutzmann.2019a}. Demnach wurde der Grad an Skalarität bei Sätzen mit expressivem Intensivierer von den Versuchspersonen de facto als bedeutend höher wahrgenommen als der von Sätzen mit deskriptivem Intensivierer oder Sätzen, die im Experiment unintensiviert präsentiert wurden. Dies liefert die erste empirische Evidenz für den in der theoretischen Literatur angenommenen Unterschied zwischen den Intensiviererklassen und zeigt außerdem, dass ein Skalenrating ein geeignetes Instrument zur Erhebung des Merkmals Skalarität ist.



%Aufgrund des Umstands, dass der linke Skalenrand mit 'normal' beschriftet war, würde man die unintensivierten Sätze intuitiv im untersten Skalenbereich verortet sehen. Da die genauen Gründe hierfür auf Basis der Daten nicht zweifelsfrei benannt werden können, lässt sich an dieser Stelle lediglich spekulieren, dass die Einschätzung auf der Skala durch eine falsche Wahrnehmung der Skalenbeschriftung beeinflusst gewesen sein könnte. So ist bspw. denkbar, dass die Versuchspersonen die Reichweite der Skala nicht als von 'normal' bis 'maximal', sondern vielmehr als von 'niedrig' bis 'hoch' interpretiert und die unintensivierten Sätze in der Folge mit einem höheren Skalaritätsgrad als erwartet evaluiert haben. Möglich ist aber auch, dass die Versuchspersonen Schwierigkeiten bei der Orientierung innerhalb der Skala hatten. Abgesehen von den Rändern war diese nämlich vollkommen leer, was dazu diente, Priming-Effekte zu umgehen. Die fehlenden Abstufungen können die Beurteilung der Items auf der Skala allerdings erschwert haben. Insofern erwächst die Frage, ob numerische Skalenbeschriftungen oder Binnendifferenzierungen (bspw. in Form von einzelnen Segmenten) die Orientierung während der Evaluation erleichtert und somit zu anderen Ergebnissen geführt hätten.


%Den höchsten Grad an Merkmalsausprägung drückte nach Einschätzung der Versuchspersonen der neue, im Korpus mit 11,43 Treffern pro Million Textwörter niedrigfrequente expressive Intensivierer \textit{mega} aus. Ein möglicher Grund hierfür bietet die noch anhaltende Aktualität sowie die dadurch bedingte hohe Salienz des Ausdrucks: Angesichts dessen, dass man \textit{mega} im Alltag überall und auf unzählige Arten hört (man denke z.\,B. an Dieter Bohlen), verbindet man damit eventuell einen höheren Skalaritätsgrad als dies bei weniger aktuellen Ausdrücken der Fall ist.\footnote{Dies würde auch im Einklang stehen mit Annahme, dass expressive Intensivierer neben einem erhöhten Grad an Sprechereinstellung ebenso mit einem höheren Skalaritätsgrad einhergehen (vgl. u.\,a. {\citealt[315]{Lang.1983}}; {\citealt[22]{Paradis.1997}}; \citealt[57]{Ortner.2014}; {\citealt[133]{Gutzmann.2019a}}).} Ein weiteres Indiz für diese Idee liefert der zweitplatzierte expressive Intensivierer \textit{super}. Auch hierbei handelt es sich um einen gerne gebrauchten Ausdruck, der insbesondere in der Werbung allgegenwärtig erscheint. So macht sich dies bspw. die deutsche Einzelhandelskette \textsc{EDEKA} mithilfe des von Friedrich Liechtenstein besungenen Werbespots \glqq{}Supergeil\grqq{} zunutze,\footnote{\citet[]{Cornelius.2014}: \textit{https://www.youtube.com/watch?v=v2GVEpGqTX0} [letzter Aufruf: 21.11.2021].} was zu einer schätzungsweise höheren Prominenz des Ausdrucks führt.\footnote{Von einem großen Einfluss der Werbesprache auf den Gebrauch von \textit{super} berichten auch \citet[193]{Wellmann.1978}.} Ähnlich wie bei \textit{mega} kann besagte Aktualität zu einem als höher empfundenen Skalaritätsgrad beigetragen haben, auch wenn es sich bei \textit{super} um einen mit 391,26 Treffern pro Million Textwörter im Korpus hochfrequenten und lange verwendeten Ausdruck handelt. Das Schlusslicht unter den expressiven Intensivierern bildete \textit{arg}. Gemäß der Ergebnisse aus der Korpusstudie handelt es sich dabei um einen hochfrequenten Intensivierer (185,50 Treffer pro Million Textwörter), der über den gesamten Untersuchungszeitraum stetig im Gebrauch war. Auch wurden die niedrigfrequenten, älteren Intensivierer \textit{derb}, \textit{ober} und \textit{voll} im Hinblick auf Skalarität vergleichbar niedrig bewertet. In dieser Gruppe war vor allem \textit{derb} auffällig, der im Korpus der einzige Ausdruck war, der über die Zeit kontinuierlich an Frequenz eingebüßt hat und einen hinreichenden Grund zur Annahme des durch \citet[2]{Hauschild.1899} postulierten Expressivitätsverlusts bietet (vgl. H1). Bei diesen Ausdrücken ist vorstellbar, dass der durch sie ausgedrückte Skalaritätsgrad infolge ihres persistenten Gebrauchs als verhältnismäßig niedrig wahrgenommen wurde.

\subsection{Aufgabe II -- Sprechereinstellung}
Daneben gaben die Versuchspersonen in Aufgabe II auf einer zweiten Skala von `gar nicht' bis `maximal' an, wie stark der Sprecher seine persönliche Haltung zur beschriebenen Sachlage, bspw. zur Fongigkeit des Nümigs, ausdrückt. Dies wurde vor dem Hintergrund der Hypothese H4 abgefragt, nach der Sätze mit expressivem Intensivierer neben Skalarität auch die subjektive Einstellung des Sprechers ausdrücken (vgl. \citealt[315]{Lang.1983}; {\citealt[57]{Ortner.2014}}). Erneut wurden mittels R mehrere lineare gemischte Modelle berechnet, die den Einfluss der unabhängigen Variable auf die abhängige Variable angeben.


Den Ergebnissen zu Aufgabe II nach ist die Hypothese H4 im Wesentlichen korrekt, wird der Grad an Sprechereinstellung bei Sätzen mit expressivem Intensivierer doch als signifikant höher eingeschätzt als bei Sätzen mit deskriptivem Intensivierer oder unintensivierten Sätzen (vgl. {Tabelle \ref{LinReg0002}}). Dass letztere beiden von den Versuchspersonen somit mit einem bedeutend niedrigeren Grad an Sprechereinstellung assoziiert wurden, belegt auch hier, dass das Skalenrating funktioniert hat. Bei der Untersuchung von Sprechereinstellung ist der Unterschied zwischen den drei Stufen deutlich stärker ausgeprägt als bei der in Aufgabe I ermittelten Skalarität -- es zeichnet sich also eine klare Trennung nach Intensiviererklasse bzw. Stufe ab. Trotzdem ist auch bei {Aufgabe II} zu bemerken, dass die ausgedrückte Sprechereinstellung bei der unmarkierten Stufe von den Versuchspersonen als recht hoch wahrgenommen wurde. Wie oben erwähnt, lässt sich dieser Befund jedoch empirisch nicht erklären. Auch wurde die deskriptive Stufe unerwartet hoch bewertet. Dies ist in Anbetracht der semantischen Analyse deskriptiver Intensivierer insofern überraschend, als darin kein expressiver Anteil vorgesehen ist, weswegen \textit{sehr} damit eigentlich als nur wenig einstellungsausdrückend evaluiert werden sollte. Da im Vorfeld der Untersuchung dazu allerdings keine Vortests durchgeführt wurden, die Aussagen über die Erwartungshaltung gegenüber der Einschätzung der Stufen erlaubten, wird von einer Spekulation der möglichen Gründe für die Skaleneinschätzung der Versuchspersonen jedoch abgesehen. In einem weiteren Analyseschritt wurde untersucht, ob sich die unmarkierte Variante des Satzes signifikant von mittels deskriptivem Intensivierer modifizierten Sätzen unterscheidet. Dies ist augenscheinlich der Fall (vgl. {Tabelle \ref{LinReg0020}}), was bedeutet, dass bei dieser Aufgabe ebenfalls die oben beschriebene Rangfolge der drei Stufen auszumachen ist.


Insgesamt bekräftigen die Ergebnisse der in Aufgabe II erfragten Sprechereinstellung die Hypothese H4 nach \citet[315]{Lang.1983} und {\citet[57]{Ortner.2014}}. So wurden Sätze mit expressivem Intensivierer von den Versuchspersonen tatsächlich mit einem wesentlich höheren Grad an Sprechereinstellung assoziiert als Sätze mit deskriptivem Intensivierer oder unintensivierte Sätze. Dieser Befund demonstriert, dass expressive Ausdrücke sowohl mit einem höheren Grad an Skalarität als auch mit einem höheren Grad an Sprechereinstellung verbunden werden und untermauert den in der Literatur vertretenen Unterschied zwischen den Intensiviererklassen. Weiterhin wird klar, dass die Einteilung der Ausdrücke in deskriptiv und expressiv {treffend war}.
 

%Im Gegensatz dazu war das im Korpus niedrigfrequente \textit{scheiß} der Ausdruck, der von den Versuchspersonen als am expressivsten eingestuft wurde (41,96 Treffer pro Million Textwörter). Ein Grund hierfür findet sich mutmaßlich in der dem Ausdruck zugrunde liegenden klar erkennbaren Semantik, durch die ihm eine negative Konnotation anhaftet. Diese hat gegebenenfalls dazu geführt, dass er als vergleichsweise ausdrucksstärker bewertet wurde und demnach besser geeignet zu sein scheint, um seine persönliche Einstellung offensiv zum Ausdruck zu bringen. Daher ist denkbar, dass die semantische Transparenz eines Intensivierers prinzipiell zu einem als höher wahrgenommenen Grad an Sprechereinstellung führt.



%So gab es keinen expressiven Intensivierer, dessen Grad an Sprechereinstellung als niedriger empfunden wurde als der des deskriptiven Intensivierers \textit{sehr} -- es zeichnete sich also eine klare Trennung nach Intensiviererklasse ab. Trotzdem war es auch bei {Aufgabe II} erstaunlich, dass die ausgedrückte Sprechereinstellung bei der unmarkierten Stufe nicht, wie eigentlich erwartet, von den Versuchspersonen als niedriger wahrgenommen wurde. Dies erhärtet den oben dargelegten Verdacht, dass die Gestaltung der Skala eventuell auf die Einschätzung der Items {eingewirkt hat}. In diesem Kontext besteht die Möglichkeit, dass die Versuchspersonen die Reichweite der verwendeten Skalen fälschlicherweise ähnlich interpretiert haben: Während die Skala zur Untersuchung von Skalarität von 'normal' bis 'maximal' reichte, reichte die Skala zur Untersuchung von Sprechereinstellung von 'gar nicht' bis 'maximal'. Vor diesem Hintergrund ist denkbar, dass die Versuchspersonen die linken Skalenränder einander angeglichen und keinen Unterschied bei der Interpretation der Skalen {gemacht haben}.



%Auffällig war zudem, dass \textit{arg} neben einem niedrigen Skalaritätsgrad ebenso mit einem niedrigen Grad an Sprechereinstellung verbunden und somit weder als übermäßig skalar noch als stark expressiv empfunden wurde. Bedenkt man, dass es sich hierbei um einen Ausdruck mit klar negativem Ursprung handelt, der im Korpus, wie oben beschrieben, zum einen über eine hohe Frequenz verfügt und zum anderen auch seit langer Zeit als Intensivierer existiert, vermag dieses Ergebnis kaum zu überraschen. Dennoch trifft dies gleichfalls auf andere Vertreter der \glqq{}schrecklichen\grqq{} Intensivierergruppe zu, die dagegen als deutlich expressiver bewertet wurden (bspw. \textit{furchtbar} oder \textit{schrecklich}). Dadurch wird klar, dass Frequenz und Persistenz nicht die alleinigen Gründe für die Einschätzung von \textit{arg} sein können. Ein möglicher Erklärungsansatz ist, dass \textit{arg} unterdessen semantisch ausgebleicht ist und daher als weniger ausdrucksstark empfunden wird als die anderen Ausdrücke. Die Frage, ob diese Vermutung der Realität entspricht, muss an dieser Stelle allerdings offenbleiben. Ähnlich niedrig wurde die Sprechereinstellung bei den expressiven Intensivierern \textit{voll}, \textit{wahnsinnig}, \textit{ober} und \textit{unheimlich} beurteilt. Dabei waren vor allem \textit{voll} und \textit{ober} auffällig, deren Skalaritätsgrad ebenfalls als niedrig wahrgenommen wurde. Dies liefert einen Grund zur Annahme, dass die im Korpus zwar niedrigfrequenten (36,65 bzw. 20,53 Treffer pro Million Textwörter), jedoch lange verwendeten Ausdrücke offenbar pragmatisch zu abgenutzt sind, um als stark skalar oder expressiv empfunden zu werden. Darüber hinaus legt die Überschneidung deren niedrigen Grade einen korrelativen Zusammenhang zwischen Skalarität und Sprechereinstellung nahe. Von besonderem Interesse war der Ausdruck \textit{derb}, der im Korpus der einzige Intensivierer mit über die Zeit rückläufigem Gebrauch war. Die diachrone Entwicklung wurde im Rahmen der Ergebnisinterpretation der Korpusuntersuchung als hinreichendes Zeichen für einen Expressivitätsverlust gesehen. Weil der durch \textit{derb} ausgedrückte Grad an Sprechereinstellung als nicht sonderlich hoch bewertet wurde, ist davon auszugehen, dass der Ausdruck grundsätzlich als nicht stark einstellungsausdrückend wahrgenommen wird. Wie bei \textit{voll} und \textit{ober} ist anzunehmen, dass der Intensivierer inzwischen zu abgenutzt und in der Folge nicht mehr expressiv genug ist. Obwohl die diachrone Entwicklung von \textit{derb} in einer früheren Zeitstufe aufgrund der zeitlichen Begrenzung des Korpus genauso unklar ist wie der damalige Grad an Sprechereinstellung, lässt sich dieses Ergebnis in Kombination mit dem kontinuierlichen Frequenzrückgang als fortgeschrittener Expressivitätsverlust deuten. Weil \textit{derb} im Korpus jedoch bereits zu Beginn des Untersuchungszeitraums niedrigfrequent war, muss der dem Anschein nach andauernde Expressivitätsschwund durch andere Faktoren (bspw. die Persistenz) zustande gekommen sein. An dieser Stelle kommt die Frage auf, wieso \textit{derb}, wenn es denn verbraucht ist, nicht mit einem noch niedrigeren Grad an Sprechereinstellung assoziiert wurde. Dabei ist vorstellbar, dass der Ausdruck den Versuchspersonen in intensivierender Funktion weniger geläufig, er im Experiment somit überraschender war und deswegen im direkten Vergleich mit den restlichen Ausdrücken als expressiver empfunden wurde. Dies spricht sich für die Möglichkeit eines erneuten Expressivitätsanstiegs aus, bei dem ein abgenutzter, allerdings selten verwendeter Intensivierer im Laufe der Zeit wieder als expressiv(er) wahrgenommen wird -- jedenfalls sofern der Ausdruck im Sprachgebrauch nicht völlig obsolet geworden ist.\footnote{Vgl. die historische Entwicklung der ehemaligen Intensivierer \textit{bäumig}  bzw. \textit{boomig}, \textit{gletscherhaft} und \textit{pyramidal} (vgl. \citealt[135]{Tobler.1868}; {\citealt[7f.]{Mueller.1899}}; {\citealt[268]{Becher.1907}}).} Im Gegensatz dazu war das im Korpus niedrigfrequente \textit{scheiß} der Ausdruck, der von den Versuchspersonen als am expressivsten eingestuft wurde (41,96 Treffer pro Million Textwörter). Ein Grund hierfür findet sich mutmaßlich in der dem Ausdruck zugrunde liegenden klar erkennbaren Semantik, durch die ihm eine negative Konnotation anhaftet. Diese hat gegebenenfalls dazu geführt, dass er als vergleichsweise ausdrucksstärker bewertet wurde und demnach besser geeignet zu sein scheint, um seine persönliche Einstellung offensiv zum Ausdruck zu bringen. Daher ist denkbar, dass die semantische Transparenz eines Intensivierers prinzipiell zu einem als höher wahrgenommenen Grad an Sprechereinstellung führt.
%Dass die restlichen expressiven Intensivierer in Bezug auf den ihnen zugesprochenen Skalenwert wiederholt recht nahe beinander lagen, könnte dem uniformen Untersuchungsaufbau geschuldet sein: So ging es für die Versuchspersonen bei den Aufgaben I und II alleine darum, die Items der Reihe nach auf einer Skala zu verorten. Daher ist nicht auszuschließen, dass sich die Uniformität der Aufgabenstellung trotz unterschiedlichen lexikalischen Materials in ähnlichen Skalenwerten niedergeschlagen hat.


%Bei den Ergebnissen zur in Aufgabe II abgefragten Sprechereinstellung kristallisierte sich auch hier die zuvor skizzierte Rangordnung heraus: expressive Intensivierer $\succ$ deskriptive Intensivierer $\succ$ unintensivierte Sätze.



\subsubsection[Einfluss von Sprechereinstellung auf Skalarität]{Beeinflusst die Sprechereinstellung die Skalarität?}


Bezüglich Skalarität (Aufgabe I) und Sprechereinstellung (Aufgabe II) schloss sich die Frage einer Korrelation an. Auf Grundlage der Definition von Expressivität nach \citet[133]{Gutzmann.2019a} ist vorstellbar, dass es einen systematischen Zusammenhang zwischen beiden Variablen gibt, sodass ein hoher Grad an Sprechereinstellung immer von einem hohen Grad an Skalarität begleitet wird. Das signifikante Ergebnis der Korrelationsanalyse bestätigt diese Annahme: Wenn der Grad an Sprechereinstellung von den Versuchspersonen als hoch empfunden wurde, wurde auch der Grad an Skalarität höher eingeschätzt. Dies lässt den Schluss zu, dass die Expressivität der Ausdrücke auf deren Skalarität einwirkt, die in der Folge als signifikant höher wahrgenommen wird. Dieser Befund steht zudem im Einklang mit der in der Literatur vorherrschenden Auffassung, dass Expressivität gleichermaßen mit einem hohen Grad an Skalarität und Sprechereinstellung einhergeht und stützt somit ebenfalls die zu klärenden {Hypothesen H3 und H4} nach \citet[22]{Paradis.1997} und {\citet[133]{Gutzmann.2019a}} sowie \citet[315]{Lang.1983} und {\citet[57]{Ortner.2014}}. Von diesem Resultat ausgehend lässt sich an dieser Stelle deutlich festhalten, dass die beiden Variablen \textsc{Skalarität} und \textsc{Sprechereinstellung} in diesem Experiment in der Tat miteinander korrelieren.


Um den Einfluss von Sprechereinstellung auf Skalarität zu überprüfen, wurde ergänzend zur Korrelationsanalyse ein lineares Regressionsmodell gerechnet. Dessen Ergebnis nach hat die Sprechereinstellung tatsächlich einen signifikanten Effekt auf die Skalarität (vgl. Tabelle \ref{LM01}). Daraus lässt sich schlussfolgern, dass der Grad an Skalarität mit steigendem Grad an Sprechereinstellung größer wird. Dieser Befund stimmt mit der zur Diskussion stehenden Vermutung von {\citet[133]{Gutzmann.2019a}} überein, auch wenn er keine Aussage über deren grundsätzliche Gültigkeit macht. Dessen ungeachtet ist davon auszugehen, dass der beobachtete Einfluss faktisch durch die Sprechereinstellung bedingt gewesen ist und sich die Expressivität der Ausdrücke daher wie erwartet in deren Skalaritätsgrad {niedergeschlagen hat}.




\subsection{Aufgabe III -- Einstellungspolarität}
Aufgabe III erfragte zudem die Polarität der Sprechereinstellung. Um ihre Einschätzung abzugeben, teilten die Versuchspersonen 100 virtuelle Taler auf die Auswahlfelder `Positiv', `Neutral' und `Negativ' einer komparativen Konstant-Summen-Skala auf. Aufgrund des Umstands, dass die Personen gezwungen waren, die Taler restlos auf die Optionen zu verteilen, kamen feinkörnige Abstufungen zustande, die Rückschlüsse auf die Rolle der Expressivität bei der Beurteilung der Einstellungspolarität zuließen. Die Grundannahme bestand gemäß der Hypothese H5 im Sinne von {\citet[135]{Gutzmann.2019a}} darin, dass expressive Intensivierer die evaluative Einstellung des Sprechers zum Geäußerten ausdrücken, die als positiv oder negativ zu interpretieren ist. Zur statistischen Überprüfung wurde mit R ein lineares gemischtes Modell berechnet, das das Referenzlevel `Neutral' mit `Positiv' und `Negativ' verglich.


Die Ergebnisse der Untersuchung lassen den Schluss zu, dass sowohl ein hoher Grad an Sprechereinstellung als auch ein hoher Grad an Skalarität zu einer signifikanten Abweichung von einer als neutral wahrgenommenen Sprechereinstellung führt (vgl. Tabelle \ref{LinReg0004}). Mit anderen Worten: Je höher der Grad an Skalarität und Sprechereinstellung eines Intensivierers eingeschätzt wird, desto stärker wird die Sprechereinstellung als positiv oder negativ empfunden. Dabei ist auffällig, dass die Beurteilung der Einstellungspolarität bei vielen expressiven Intensivierern deutlich stärker in Richtung `Negativ' tendiert, was sich in einem Übergewicht der negativen Polarität äußert (vgl. Tabelle \ref{Bedin}). Dieses Ergebnis ist möglicherweise auf die häufig transparente negative Semantik der expressiven Ausdrücke (z.\,B. \textit{fürchterlich} oder \textit{verdammt}) zurückzuführen, die sich dem Anschein nach auf die Beurteilung der Sprechereinstellung auswirkt.
%Unerwartet war allerdings, dass der expressive Intensivierer \textit{unheimlich} im Hinblick auf die Einstellungspolarität von den Versuchspersonen als weitgehend ausgewogen beurteilt wurde, ist der Ausdruck wie die anderen Vertreter der \glqq{}schrecklichen\grqq{} Intensivierergruppe doch ebenso wahrnehmbar negativ. Im Korpus war \textit{unheimlich} hochfrequent (161,89 Treffer pro Million Textwörter) und lange im Gebrauch. Demzufolge ist nicht ausgeschlossen, dass seine Semantik in intensivierender Funktion in der Zwischenzeit verblasst ist, wie {\citet[15]{Nouwen.2020}} am Beispiel des englischen Intensivierers \textit{terribly} bespricht. Das Ausgebleichtsein kann also dazu geführt haben, dass ein mit \textit{unheimlich} modifizierter Intensivausdruck entgegen den Erwartungen in Bezug auf die Einstellungspolarität als weitgehend ausgewogen interpretiert wird.
Bei den als positiv bewerteten expressiven Intensivierern handelt es sich mit Ausnahme von \textit{ober} um solche, die als Interjektionen bzw. innerhalb von Ausrufesätzen in prädikativer Verwendung vorkommen können und die in dieser Stellung ausschließlich positiv konnotiert sind: \textit{Super!} bzw. \textit{Das ist ja mega!} Der Ausdruck \textit{ober} verfügt hingegen über eine offenkundig dimensionale, aufwärtsgerichtete Bedeutung, die der traditionellen Metapherntheorie nach \citet[14ff.]{Lakoff.1980} zufolge ebenfalls mit Positivem in Verbindung steht. Auch dies legt einen Einfluss der Semantik und der damit verbundenen Konnotation des Ursprungslexems nahe. Anders als die expressiven Intensivierer werden der deskriptive Intensivierer \textit{sehr} und die unmarkierte Variante als überwiegend neutral empfunden, sprich letztere drücken nach Einschätzung der Versuchspersonen keine evaluative Sprechereinstellung aus.  Dieser Befund erhärtet die zu überprüfende {Hypothese H5} nach {\citet[135]{Gutzmann.2019a}} und beweist außerdem, dass die Erhebung der Einstellungspolarität mittels komparativer Konstant-Summen-Skala als Untersuchungsmethode funktioniert hat.


Zusammengefasst weisen die Ergebnisse von Aufgabe III darauf hin, dass expressive Intensivierer, die den Resultaten der Untersuchung zufolge im Allgemeinen sowohl mit einem signifikant höheren Grad an Skalarität als auch einem signifikant höheren Grad an Sprechereinstellung einhergehen, faktisch die evaluative Einstellung des Sprechers zum beschriebenen Sachverhalt kommunizieren. Die ausgedrückte Sprechereinstellung kann vom Adressaten demnach gleichermaßen positiv wie negativ aufgefasst werden. Dies spricht klar für das Zutreffen der Hypothese H5 und kann somit als erste empirische Evidenz für die Annahme herangezogen werden, dass die Einstellungsspiegelung, die expressive Intensivierer inhärent zum Ausdruck bringen, evaluativen Charakters und somit alles andere als neutral ist.





\subsection{Die Rolle der Frequenz}

\largerpage
Ergänzend zu den Merkmalen Skalarität, Sprechereinstellung und Einstellungspolarität wurde im Experiment der Rolle von Frequenz nachgegangen. Die zu testende Hypothese {H1-Freq} stammt von {\citet[2]{Hauschild.1899}} und besagt, dass sich ein häufiger Gebrauch von Intensivierern nachteilig auf ihre Expressivität auswirkt. Davon ausgehend ist in erster Linie die Untersuchung der in Aufgabe II erfragten Sprechereinstellung relevant. Dennoch besteht hinsichtlich Skalarität die Möglichkeit, dass expressive Intensivierer mit frequentem Gebrauch zwar an Expressivität verlieren, gleichsam jedoch eine eher skalierende Bedeutung annehmen.\footnote{Eine solche Entwicklung könnte auch \textit{sehr} durchlaufen haben, dessen Skalaritätsgrad von den Versuchspersonen ähnlich hoch wie der der expressiven Intensivierer {eingestuft wurde}.} Durch den Expressivitätsverlust sind die Ausdrücke eventuell textsortenadäquater und können auch in als sprachlich konservativ geltenden Textformen verwendet werden, was unter Umständen eine noch höhere Frequenz nach sich zieht.\footnote{Diese Idee wurde im Rahmen der Korpusuntersuchung als mögliche Erklärung für die hohe Gebrauchsfrequenzen der \glqq{}schrecklichen\grqq{} Intensivierer im Zeitschriftenkorpus diskutiert.} Dies würde dazu führen, dass manche abgenutzten Intensivierer zwar mit einem niedrigeren Grad an Sprechereinstellung in Zusammenhang gebracht werden, der durch sie ausgedrückte Skalaritätsgrad allerdings beibehalten bleibt. Um Hinweise auf die Rolle der Frequenz zu erhalten, wurden die expressiven Intensivierer gemäß der im Korpus ermittelten Gebrauchsfrequenz des letzten Zeitabschnitts (d.\,h. dem Zeitraum von 1990 bis 2017) in eine niedrig- und eine hochfrequente Gruppe eingeteilt (vgl. Tabelle \ref{Fre3.1}), die als Grundlage für die statistischen Auswertungen des Experiments {herangezogen wurden}.


Die Ergebnisse zur in Aufgabe I untersuchten Skalarität zeigen in diesem Punkt einen signifikanten Unterschied zwischen den beiden Frequenzgruppen (vgl. Tabelle \ref{LinRea0100}). Dieser äußert sich darin, dass die niedrigfrequenten Intensivierer mit einem bedeutend niedrigeren Skalaritätsgrad assoziiert werden als die hochfrequenten Intensivierer. Dieser Befund legt die Korrektheit der soeben beschriebenen Vermutung nahe, dass die wahrnehmbare Skalarität eines Ausdrucks zunimmt, je häufiger er im Gebrauch ist. Auch die Ergebnisse der Regressionsanalyse zur in {Aufgabe II} ermittelten Sprechereinstellung sind signifikant (vgl. {Tabelle \ref{LinRegzaz0}}). Hier verbanden die Versuchspersonen die niedrigfrequenten Intensivierer grundsätzlich mit einem höheren Grad an Sprechereinstellung als die hochfrequenten Intensivierer. Dies deutet darauf hin, dass eine hohe Frequenz tatsächlich negative Auswirkungen auf den mit den Ausdrücken assoziierten Grad an Sprechereinstellung hat. Dieser Befund spricht eindeutig für das Zutreffen der {Hypothese H1-Freq} nach {\citet[2]{Hauschild.1899}}, nach der eine hohe Frequenz von Ausdrücken abträglich für ihre Expressivität ist. Erklärbar ist dies vermutlich damit, dass ein häufig gebrauchter Intensivierer vorhersehbarer, daher weniger überraschend und aufsehenerregend ist als ein selten gebrauchter Intensivierer. Darüber hinaus bietet dies auch eine potenzielle Erklärung für den unerwartet hoch eingeschätzten Grad an Sprechereinstellung des dem Anschein nach abgeschliffenen expressiven Intensivierers \textit{derb}: Auf Grundlage der Untersuchungsergebnisse ist denkbar, dass der Ausdruck den Versuchspersonen in intensivierender Funktion weniger geläufig war, er von ihnen daher als expressiver wahrgenommen wurde und sie ihm in der Konsequenz einen vergleichsweise höheren Grad an Sprechereinstellung zugesprochen haben.


\subsection{Die Rolle der Persistenz}
\largerpage
Neben der Frequenz wurde der Faktor Zeit berücksichtigt. Wie {\citet[59]{Tobler.1868}} bemerkt, verlieren Intensivierer sukzessive an Expressivität, je länger sie im Gebrauch und desto weniger aktuell sie sind. Klar ist, dass auch hier -- ähnlich wie bei Frequenz -- primär der in {Aufgabe II} ermittelte Grad an Sprechereinstellung im Vordergrund steht, während die Ergebnisse zur in Aufgabe I untersuchten Skalarität lediglich bedingt Rückschlüsse auf die Rolle der Persistenz zulassen. Zur inferenzstatistischen Überprüfung wurden analog zum oben beschriebenen Vorgehen basierend auf den Ergebnissen der Korpusstudie zwei Persistenzgruppen gebildet. Im Zuge dessen wurden diejenigen Intensivierer, die erst im Laufe des Untersuchungszeitraums erstmals im Korpus aufgetreten sind, von den restlichen expressiven Intensivierern separiert: \textit{hammer}, \textit{krass} und \textit{mega}.\footnote{Diese Ausdrücke wurden ferner als Evidenz für die Gültigkeit der Hypothese H2 gewertet, nach der es infolge des Expressivitätsverlusts einen ständigen Bedarf an neuen Ausdrücken gibt, die den Ausfall abgenutzter Intensivierer kompensieren {\citep[vgl.][274]{Baumgarten.1908}}.}


Die Ergebnisse zur in Aufgabe I erfragten Skalarität sind in dem Sinne signifikant, dass die neuen Intensivierer von den Versuchspersonen mit einem bedeutend höheren Grad an Skalarität assoziiert wurden als die älteren Intensivierer (vgl. Tabelle \ref{LinRea0100}). Dies liegt mutmaßlich im Neuheitsgrad ersterer begründet, in dessen Folge ein im Kontext eher überraschender Ausdruck einen als höher empfundenen Skalaritätsgrad erhält. Je neuer und weniger vorhersehbar ein Intensivierer also ist, desto höher ist der Grad des durch ihn bezeichneten Merkmals anzusehen. Im Gegensatz dazu liefern die Ergebnisse zur in {Aufgabe II} untersuchten Sprechereinstellung keine empirische Evidenz für die Annahme, dass die älteren Intensivierer prinzipiell weniger expressiv sind als die neuen Intensivierer (vgl. {Tabelle \ref{LinRegzaz0}}), weshalb die {Gültigkeit} der Hypothese H1-Temp nach {\citet[59]{Tobler.1868}} verworfen werden muss. Aufgrund des Umstands, dass die Gruppe der neuen Intensivierer mit lediglich drei Mitgliedern (d.\,h. \textit{hammer}, \textit{krass} und \textit{mega}) stark {unterrepräsentiert} ist, ist dieser Befund allerdings mit Vorsicht zu interpretieren. Während diese Ausdrücke im Korpus erst im Laufe des Untersuchungszeitraums erstmals in intensivierender Funktion aufgetreten sind, ließen sich die restlichen Ausdrücke {(n = 13)} schon frühzeitig darin nachweisen. Dies resultierte im Experiment wiederum in einem Ungleichgewicht zwischen den Persistenzgruppen, das mit hoher Wahrscheinlichkeit auf die inferenzstatistischen Auswertungen eingewirkt hat. Da einige der älteren Intensivierer (z.\,B. \textit{fürchterlich} oder \textit{scheiß}) als mindestens genauso einstellungsausdrückend wahrgenommen wurden wie die neuen, ist anzunehmen, dass die stark ausgeprägte Dominanz von ersteren den erwarteten Effekt überlagert und ihm dadurch entgegengewirkt hat. Demnach besteht die Möglichkeit, dass das Experiment mit einem ausbalancierten Verhältnis aus alten und neuen Ausdrücken durchaus zu einem signifikanten Ergebnis geführt hätte. So ist es trotz des aufgetretenen Nulleffekts keinesfalls abwegig, dass abgesehen von häufigen auch lange verwendete Intensivierer weniger überraschend und daher als weniger expressiv empfunden werden als neuere Intensivierer. Das Zutreffen dieser Mutmaßung bleibt jedoch über die Untersuchung hinaus ungeklärt und kann allenfalls in Folgestudien durch ein ausgewogenes Verhältnisses aus alten und neuen Ausdrücken {überprüft werden}.




\subsection[Zur Expressivität von \textit{derb}]{Wie steht es um die Expressivität von \emph{derb}?}
\largerpage
Von besonderem Interesse war der Ausdruck \textit{derb}, der im Korpus der einzige Intensivierer mit über die Zeit rückläufigem Gebrauch war. Die diachrone Entwicklung wurde im Rahmen der Ergebnisinterpretation der Korpusuntersuchung als hinreichendes Zeichen für einen Expressivitätsverlust gesehen. Weil der durch \textit{derb} ausgedrückte Grad an Sprechereinstellung als nicht sonderlich hoch bewertet wurde, ist davon auszugehen, dass der Ausdruck grundsätzlich als nicht stark einstellungsausdrückend wahrgenommen wird. Von diesem Befund ausgehend ist anzunehmen, dass der Intensivierer inzwischen zu abgenutzt und in der Folge nicht mehr expressiv genug ist. Obwohl die diachrone Entwicklung von \textit{derb} in einer früheren Zeitstufe aufgrund der zeitlichen Begrenzung des Korpus genauso unklar ist wie der damalige Grad an Sprechereinstellung, lässt sich dieses Ergebnis in Kombination mit dem kontinuierlichen Frequenzrückgang als fortgeschrittener Expressivitätsverlust deuten. Weil \textit{derb} im Korpus jedoch bereits zu Beginn des Untersuchungszeitraums niedrigfrequent war, muss der dem Anschein nach andauernde Expressivitätsschwund durch andere Faktoren (bspw. die Persistenz) zustande gekommen sein. An dieser Stelle kommt die Frage auf, wieso \textit{derb}, wenn es denn verbraucht ist, nicht mit einem noch niedrigeren Grad an Sprechereinstellung assoziiert wurde. Dabei ist vorstellbar, dass der Ausdruck den Versuchspersonen heute in intensivierender Funktion weniger geläufig, er im Experiment somit überraschender war und deswegen im direkten Vergleich mit den restlichen Ausdrücken als expressiver empfunden wurde. Dies legt die Möglichkeit eines erneuten Expressivitätsanstiegs nahe, bei dem ein abgenutzter, allerdings selten verwendeter Intensivierer im Laufe der Zeit wieder als expressiv(er) wahrgenommen wird\footnote{Die Beobachtung, dass Intensivierer an Popularität verlieren, ihre Häufigkeit dadurch abnimmt, und sie erst sehr viel später auf bemerkenswerte Weise „wiederbelebt“ werden, ist keineswegs neu. Eine solche Entwicklung ist für den britischen Ausdruck \textit{well} dokumentiert: Wie \citet[]{Stratton.2020d,Stratton.2022a} beschreibt, ist das intensivierende \textit{well} bereits im Alt- und Mittelenglischen belegt. Seine Frequenz ging jedoch zurück, sodass der Ausdruck nur noch randständig in peripheren Varietäten mit äußerst geringer Häufigkeit auftrat, ehe er später im britischen Englisch wieder häufiger verwendet wurde \citep[vgl.][]{Stratton.2018,Stratton.2020d,Stratton.2020a}} -- jedenfalls sofern der Ausdruck im Sprachgebrauch nicht völlig obsolet {geworden ist.}\footnote{Siehe die historische Entwicklung der ehemaligen Intensivierer \textit{bäumig}  bzw. \textit{boomig}, \textit{gletscherhaft} und \textit{pyramidal} (vgl. \citealt[135]{Tobler.1868}; {\citealt[7f.]{Mueller.1899}}; {\citealt[268]{Becher.1907}}).}

%
\section{Fazit}\label{abs:Fazit835}
Das in diesem Kapitel vorgestellte Experiment beruhte auf der Intention, Rückschlüsse auf die noch offenen Fragen der Korpusuntersuchung ziehen zu können. Vor diesem Hintergrund wurde der Gültigkeit von zwei Hypothesen vertieft auf den Grund gegangen. Die erste Hypothese geht u.\,a. auf {\citet[2]{Hauschild.1899}} zurück, der postuliert, dass eine häufige Verwendung von Intensivierern zulasten ihrer Expressivität gehe {(vgl. H1-Freq)}. Damit verbunden scheint der Faktor Zeit zu sein, wie schon {\citet[59]{Tobler.1868}} erkennt und der meinem Wissen nach bislang keine Beachtung fand ({vgl. H1-Temp)}. Demnach hat nicht bloß die Häufigkeit eines Ausdrucks einen negativen Effekt auf die Expressivität, sondern ebenso die Dauer seines Gebrauchs. Auf Basis der Ergebnisse der Korpusanalyse wurde erwogen, dass die im Korpus sowohl hochfrequenten wie auch älteren Intensivierer grundsätzlich weniger expressiv sind als die niedrigfrequenten oder neuen Intensivierer. Aus diesem Grund zeichnet sich in Bezug auf Skalarität und Sprechereinstellung vermutlich ein Unterschied zwischen diesen Gruppen ab. Zusätzlich wurde in der Untersuchung erstmals dem wahrnehmbaren Unterschied zwischen deskriptiven und expressiven Intensivierern nachgegangen. So wird in der theoretischen Literatur die Annahme vertreten, dass Expressivität gleichermaßen mit einem höheren Grad an Skalarität und Sprechereinstellung einhergeht (vgl. \citealt[315]{Lang.1983}; \citealt[22]{Paradis.1997}; \citealt[57]{Ortner.2014}; \citealt[133]{Gutzmann.2019a}). In diesem Sinne drücken expressive Intensivierer mutmaßlich nicht nur einen höheren Grad einer Merkmalsausprägung aus, sondern geben ferner auch die persönliche Einstellung des Sprechers zum beschriebenen Sachverhalt wieder (vgl. H3 und H4). Im Hinblick auf die in der Hypothese H4 angenommene Sprechereinstellung war ferner die Einstellungspolarität von Interesse, die in einer weiterführenden Hypothese beleuchtet wurde (vgl. H5). In diesem Zusammenhang ist {\citet[135]{Gutzmann.2019a}} folgend anzunehmen, dass die subjektive Sprechereinstellung, die expressive Intensivierer kommunizieren, evaluativer Natur ist und deshalb als positiv oder negativ verstanden wird, wohingegen deskriptive Intensivierer oder unintensivierte Sätze aufgrund des Fehlens an Expressivität grundsätzlich stärker als neutral empfunden werden.


%In diesem Sinne hatte jede Versuchsperson im Experiment drei Aufgaben zu erfüllen: Bezog sich die erste auf die Skalarität und erfragte die Merkmalsausprägung eines im Satz thematisierten Gegenstands, ging es bei der zweiten um den Grad der Einstellung, die der Sprecher zum beschriebenen Sachverhalt hat. Die dritte Aufgabe kam in Gestalt einer komparativen Konstant-Summen-Skala daher, bei der die Versuchspersonen {100 Taler} gemäß Einschätzung der Einstellungspolarität auf die zur Verfügung stehenden Optionen 'positiv', 'neutral' und 'negativ' zu {verteilen hatten}.


Als erstes Untersuchungsergebnis kann festgehalten werden, dass sich die drei Stufen sowohl hinsichtlich der in Aufgabe I erfragten Skalarität als auch der in Aufgabe II untersuchten Sprechereinstellung erwartungsgemäß {verhalten}: Brachten die Versuchspersonen Sätze mit expressivem Intensivierer in beiden Fällen mit dem höchsten Grad in Verbindung, so empfanden sie den Grad von Sätzen mit deskriptivem Intensivierer als niedriger, genauso wie sie den Grad der unintensivierten Sätze entsprechend niedriger einschätzten. Besonders stark ausgeprägt ist der Unterschied bei der in {Aufgabe II} untersuchten Sprechereinstellung. Dies steht im Einklang mit der in der Literatur vorherrschenden Annahme, dass vor allem -- wenn nicht sogar ausschließlich -- expressive Intensivierer eine Einstellungsbekundung des Sprechers zu dem im Satz beschriebenen Sachverhalt kommunizieren \citep[vgl.][133]{Gutzmann.2019a}. Dieser Befund deutet darauf hin, dass die zu überprüfenden Hypothesen H3 und H4 nach \citet[22]{Paradis.1997} und \citet[133]{Gutzmann.2019a} sowie \citet[315]{Lang.1983} und {\citet[57]{Ortner.2014}} zutreffend zu sein scheinen: Wie die Ergebnisse der Untersuchung demonstrieren, gehen expressive Intensivierer in der Tat mit einem signifikant höheren Grad an Skalarität und Sprechereinstellung einher, was aller Wahrscheinlichkeit nach der Hauptgrund für ihren pragmatisch aufsehenerregenden Charakter ist.


%Von besonderem Interesse war \textit{derb}, der im Korpus der einzige Intensivierer war, der über die Zeit stetig seltener gebraucht wurde und einen hinreichenden Grund zur Annahme des in H1 postulierten Expressivitätsverlusts nach \citet[2]{Hauschild.1899} bot. Somit stellte sich die Frage, wie \textit{derb} von den Versuchspersonen mit Blick auf Skalarität und Sprechereinstellung bewertet. Die Ergebnisse des Experiments lieferten an dieser Stelle keine klaren Hinweise darauf, dass sich die Wahrnehmung von \textit{derb} sowohl bei Skalarität als auch bei Sprechereinstellung signifikant von der Wahrnehmung der restlichen Ausdrücke unterschied. Dennoch wurde \textit{derb} von den Versuchspersonen weder als übermäßig skalar noch als stark einstellungsausdrückend wahrgenommen. Eine Erklärung findet sich möglicherweise in der Gebrauchsfrequenz des intensivierenden \textit{derb}. Wie die Ergebnisse der Korpusanalyse nahelegten, scheint der Ausdruck unterdessen nahezu ausgestorben zu sein. Demnach ist davon auszugehen, dass er für die Versuchspersonen im direkten Vergleich mit den anderen Ausdrücken neu und überraschend war, daher als aufsehenerregender empfunden und in der Folge mit einem höheren Grad assoziiert wurde, als aufgrund des Expressivitätsverlust zu erwarten war. Wie oben beschrieben, legt dies die Möglichkeit einer erneuten Expressivitätszunahme nahe -- sofern der Ausdruck nicht völlig ausgestorben und er vereinzelt weiterhin im Gebrauch ist.


Außerdem wurde im Experiment im Rahmen der Hypothese {H1-Freq} der Einfluss von Frequenz überprüft. Nach dieser verlieren Intensivierer ihre Expressivität, wenn sie zu häufig verwendet werden {\citep[vgl.][2]{Hauschild.1899}}. Den Ergebnissen nach ist auch diese Hypothese gültig, werden die niedrigfrequenten Intensivierer doch mit einem deutlich höheren Grad an Sprechereinstellung assoziiert als die hochfrequenten Intensivierer. Dies liefert die erste empirische Evidenz für einen Expressivitätsverlust von häufig gebrauchten Ausdrücken. Zudem geben die Untersuchungsergebnisse die Möglichkeit einer erneuten Expressivitätszunahme zu bedenken. So wurde der dem {Anschein} nach abgeschliffene und im Korpus beinahe ausgestorbene Intensivierer \textit{derb} von den Versuchspersonen zwar weder als übermäßig skalar noch als stark expressiv eingeschätzt. Dennoch wurde er in Bezug auf die in Aufgabe II abgefragte Sprechereinstellung nicht so niedrig bewertet, wie man es infolge des Expressivitätsverlusts erwarten würde. Daher ist zu erwägen, dass der inzwischen selten gebrauchte Intensivierer für die Versuchspersonen im direkten Vergleich mit den restlichen Ausdrücken neu und überraschend war, somit als aufmerksamkeitsstärker empfunden und in der Folge mit einem vergleichsweise höheren Grad an Sprechereinstellung assoziiert wurde. Ist ein abgenutzter Ausdruck im Sprachgebrauch also nicht völlig obsolet und weiterhin vereinzelt im Gebrauch, ist nicht ausgeschlossen, dass er nach längerer Zeit wieder als expressiv(er) wahrgenommen wird.\footnote{Eine solche Entwicklung scheint der britische Ausdruck \textit{well} über die Jahrhunderte durchlaufen zu haben, wie die Ausführungen von \citet[]{Stratton.2018,Stratton.2020d,Stratton.2020a,Stratton.2022a} nahelegen. Dessen ungeachtet ist es nach Auffassung von \citet[95]{Dahl.2001} so gut wie ausgeschlossen, dass ein verbrauchter Ausdruck in seiner ursprünglichen Funktion wiederbelebt wird. Stattdessen sei es ökonomischer, diesen durch einen neuen Ausdruck zu ersetzen.}


Neben der Frequenz wurde ebenso die Persistenz berücksichtigt, die auf die Überprüfung der Hypothese H1-Temp abzielte. Diese geht im Kern auf \citet[59]{Tobler.1868} zurück, der postuliert, dass ein Intensivierer im Laufe der Zeit an Expressivität einbüße. Die Ergebnisse des Experiments bieten keine Anhaltspunkte für das Zutreffen von H1-Temp. Aufgrund des bestehenden Ungleichgewichts, das sich in Bezug auf den Umfang der Persistenzgruppen abzeichnet (d.\,h. drei neue Intensivierer vs. 13 ältere Intensivierer), ist dieses Ergebnis allerdings nur wenig repräsentativ und mit Vorsicht zu betrachten.


Bei Aufgabe III wurde vor dem Hintergrund der Hypothese H5 nach {\citet[135]{Gutzmann.2019a}} zudem die Polarität der Sprechereinstellung abgefragt. In diesem Kontext schlussfolgert er unter Verweis auf {\citet[681]{McCready.2008}}, dass expressive Intensivierer der Einstellungsspiegelung dienen, die evaluativ und folglich als positiv oder negativ zu beurteilen ist. Die Ergebnisse des Experiments sprechen für die Gültigkeit dieser Annahme, wurden doch vor allem diejenigen Intensivierer als von `Neutral' abweichend interpretiert, die mit einem hohen Grad an Skalarität und Sprechereinstellung verbunden wurden. Daneben wurden Sätze mit deskriptivem Intensivierer sowie unintensivierte Sätze von den Versuchspersonen vorwiegend als neutral empfunden. Dieser Befund liefert erste empirische Evidenz, dass expressive Intensivierer inhärent die evaluative Einstellung des Sprechers ausdrücken, was die Hypothese H5 untermauert.


Kurz gefasst erhärten die Ergebnisse des Experiments die Korrektheit nahezu aller zur Diskussion stehenden Hypothesen. So findet sich in dem Sinne Evidenz für die Hypothese {H1-Freq} nach {\citet[2]{Hauschild.1899}}, dass die im Korpus niedrigfrequenten Intensivierer tatsächlich mit einem höheren Grad an Sprechereinstellung in Verbindung gebracht werden als die im Korpus hochfrequenten Intensivierer. Auch bestätigen die Untersuchungsergebnisse den im Kontext von H3 und H4 angenommenen Unterschied zwischen den Intensiviererklassen (vgl. \citealt[315]{Lang.1983}; \citealt[22]{Paradis.1997}; \citealt[57]{Ortner.2014}; {\citealt[133]{Gutzmann.2019a}}): Sätze mit expressivem Intensivierer werden mit einem bedeutend höheren Grad an Skalarität und Sprechereinstellung assoziiert als Sätze mit deskriptivem Intensivierer oder unintensivierte Sätze. Obendrein konnte gezeigt werden, dass die durch expressive Intensivierer ausgedrückte Sprechereinstellung, wie {\citet[135]{Gutzmann.2019a}} beschreibt, evaluativer Natur ist und grundsätzlich als positiv oder negativ wahrgenommen wird. Keine empirische Evidenz findet sich hingegen für die {Hypothese H1-Temp} nach {\citet[59]{Tobler.1868}}, unterscheiden sich die im Korpus neu aufgetretenen Intensivierer doch im Hinblick auf den ihnen zugeschriebenen Grad an Sprechereinstellung nicht bedeutend von den älteren Intensivierern, was mit hoher Wahrscheinlichkeit der Imbalance der beiden Persistenzgruppen geschuldet und somit nur wenig aussagekräftig ist.


In Kapitel \ref{Experimente} wird eine weitere experimentelle Untersuchung vorgestellt. Darin wird getestet, ob es nicht nur zwischen, sondern auch innerhalb der Intensiviererklassen Unterschiede in Bezug auf den durch die Intensivierer ausgedrückten Intensitätsgrad gibt, wie sich aus der theoretischen Literatur (bspw. {\citealt[87]{Biedermann.1969}} oder {\citealt[408]{Breindl.2007}}) ableiten lässt. Um der zentralen Fragestellung gerecht zu werden, wird eine Stichprobe aus je vier deskriptiven und expressiven Intensivierern mit unterschiedlich wahrnehmbarer Intensität hinsichtlich des mit den Ausdrücken assoziierten Grades an Skalarität und Sprechereinstellung von Versuchspersonen in einem ähnlich aufgebauten Rating-Experiment kontrastiert. Der elementare Unterschied zwischen dem in diesem Kapitel vorgestellten und dem nachfolgend beschriebenem Experiment besteht daher vornehmlich darin, dass die Stichprobe im nächsten Experiment Intensivierer mit unterschiedlicher Ausdrucksstärke umfasst, um ebendiese potenziellen klasseninternen Abstufungen zu überprüfen.






% eine weiterführende Hypothese getestet, die sich implizit aus der Annahme eines wahrnehmbaren Unterschiedes zwischen deskriptiven und expressiven Intensivierern ergibt. So lässt sich aus der theoretischen Literatur (bspw. {\citet[87]{Biedermann.1969}} oder {\citet[408]{Breindl.2007}}) heraus ableiten, dass es Unterschiede in Bezug auf den durch die Intensivierer ausgedrückten Intensitätsgrad gibt und sich die Ausdrücke in Abhängigkeit ihrer Semantik diesbezüglich auf einer Skala anordnen lassen. Dabei besteht die Möglichkeit, dass die Expressivität Einfluss auf den Skalaritätsgrad nimmt, sich die Expressivität also in der Skalarität niederschlägt, und letztere infolgedessen als stärker empfunden wird {\citep[vgl.][133]{Gutzmann.2019a}}. Um der zentralen Fragestellung gerecht zu werden, wird eine Stichprobe aus je vier deskriptiven und expressiven Intensivierern mit unterschiedlich wahrnehmbarer Intensität hinsichtlich des mit ihnen assoziierten Grades an Skalarität und Sprechereinstellung von Versuchspersonen in einem ähnlich aufgebauten Rating-Experiment kontrastiert. Der elementare Unterschied zwischen dem in diesem Kapitel vorgestellten und dem nachfolgend dargelegten Experiment besteht daher vornehmlich darin, dass die Stichprobe im nächsten Experiment Intensivierer mit unterschiedlicher Ausdrucksstärke umfasst, um ebendiese potenziellen Abstufungen zu überprüfen. Die im Anschluss präsentierte Untersuchung setzt sich aus zwei separaten Teilen zusammen: Im ersten Teil ({Experiment 2a}) geht es um den durch die Intensivierer ausgedrückten Grad an Skalarität, im zweiten Teil {(Experiment 2b}) um den durch sie ausgedrückten Grad an {Sprechereinstellung}.




%% Ausgeschnitten, aber noch zu diskutieren!!
%\noindent Die Skalenrange reichte von 0 bis 100 Punkte ohne dass die Werte eingeblendet wurden. Vom Anzeigen sowohl des Werts der Skalenendpunkte, von festgelegten Skalenmarken wie auch der gewählten Skalenposition wurde abgesehen, um unerwünschte Priming-Effekte zu umgehen. So wurde bspw. die Möglichkeit in Betracht gezogen, dass Binnendifferenzierungen wie Skalenmarken oder segmentale Abstufungen die Versuchspersonen dazu verleiten würden, alle Ausdrücke gleichermaßen um ebendiese Marke(n) herum anzuordnen, statt sich bei der Beurteilung gänzlich auf ihre persönliche Intuition zu verlassen.


%Nichtsdestotrotz ist bei der unmarkierten Stufe auffällig, dass ihr sowohl im Kontext von Skalarität als auch von Sprechereinstellung von den Versuchspersonen eine vergleichsweise hohe Skalenposition zugeschrieben wurde. So würde man aufgrund des fehlenden Intensivierers vielmehr annehmen, dass unintensivierte Sätze weder als übermäßig skalar noch als stark expressiv wahrgenommen und diese in der Folge intuitiv im untersten Skalenbereich verortet werden. Die möglichen Gründe für diese unerwartet hohe Skaleneinschätzung werden in Abschnitt \ref{abs:Diskussion22} vertiefend betrachtet.














%Bei Aufgabe I war zu erkennen, dass sich die expressiven Intensivierer hinsichtlich des mit ihnen assoziierten Skalaritätsgrades nicht übermäßig stark voneinander unterschieden (vgl. Abbildung \ref{Skalar7}). Hervorgehoben seien an dieser Stelle diejenigen Intensivierer, die den höchsten und niedrigsten gemittelten Skalenwert aufwiesen: Während der hohe Wert bei \textit{mega} auf die anhaltende Aktualität des Ausdrucks zurückgeführt wurde, wurde der verhältnismäßig niedrige Wert bei \textit{arg} als Folgeerscheinung von Abnutzung erklärt (vgl. diesbezüglich auch \citealt[15]{Nouwen.2020}). Demnach besteht die Möglichkeit, dass das intensivierende \textit{arg} infolge seines erwiesenermaßen langen Gebrauchs nach gegenwärtigem Empfinden mit einem niedrigen Skalaritätsgrad assoziiert wird. Festzustellen war zudem, dass der deskriptive Intensivierer \textit{sehr} in Bezug auf Skalarität vergleichbar hoch positioniert wurde wie die expressiven Intensivierer, was in einem nur geringen Unterschied zwischen den Intensiviererklassen resultierte. Erklärt wurde dies durch die \textit{sehr} primär zugrunde liegende skalare Bedeutung, durch die der Ausdruck trotz fehlender Expressivität einen hohen Skalaritätsgrad angibt.


%Bei Aufgabe II war der Unterschied zwischen den beiden Intensiviererklassen stärker ausgeprägt. Den Ergebnissen nach wurde der Grad an Sprechereinstellung des deskriptiven Intensivierers \textit{sehr} von den Versuchspersonen auf der Skala bedeutend niedriger beurteilt als der der expressiven Intensivierer. Auch hier seien die expressiven Intensivierer mit dem höchsten und niedrigsten Skalenwert herausgestellt: Während es sich bei ersterem mit \textit{scheiß} um einen Ausdruck mit transparenter Semantik und negativer Konnotation handelt, ist der am niedrigsten gerankte expressive Intensivierer \textit{arg} semantisch mutmaßlich zu ausgebleicht und aufgrund eines hohen Lexikalisierungsgrades nicht mehr ausdrucksstark genug.


%Von besonderem Interesse war \textit{derb}, der im Korpus der einzige Intensivierer war, der über die Zeit weniger häufig gebraucht wurde. In der Folge bot er alleinigen Grund zur Annahme der Hypothese H1, die einen Expressivitätsverlust postuliert \citep[vgl.][2]{Hauschild.1899}. Somit stellte sich die Frage, wie \textit{derb} von den Versuchspersonen mit Blick auf Skalarität und Sprechereinstellung auf der Skala bewertet wurde. Wie in Abschnitt \ref{abs:Diskussion22} beschrieben, zeichnete sich bei beiden Aufgaben kein ausgeprägter Unterschied zwischen den expressiven Intensivierern -- darunter auch \textit{derb} -- ab. Obwohl \textit{derb} sowohl im Hinblick auf Skalarität als auch auf Sprechereinstellung höher als bspw. \textit{arg}, \textit{voll} und \textit{ober} eingeschätzt wurde, lag der Ausdruck dennoch bei beiden Aufgaben im unteren Skalenbereich. Dies bedeutet, dass \textit{derb} von den Versuchspersonen weder als übermäßig skalar noch als stark einstellungsausdrückend wahrgenommen wurde. Dabei besteht die Möglichkeit, dass das intensivierende \textit{derb} in einer früheren Zeitstufe mit einem höheren Expressivitätsgrad assoziiert wurde als es heute der Fall ist, was sich auf Grundlage der korpuslinguistischen und experimentellen Daten allerdings nicht begründen lässt. Fakt ist jedoch, dass \textit{derb} in intensivierender Funktion bei beiden Aufgaben eine niedrige Skalenposition einnahm und folglich als vergleichsweise wenig skalar und expressiv zu bewerten ist. Auch wenn sich an dieser Stelle keine stichfesten empirischen Evidenzen für einen Expressivitätsschwund finden, bietet die Rolle der Expressivität eine hinreichende Erklärung für die diachrone Frequenzabnahme von \textit{derb}.


%Außerdem wurde im Experiment im Rahmen der Hypothese {H1-Freq} dem Einfluss von Frequenz nachgegangen. Nach dieser verlieren Intensivierer ihre Expressivität, wenn sie zu häufig verwendet werden {\citep[vgl.][2]{Hauschild.1899}}. Basierend auf den Ergebnissen der Korpusuntersuchung ist davon auszugehen, dass die Gruppe der niedrigfrequenten Intensivierer demzufolge prinzipiell als expressiver wahrgenommen wird als die Gruppe der hochfrequenten Intensivierer. Wie anhand von Tabelle \ref{LinRegzaz} abgelesen werden konnte, ist diese Hypothese \textit{de facto} gültig: So wurde der Grad an Sprechereinstellung bei den niedrigfrequenten Intensivierern als signifikant höher empfunden als bei den hochfrequenten Intensivierern. Kein signifikanter Unterschied fand sich hingegen im Kontext von Skalarität, was bedeutet, dass sich die Frequenz eines Ausdrucks nachweislich nicht auf die Einschätzung dessen Skalaritätsgrades ausgewirkt hat (vgl. Tabelle \ref{LinRea}).


%Neben der Frequenz wurde auch die Persistenz berücksichtigt, die auf die Überprüfung der Hypothese H1-Temp abzielte. Diese ging im Kern auf \citet[59]{Tobler.1868} zurück, der postuliert, dass ein Intensivierer im Laufe der Zeit an Expressivität einbüßt. Ist diese Hypothese korrekt, so ist davon auszugehen, dass Intensivierer, die sich schon früh im Korpus nachweisen lassen, weniger expressiv sind als diejenigen Intensivierer, die erst im Laufe des Untersuchungszeitraums neu aufgetreten sind. Die Ergebnisse der inferenzstatistischen Auswertung lieferten keine empirische Evidenz für diese Hypothese. So enthielt das in Tabelle \ref{LinRegzaz2} gelistete finale Regressionsmodell im Kontext der in Aufgabe II erfragten Sprechereinstellung keine signifikanten Haupteffekte. Im Hinblick auf die in Aufgabe I untersuchte Skalarität zeigte sich ein signifikanter Unterschied zwischen den Persistenzgruppen (vgl. Tabelle \ref{LinRegzy}): Dementsprechend wurden die neuen Intensivierer mit einem höheren Skalaritätsgrad assoziiert als die älteren Intensivierer. Aufgrund des bestehenden Ungleichgewichts, das sich in Bezug auf den Umfang der Persistenzgruppen abzeichnete (d.\,h. drei neue Intensivierer vs. 13 ältere Intensivierer), sind diese Ergebnisse allerdings mit Vorsicht zu betrachten.


%Bei Aufgabe III wurde im Rahmen der Hypothese H3 ergänzend die Polarität der Sprechereinstellung abgefragt. Dabei wurde davon ausgegangen, dass die Interpretation des Intensivausdrucks von der Semantik des Ursprungslexems beeinflusst wird, wenn der Intensivierer ein semantisch leeres Pseudoadjektiv modifiziert. Dies bedeutet, dass die Beurteilung der Versuchspersonen durch die Polarität (und die oftmals damit verbundene Konnotation) des Ursprungslexems in die jeweils einschlägige Richtung gelenkt wird. In besonderem Maße sind hiervon vermutlich die Mitglieder der \glqq{}schrecklichen\grqq{} Intensivierergruppe betroffen, wohingegen opake Intensivierer wie \textit{krass} oder \textit{sehr} tendenziell verstärkt als neutral empfunden werden. Die Ergebnisse  von Aufgabe III deuteten darauf hin, dass diese Annahme grundsätzlich korrekt ist. So wurden Sätze mit deskriptivem Intensivierer sowie unintensivierte Sätze gemäß den Erwartungen vorwiegend als neutral beurteilt, während Sätze mit expressivem Intensivierer dagegen verstärkt in Richtung 'negativ' tendierten. Erklärt wurde dies mit der negativen Semantik, die vielen Mitgliedern der expressiven Intensiviererklasse zugrunde liegt. Diese Vermutung konnte durch die Übersicht in Tabelle \ref{ValenzInten} sowie Abbildung \ref{Aufg3.1} gestützt werden, nach denen Intensivierer mit negativem semantischem Ursprung eine ausgeprägte Tendenz für die negative Polarität aufwiesen. Bemerkenswert war auch das Ergebnis der verstärkt als positiv eingeschätzten Intensivierer: Bei ihnen handelte es sich nahezu ausschließlich um Ausdrücke, die bspw. in Ausrufesätzen in prädikativer Verwendung stehen können und in dieser Stellung durchgehend mit einer positiven Konnotation verbunden werden.


%Das berechnete Modell zeigte, dass vor allem diejenigen Intensivierer als positiv oder negativ wahrgenommen wurden, die mit einem hohen Grad an Sprechereinstellung einhergingen (vgl. Tabelle \ref{LinReg004}). Mit anderen Worten: Je expressiver ein Ausdruck war, desto größer war die Abweichung von 'neutral'. Ähnliches traf auf die Skalarität zu. Ein hoher Skalaritätsgrad führte zu einer stärkeren Abweichung von 'neutral' als ein niedriger. Dies lieferte empirische Evidenz für die Annahme, dass die Semantik des Ursprungslexems die Einschätzung der Einstellungspolarität im Falle der Modifikation eines Pseudowortes beeinflusst, was die zu testende Hypothese H3 untermauert.


%Kurz gefasst konnte durch die Ergebnisse des Experiments die Gültigkeit der meisten Hypothesen erhärtet werden. So fand sich in dem Sinne empirische Evidenz für die Hypothese {H1-Freq} nach \citet[2]{Hauschild.1899}, dass die im Korpus niedrigfrequenten Intensivierer auf der Skala signifikant höher verortet wurden als die im Korpus hochfrequenten Intensivierer. Daneben konnte gezeigt werden, dass die Interpretation des Intensivausdrucks \textit{de facto} durch die Semantik des Ursprungslexems gesteuert wird, wenn der Intensivierer mit einem semantisch leeren Pseudoadjektiv kombiniert wird. Dieser Befund stimmte mit der Hypothese H3 überein. Interessanterweise war die Abweichung von 'neutral' umso stärker, je höher der Grad an Skalarität und Sprechereinstellung evaluiert wurde. Keine empirische Evidenz fand sich hingegen für die Hypothese H1-Temp nach {\citet[59]{Tobler.1868}}. So unterschieden sich die im Korpus neu aufgetretenen Intensivierer im Hinblick auf den ihnen zugeschriebenen Grad an Sprechereinstellung nicht von den älteren Intensivierern. Auch die Vermutung, dass \textit{derb} aufgrund seiner diachronen Frequenzabnahme generell weniger expressiv ist als die anderen expressiven Intensivierer, konnte durch die Ergebnisse nicht gestützt werden. Wenngleich \textit{derb} sowohl hinsichtlich der in Aufgabe I abgefragten Skalarität wie auch der in Aufgabe II untersuchten Sprechereinstellung von den Versuchspersonen auf der Skala tendenziell niedrig positioniert wurde, unterschied es sich nicht merklich von anderen expressiven Intensivierern (bspw. \textit{arg}, \textit{ober} oder \textit{voll}). Nichtsdestotrotz deutet die diachrone Frequenzabnahme sowie die niedrige Skalenposition von \textit{derb} auf einen fortgeschrittenen Expressivitätsverlust hin, der sich vermutlich fortsetzen wird.


%In Kapitel \ref{Experimente} wird eine weitere experimentelle Untersuchung vorgestellt. Darin wird die theoretische Grundannahme hinterfragt, dass expressive Intensivierer gleichsam mit einem höheren Grad an Skalarität und Sprechereinstellung einhergehen als deskriptive Intensivierer (vgl. {\citealt[315]{Lang.1983}}; \citealt[22]{Paradis.1997}; \citealt[57]{Ortner.2014}; \citealt[133]{Gutzmann.2019a}). Ist diese faktisch gültig, so ist anzunehmen, dass Intensivausdrücke mit expressivem Intensivierer sowohl hinsichtlich Skalarität als auch Sprechereinstellung auf der Skala grundsätzlich höher positioniert werden als Intensivausdrücke mit deskriptivem Intensivierer (vgl. {\citealt[315]{Lang.1983}}; \citealt[22]{Paradis.1997}, {\citealt[57]{Ortner.2014}}; {\citealt[133]{Gutzmann.2019a}}). Um der zu testenden Fragestellung gerecht zu werden, wird eine Stichprobe aus jeweils vier deskriptiven und expressiven Intensivierern in Bezug auf den mit ihnen assoziierten Grad an Skalarität und Sprechereinstellung in einem ähnlich aufgebauten Rating-Experiment kontrastiert. Das Experiment setzt sich aus zwei Teilen zusammen: Im ersten Teil geht es um den durch die Intensivierer ausgedrückten Grad an Skalarität, im zweiten um den ausgedrückten Grad an Sprechereinstellung.
%Dabei wird eine Stichprobe der untersuchten expressiven Intensivierer mit der gleichen Anzahl an ausgewählten deskriptiven Intensivierern kontrastiert.


%Ergebnis S. 174, dass Intensivierer, wenn abgenutzt, nicht notwendig skalarer sind




%vielleicht Fabians Erklärung für \textit{derb} ausm Kolloquium mit einbeziehen






%Wright, Gaskell & Muircheartaigh (1995: 173): "An alternative explanation is that people do not have specific scalar values associated with each intensifier but they do realize the ordering of common intensifiers. This hypothesis has been put forward with regard to verbal quantifiers (Newstead, 1988; Pepper, 1981), and seems plausible to apply also to intensifiers. When asked the ‘satisfied’ question and then the ‘very satisfied’ question, the subjects in condition 1 appreciate that the latter means more satisfied than the former, but not how much more. When making a judgement, subjects use the answers to the first question as an anchor for the second; they shift their responses downwards irrespective of the intensifier used."



